\documentclass[usenatbib,twoside,twocolumn]{mnras}

% Paper packages
\usepackage{graphicx}
\usepackage{amsmath}
\usepackage{amssymb}
\usepackage{booktabs}

% Denser float packing (reduce float-only pages; no content change)
\renewcommand{\topfraction}{0.9}
\renewcommand{\bottomfraction}{0.85}
\renewcommand{\textfraction}{0.07}
\renewcommand{\floatpagefraction}{0.75}
\renewcommand{\dbltopfraction}{0.9}
\renewcommand{\dblfloatpagefraction}{0.75}
\setcounter{topnumber}{3}
\setcounter{bottomnumber}{2}
\setcounter{totalnumber}{5}
\setcounter{dbltopnumber}{3}

% Define astronomical units
\newcommand{\msun}{\ensuremath{M_{\odot}}}
\newcommand{\pc}{\ensuremath{\rm pc}}
\newcommand{\kms}{\ensuremath{\rm km\,s^{-1}}}

\title[Core spacing and filament occupancy in HGBS filaments]{Core Spacing in Herschel Gould Belt Filaments: Short Local Separations, Inhomogeneous Occurrence, and the Filament Length Eligible to Fragment\thanks{Both results are established for the nearby, low-mass Gould Belt sample ($N=4$ clouds) and are not offered as a universal fragmentation law.}}

\author[G.~J.~White]{G.~J.~White$^{1,2}$\\
$^{1}$School of Physical Sciences, The Open University, Walton Hall, Milton Keynes, MK7 6AA, UK\\
$^{2}$Rutherford Appleton Laboratory, Harwell Science and Innovation Campus, Didcot, OX11 0QX, UK}

\begin{document}

\maketitle

\begin{abstract}
Dense cores form along the filaments that thread star-forming clouds. If they come from the classical gravitational fragmentation of a cylinder they should be separated by about five to six filament widths. We re-examine that expectation in four \textit{Herschel} Gould Belt Survey clouds using \textit{Gaia} Data Release 3 distances.

The widely used median of all pairwise separations cannot measure a local spacing: it converges on a fixed fraction of the region's length, measuring the area selected rather than the physics. What a one-fragment-per-wavelength theory predicts is the filament length per core; on that quantity two of our four clouds agree with the classical value and two fall below it. Neighbouring-core separations, measured only where cores occur, are several times shorter under every filament and width definition tried. The gap between the two is itself the measurement: most of the crest is bare, so cores sit in tight groups separated by long empty stretches.

Restricting to crest above the critical line mass brings the two together, but the correction weakens once the cores' own contribution to the selecting column density is removed, after which only two clouds retain enough cores to decide. A deficit of the smallest separations is indistinguishable from the extraction resolution limit.

Simulations show a contracting filament need not fragment at the equilibrium wavelength, but resolve no growth-rate maximum and give only an upper bound: a counterexample, not an explanation. With four clouds these are sample statements, and the pipeline has so far been run only by the author.
\end{abstract}

\begin{keywords}
stars: formation -- ISM: structure -- ISM: filaments -- MHD -- turbulence
\end{keywords}

\section{INTRODUCTION}

The \textit{Herschel} Space Observatory demonstrated the ubiquity and importance of filamentary structure in nearby star-forming clouds \citep{Andre2010}, with a characteristic width of $\sim$0.1\,pc across diverse environments \citep{Arzoumanian2011}, though the universality, and even the physical interpretation, of this scale is contested \citep{Panopoulou2017,Panopoulou2022}, and a critical mass-per-unit-length $M_{\rm line,crit}\approx16~M_\odot\,{\rm pc^{-1}}$ \citep{Ostriker1964}.

\emph{The observation.} We quantify the core-to-core spacing homogeneously in four nearby HGBS clouds with Gaia DR3 distances \citep{Zhang2023}. Throughout we keep three quantities distinct and do \emph{not} equate them: the observed adjacent-core arclength spacing (hereafter adjacent-core spacing, ACS) statistic $s_{\rm ACS,med}$ (the median adjacent-core gap, a line-density scale, \emph{not} itself a fragmentation wavelength); the fragmentation wavelength of our simulated filaments. For the contracting runs no maximum of $\Gamma(\lambda)$ is resolved, so we never quote one and report only the upper bound $\lambda_{\rm turn}$ on where a turnover could lie; and the classical Inutsuka--Miyama equilibrium value $\lambda_{\rm m}\approx5$--$6\times$FWHM \citep{Inutsuka1992,Fischera2012}. Two empirical results anchor the paper. First, the appropriate estimator is the ACS spacing: the widely-used all-pairs (pairwise) median is a region-extent-biased estimator \emph{of local core spacing}, it validly estimates the segment extent but is inappropriate as a local-spacing estimator, converging to $\sim L/3$, the extent of whatever segment it is applied to (an individual filament, $L_{\rm fil}\sim1$\,pc, or a whole region), rather than to any local spacing (Section~\ref{sec:statistics}); the bias applies wherever an all-pairs or pooled separation statistic is adopted, for instance in cross-region compilations, and we make no claim that published HGBS spacings generally suffer from it. Table~\ref{tab:literature} shows the opposite: the values we could check are local or along-fibre statistics that agree with our ACS measurements. Second, on that estimator the observed conditional spacing $s_{\rm ACS,med}/W$ lies substantially below the equilibrium-cylinder wavelength $\lambda_{\rm m}$ under every skeleton and width convention. The deficit is a factor of roughly $2.4$--$4.3$ below $\lambda_{\rm m}$ for the DisPerSE conventions, and further still under FilFinder; the residual magnitude is set by the skeleton algorithm and the width convention. Because several different ranges appear in the literature and in this paper, we state once here how they relate: $2.4$--$4.3$ is the factor by which $S_{\rm local}$ falls below $\lambda_{\rm m}$ for the DisPerSE conventions; $2$--$13$ is the same factor across the full convention envelope including FilFinder; and $0.8$--$1.9$ (DisPerSE) or $0.4$--$1.9$ (including FilFinder) are the values of $S_{\rm local}$ itself. Table~\ref{tab:headline} is the decoder. The primary observed result is this convention \emph{envelope}: $s_{\rm ACS,med}/W\simeq0.8$--$1.9$ for the DisPerSE conventions, widening to $\simeq0.4$--$1.9$ once the FilFinder skeleton is included (Fig.~\ref{fig:envelope}); one illustrative point within it is the region-Plummer/DisPerSE ACS-spacing ratio $(s_{\rm ACS,med}/W)_0\approx1.4$ (Table~\ref{tab:headline}), which we carry only as a bookkeeping reference convention, not the observed value. Table~\ref{tab:headline} is a reader's guide to every ratio quoted in this paper, grouped by whether it is a conditional local descriptor or an unconditional line density; the notation separating the observed statistic from the simulated wavelength is defined in Section~\ref{sec:statistics}. We state the limit of the interpretation at the outset rather than deferring it to the discussion: three mechanisms could produce a sub-classical spacing, namely fragmentation from a contracting rather than an equilibrium state, projection of unresolved velocity-coherent fibres, and helical magnetic tension. The weight of simulated evidence is not evenly distributed among them: the first is the one the campaign actually addresses, the second is treated only by a geometric toy model, and the third is untested, because the single field configuration that could shorten the axial wavelength, a force-balanced helical equilibrium, was not simulated. The data presented here cannot discriminate between the three. Nothing in this paper should be read as identifying the cause. One further limitation belongs here rather than in the data-availability statement: the skeleton extraction, the injection--recovery corrections, the growth-rate measurement and the forward-modelling to observable units all rest on a single pipeline, written and run by one person. Only the adjacent-core arithmetic has been re-derived independently, and that by a second implementation written by the same author. The code and catalogues are deposited and public, and we regard external re-execution as a precondition for treating the pairwise-median result as settled rather than as reported.

\emph{Candidate explanations.} Such a deficit relative to $\lambda_{\rm m}$ has been read as evidence for additional physics. \citet{Arzoumanian2019} established the $\sim$0.1\,pc \emph{width}, not a spacing; the sub-classical \emph{spacing} and its usual fibre/magnetic reading come from the fragmentation literature \citep{Fischera2012,Zhang2023cal}. Three mechanisms are candidates. (i) \emph{Hierarchical fragmentation}, in which filaments are bundles of velocity-coherent fibres \citep{Hacar2013,Hacar2018,Yang2024}, so that measuring core spacings across the bundle compresses the apparent filament-level spacing. (ii) \emph{Magnetic tension} \citep{Nakamura1993}: only a \emph{helical}/toroidal field shortens the axial mode. A static longitudinal field leaves it unchanged and a perpendicular field lengthens it (Section~\ref{sec:magnetic}), so a magnetic route to a shorter spacing is a specifically helical effect. Its force-balanced case is untested here, and we deliberately leave the helical/toroidal magnetic mechanism as an identified open problem that our MHD campaign does not close (Section~\ref{sec:optional_tests}). (iii) \emph{Dynamical fragmentation}: a still-collapsing filament fragments at a scale set by its current, contracting state rather than by the marginally-stable equilibrium mode.

That real, still-accreting filaments need not fragment at the equilibrium wavelength is already qualitatively expected in the dynamical-ISM paradigm, in which filaments form and evolve out of equilibrium while accreting \citep{Andre2014,Clarke2016,Clarke2017,Clarke2020}. Our aim is therefore not to propose a new mechanism but to quantify the observation and test whether the simplest dynamical reading is even viable. The genuinely new contributions are three: (i) a homogeneous, Gaia-DR3, ACS-corrected, multi-convention observational quantification of the spacing across the four primary clouds; (ii) identifying the region-extent ($\approx0.293\,L$) bias as a general bias in all-pairs and pooled separation statistics; the underlying triangular-distribution order statistic is elementary and textbook, and region-extent biases of this kind are already familiar in spatial point-process statistics (the same Clark--Evans nearest-neighbour context used in Section~\ref{sec:clustering}), so we claim no new mathematics, only the identification of this estimator bias in the filament-spacing context; and (iii) a specific longitudinal near-critical MHD simulation that fragments sub-classically, with an equilibrium-recovery control that recovers the classical wavelength under the identical pipeline.

We present the full-coverage ACS analysis for the four primary HGBS regions (Orion~B, Aquila, Perseus, Taurus; the eight-region sample is retained only for the legacy pairwise-median statistic), followed by the simulation campaign, its validation and the interpretation. Two groups of statements should be kept apart. \emph{Established here:} the pairwise median is a region-extent-biased estimator of local core spacing (converging to $\approx0.293L$); the ACS spacing is the appropriate observable; the spacing is sub-classical in \emph{sign} under every convention, while its magnitude is convention-dominated, the skeleton choice being the single largest term. \emph{Not established:} that dynamical collapse is the dominant cause, that hierarchical fibres are unimportant, or that magnetic tension is excluded. Symbols are collected in Table~\ref{tab:symbols}. With only four independent clouds, these results characterise the nearby Gould-Belt sample alone, not a universal fragmentation law.

\begin{table}
\caption{Principal symbols.}
\label{tab:symbols}
\footnotesize
\setlength{\tabcolsep}{4pt}
\begin{tabular}{lp{6.1cm}}
\toprule
Symbol & Meaning \\
\midrule
$s_{\rm ACS,med}$ & observed adjacent-core spacing (ACS) statistic (median adjacent gap; a line-density scale, \emph{not} a resonant/fragmentation wavelength; cf.\ Section~\ref{sec:clustering}) \\
$W_{\rm fil}$ & observed filament FWHM (beam-deconvolved Plummer; the reference width) \\
$W_{\rm core}$ & simulation initial Gaussian half-width $\sigma$ ($0.3\,\lambda_{\rm J}$) \\
$W_{\rm form}$ & FWHM of the evolved, projected simulation profile \\
$T_1\!\equiv\!\eta_W$ & $W_{\rm form}/W_{\rm fil}$, the simulation-to-observation width normalisation (forward-model band $0.59$--$0.78$; adopted $0.61$) \\
$\lambda_{\rm J}$ & Jeans length at the \emph{initial background} density $\rho_0$ ($=1$ in code units) and never at the evolving central density, for which we write $\lambda_{\rm J}(\rho_c)$ \\
$\lambda_{\rm peak}$ & wavelength at which $\Gamma(\lambda)$ is maximal. Resolved \emph{only} for the relaxed equilibrium control (at $22H$); \emph{not} resolved for any contracting run, for which we quote $\lambda_{\rm turn}$ instead \\
$\lambda_{\rm turn}$ & upper bound on where a turnover in $\Gamma(\lambda)$ could lie for a contracting filament ($\lesssim1.1\,\lambda_{\rm J}$); the only wavelength we quote for those runs \\
$\lambda_{\rm m}\!\equiv\!\lambda_{\rm IM92}$ & classical Inutsuka--Miyama equilibrium-cylinder wavelength $\approx5$--$6\,W_{\rm FWHM}$ ($=4\times$ flat diameter) \\
$\Gamma(\lambda)$ & linear growth rate of axial mode of wavelength $\lambda$ \\
$f$ & thermal line-mass fraction $M_{\rm line}/M_{\rm line,crit}$ \\
$\beta$ & plasma beta $8\pi P/B^2$ \\
$\mathcal{M}$ & amplitude label of the seed perturbation ($\delta v=\mathcal{M}c_s\times10^{-4}$; a linear seed, \emph{not} a turbulent Mach number; driven transonic turbulence is tested separately in Appendix~\ref{sec:supplementary_campaigns}) \\
$\mu_\Phi/\mu_{\Phi,\rm crit}$ & normalised mass-to-flux ratio \\
$\theta$ & angle between field and filament axis \\
$N_J$ & cells per local Jeans length (Truelove criterion) \\
$C$ & density contrast $\rho_{\rm max}/\langle\rho\rangle$ \\
\bottomrule
\end{tabular}
\end{table}

\noindent\textbf{The result in one sentence, without symbols.} In each of the four clouds examined here, the cores that do form are found in tight groups; the groups are far apart; and most of the filament crest carries no core at all. Everything below is an attempt to measure those three statements and to decide which of them a fragmentation theory is entitled to be tested against.

\smallskip\noindent\textbf{Notation, and which quantity tests the theory.} Four normalised quantities carry the argument, and we use the capital symbols whenever a ratio to the classical benchmark is meant, reserving $s_{\rm ACS,med}$ for raw separations in parsecs:
\begin{itemize}\setlength{\itemsep}{1pt}
\item $S_{\rm global}=(L/N)/W$, the unconditional filament length per core. \emph{This is the quantity a one-fragment-per-wavelength theory predicts}, and it is therefore the primary comparison with $\lambda_{\rm m}/W\approx5$--$6$.
\item $S_{\rm local}=s_{\rm ACS,med}/W$, the median adjacent-core separation measured \emph{only within core-bearing stretches}. This is a descriptor of the internal spacing of core-bearing regions, not a fragmentation wavelength, and it is not the primary test.
\item $S_{\rm elig}=(L_{\rm elig}/N_{\rm elig})/W$, $S_{\rm global}$ restricted to filament above the critical line mass.
\item $S_{\rm sim}^{\rm upper}=\lambda_{\rm turn}/W$, an \emph{upper bound} from the simulations, never a measured wavelength.
\end{itemize}
The difference between $S_{\rm global}$ and $S_{\rm local}$ is not a discrepancy to be explained away: it is itself the measurement of how unevenly cores occupy the crest. Every other ratio quoted here is a sensitivity test or alternative convention on one of these four, decoded in Table~\ref{tab:headline}.

\section{OBSERVATIONAL ANALYSIS}

\subsection{Complete HGBS Sample}

Our sample comprises 8 high-confidence HGBS regions (four \emph{primary}, Table~\ref{tab:sample}; four \emph{limited}, Table~\ref{tab:sample_limited} in Appendix~\ref{sec:limited_sample}). The original catalogues \citep{Konyves2015, Konyves2020, Ladjelate2020, Pezzuto2021} adopted distances of $130$--$400$\,pc from the best estimates available in 2010--2015 (Table~\ref{tab:literature} lists the pre-Gaia distance and its source per region; note that \citealt{Konyves2020} already adopted $\sim400$\,pc for Orion~B); Gaia DR3 parallaxes \citep{Zhang2023} have since enabled substantial revisions, largest for the more distant regions, whose pre-Gaia distances were the most uncertain.

We adopt the Gaia DR3 YSO-based distances from \citet{Zhang2023} (Table~4), with typical uncertainties of 3\%--5\%. The revisions are strongly region-dependent (largest for Serpens $260\to458$\,pc, $+76\%$, and Aquila $260\to436$\,pc, $+68\%$; Orion~B, Taurus and Ophiuchus essentially unchanged); since spacings scale linearly with distance, the core-count-weighted mean rises only modestly ($\sim10\%$ for the primary regions), dominated by Aquila rather than being a uniform $+32\%$ shift. Because Gaia DR3 thus compresses or expands each cloud non-uniformly, the revision is not a single global rescaling, which is a further reason the clouds must be treated as $N=4$ independent units rather than by pooling cores that carry very different distance corrections. The observed spacing is a line-density statistic, not a resonance, and coincides with the theoretical wavelength only to order of magnitude.

The analysis includes all HGBS-identified cores (starless, prestellar and protostellar), extracted and classified within the consistent HGBS framework \citep{Konyves2015, Konyves2020}. Including unbound starless and protostellar sources may introduce some bias, but fragmentation produces a continuum of boundness states, and restricting the sample to prestellar cores would shrink several regions (notably TMC1, CRA and Serpens) below statistical usefulness. Core positions and properties are taken directly from the published catalogues.

Protostellar migration displacements of $0.01$--$0.05$\,pc \citep{Kirk2016, Mattern2018} are small against the $\sim0.2$--$0.3$\,pc spacing; we estimate the resulting ACS bias at $\sim5$--$10\%$ and carry it as a systematic (Section~\ref{sec:error_budget}).

\begin{table}
\caption{Primary HGBS sample with Gaia DR3 distances. Per-region spacings are pairwise-median values rounded to three significant figures; the core-count-weighted mean of the four primary regions is $0.290$\,pc. The four \emph{limited} regions (Ophiuchus, Serpens, TMC1, CRA) and the full 8-region weighted mean ($0.279$\,pc) are listed in Appendix~\ref{sec:limited_sample}, Table~\ref{tab:sample_limited}.}
\label{tab:sample}
\footnotesize
\setlength{\tabcolsep}{3pt}
\begin{tabular}{lccccc}
\toprule
Region & $D$ & $N_{\rm core}$ & Pairwise med. & S.E. & Status \\
       & (pc) &  & (pc) & (pc) &  \\
\midrule
Orion B  & 386 & 1,844 & 0.313 & 0.047 & Primary \\
Aquila   & 436 & 749   & 0.346 & 0.047 & Primary \\
Perseus  & 300 & 816   & 0.248 & 0.040 & Primary \\
Taurus   & 135 & 536   & 0.198 & 0.040 & Primary \\
\midrule
\textbf{Weighted Mean}$^\dagger$ &  & \textbf{3,945} & \textbf{0.290} &, &  \\
\bottomrule
\end{tabular}
\par\vspace{2pt}
\footnotesize $^\dagger$Core-count-weighted mean of the four primary regions. A formal core-count-weighted standard error (treating cores as independent) is \emph{not} the relevant uncertainty, since the four clouds are the independent units (Section~\ref{sec:error_budget}); the cloud-to-cloud scatter is carried instead.
\end{table}

\subsection{Measurement Methodology}
\label{sec:methodology}

The analysis pipeline, from raw HGBS maps to the final characteristic $s_{\rm ACS,med}/W$, is summarised in Fig.~\ref{fig:pipeline}. Filament skeletons were identified using DisPerSE \citep{Sousbie2011} following standard HGBS methodology \citep{Arzoumanian2011}, with a $3\sigma$ persistence threshold and manual editing restricted to removing obvious artefacts (e.g.\ spurs connecting unrelated structures). Its effect on the spacing is small ($\lesssim5\%$), bounded by the junction-retention and association-threshold sensitivity tests as a \emph{proxy} (Table~\ref{tab:errorbudget}). Because the skeleton algorithm is the dominant systematic in this paper, robustness to the persistence threshold of the adopted ontology is part of the measurement and not future work, and we have therefore tested it in two ways (Section~\ref{sec:persistence}). A blind, multi-editor re-skeletonisation remains outstanding (Section~\ref{sec:limitations}). Because skeletons and core positions derive from distance-independent angular positions, the Gaia DR3 revisions require no re-extraction. The measured spacing is insensitive to these choices: excluding junctions changes it by $<5\%$ and varying the association threshold between $1\times$ and $1.5\times$ the width by $<3\%$ (a coarse two-point check, distinct from the finer $0.04$--$0.08$\,pc half-width sweep of Section~\ref{sec:injection_recovery} that gives a $4$--$7\%$ per-region spread; the primary ACS pipeline associates cores within a half-width, $\sim0.06$\,pc, of a spine).

Core-to-core spacings were also measured using the pairwise median statistic, which Section~\ref{sec:statistics} shows converges to $\approx0.293L$ rather than the fragmentation wavelength for $N\gtrsim5$; we retain it only as a comparability statistic with prior HGBS work.

Because physical spacing is angular separation times distance, the $\lambda\propto d$ scaling is tautological and we do not present it as corroboration; the relevant check is that the \emph{residuals} about that scaling show no trend with the absolute revision fraction $|\Delta D/D_{\rm original}|$, the two most-revised clouds (Aquila, Serpens) do not stand out as having the largest or smallest residuals, so the revision does not preferentially bias the most-revised clouds. With only four clouds a formal correlation test has negligible power and we do not quote one. Two caveats remain. The YSO-based distances \citep{Zhang2023} are assumed to apply to the dense gas measured here, and the Aquila/Serpens Rift has substantial line-of-sight depth ($\sim$\,tens of pc in the 3-D dust reconstructions of \citealt{Zucker2020} and \citealt{Green2024}), so a single assigned distance may blend populations. We quantified the blending with a geometric Monte-Carlo model in which two crests at different depths are projected onto a common spine, taking an intrinsic spacing of $0.20$\,pc at $d=436$\,pc. The model draws each crest's line-of-sight offset from a Gaussian of dispersion $\sigma_z$, with the two crests separated by $2\sigma_z$ in depth, and we scan $\sigma_z=5$, $10$, $15$ and $30$\,pc (total separations of $10$--$60$\,pc, spanning the $\sim$\,tens of parsecs quoted by \citealt{Zucker2020} and \citealt{Green2024} for the Rift). Superposition interleaves the two chains and \emph{shortens} the measured median, to $0.41$--$0.49$ of the true value across that range, whereas the distance spread alone (which rescales each spacing by $d/(d+z)$) contributes only $\pm1$ to $\pm7$\,per cent. Blending therefore biases the ACS towards \emph{more} sub-classical values rather than less. The $\sim50$\,per cent figure is a worst case in which two crests overlap completely; the realistic effect is smaller because most associated cores lie on coherent crests. We carry it as a one-sided systematic in Table~\ref{tab:errorbudget}. As a floor, retaining the pre-Gaia distances lowers the primary-region characteristic to $s_{\rm ACS,med}/W\approx1.4$ (still sub-classical), so the sub-classical \emph{sign} does not depend on the distance revision, though its magnitude does.

\begin{figure}
\centering
\begin{minipage}{0.95\columnwidth}
\footnotesize
\begin{enumerate}\setlength{\itemsep}{2pt}
\item HGBS column-density maps ($250\,\mu$m reference) + Gaia DR3 distances
\item DisPerSE skeleton extraction ($3\sigma$ persistence threshold)
\item Morphological closing (8 iterations) $+$ re-skeletonisation (bridge fragmented segments; components $\geq40$\,px)
\item Core association: cores within $0.06$\,pc of a spine are assigned; others rejected
\item Graph ordering: cores ordered by double-BFS arclength along each component
\item Adjacent-core ACS spacing (in pc, via Gaia DR3 distance)
\item Injection--recovery bias correction ($\div\,1.11$--$1.29$, region-dependent)
\item Characteristic $s_{\rm ACS,med}/W = $ bias-corrected median spacing $/\,W_{\rm fil}$
\end{enumerate}
\end{minipage}
\caption{Analysis pipeline from raw HGBS maps to the final characteristic $s_{\rm ACS,med}/W$ (a line-density statistic, not a wavelength). Each step's parameters are given in the text (Sections~\ref{sec:statistics}--\ref{sec:injection_recovery}).}
\label{fig:pipeline}
\end{figure}

The measurement convention, above all the skeleton algorithm, dominates the numerical value: the DisPerSE and FilFinder skeletons of the identical Orion~B map (Fig.~\ref{fig:skeleton_compare}) trace substantially different networks, moving the inferred $s_{\rm ACS,med}/W$ by a factor $\sim$2 while the sub-classical \emph{sign} holds for both (Section~\ref{sec:error_budget}).

\begin{figure}
\centering
\includegraphics[width=0.99\columnwidth]{figures/fig_skeleton_compare.pdf}
\caption{The dominant observational systematic, shown directly: the DisPerSE (blue) and FilFinder (red) skeletons extracted from the \emph{same} \textit{Herschel} column-density map of Orion~B (greyscale, $\log N_{\rm H_2}$; a representative sub-region). The FilFinder skeleton is produced by \textsc{FilFinder2D} with the cloud distance set per region and an \emph{externally supplied} mask (\texttt{use\_existing\_mask=True}) followed by \texttt{medskel()}; the mask is a simple column-density cut at the $92$nd percentile of the finite pixels, which corresponds to $N_{\rm H_2}=2.5$, $3.8$, $3.1$ and $2.9\times10^{21}$\,cm$^{-2}$ in Orion~B, Aquila, Perseus and Taurus respectively. Because the mask is supplied externally, \textsc{FilFinder}'s own adaptive and branch thresholds are not invoked; earlier versions of this caption quoted nominal adaptive and branch thresholds that were not in fact used, including a ``$40$\,pc'' branch threshold that was neither used nor physically sensible. A more permissive $70$th-percentile mask is also computed and is what sets the upper end of the FilFinder range. The DisPerSE skeletons are the published HGBS products, extracted at a $3\sigma$ persistence cut. The two algorithms trace substantially different filament networks, FilFinder's medial axis of a column-density mask is denser and more branched and follows more low-contrast material, whereas DisPerSE follows the persistent crests, which is why the inferred $s_{\rm ACS,med}/W$ differs by a factor of $\sim$2 between them (Section~\ref{sec:error_budget}). This is intrinsic to how each algorithm defines a ``filament'' and is the largest single uncertainty on the \emph{magnitude} of the spacing, though the sub-classical \emph{sign} holds for both.}
\label{fig:skeleton_compare}
\end{figure}

\subsection{Results}

The four primary regions are the independent statistical units throughout this paper; individual cores and spacings within a cloud are not independent, and we propagate this nesting in the hierarchical bootstrap (Section~\ref{sec:error_budget}). With only four independent clouds we treat the cloud-to-cloud scatter ($\sigma\approx0.4$) as a descriptive range, not a formal confidence interval; we do not attach a bootstrap CI to an $N=4$ sample.

\smallskip\noindent\textbf{Two statistics, two different answers.} Because this distinction governs everything that follows, we state it before any numbers. The adjacent-core spacing $s_{\rm ACS,med}$ is a \emph{conditional} descriptor: it is measured only where cores are adjacent, so it characterises the interior of core-bearing stretches and says nothing about filament that has formed no cores. The core line density is $\rho_{\rm core,1D}=N/L$, cores per unit length of the whole skeleton. For comparison with a wavelength we use its reciprocal, $L/N$, the mean filament length per associated core, and it is this \emph{unconditional} quantity, not $s_{\rm ACS,med}$, that a one-fragment-per-wavelength theory predicts. The two give different answers, and that difference is the main observational result of this paper. Taking the theoretically predicted quantity first: unconditionally, $S_{\rm global}=(L/N)/W\approx3.4$, $3.4$, $4.5$ and $5.8$ for Orion~B, Aquila, Taurus and Perseus, so Perseus is essentially consistent with the classical expectation, Taurus is close to it, and only Orion~B and Aquila fall clearly below the nominal $5$--$6\,W$ (Section~\ref{sec:linedensity}). \emph{On the comparison that classical theory actually makes, therefore, two of the four clouds are not discrepant at all.} Conditionally, by contrast, $S_{\rm local}$ is sub-classical by a factor $\approx2.4$--$4.3$ (Table~\ref{tab:headline}). The reconciliation is that cores occupy only $6$--$27$ per cent of the filament length (Section~\ref{sec:intermittency}). For brevity we refer to the conditional local ratio as $S_{\rm local}\equiv s_{\rm ACS,med}/W$, the unconditional one as $S_{\rm global}\equiv(L/N)/W$, the global mean length per core in units of $W$, the same quantity restricted to fragmentation-eligible filament as $S_{\rm elig}\equiv(L_{\rm elig}/N_{\rm elig})/W$ (Section~\ref{sec:eligible}), and the simulated bound as $S_{\rm sim}^{\rm upper}\equiv\lambda_{\rm turn}/W$. Table~\ref{tab:principal} collects all four per cloud, and every other number quoted in this paper is a sensitivity test or an alternative convention on one of them. Every comparison of $S_{\rm local}$ with $5$--$6\,W$ in what follows, including Table~\ref{tab:headline} and Fig.~\ref{fig:spacing}, is therefore a comparison between a conditional local descriptor and a fragmentation wavelength. It is not a measurement of a shortened fragmentation wavelength, and we do not use it as one.

\begin{table}
\caption{Principal results, one row per cloud. $S_{\rm local}=s_{\rm ACS,med}/W$ is the conditional adjacent-core spacing, measured only within core-bearing stretches; $S_{\rm global}=(L/N)/W$ is the unconditional filament length per core; $S_{\rm elig}=(L_{\rm elig}/N_{\rm elig})/W$ restricts the latter to filament above the critical line mass, selected on a core-masked column-density reconstruction (Section~\ref{sec:eligible}), with $f_{\rm elig}$ the eligible length fraction and $N_{\rm elig}$ the cores on it; $f_{\rm occ}$ is the $0.06$\,pc occupancy fraction. All are normalised by the fixed $W=0.10$\,pc width so the columns are directly comparable, and all use the DisPerSE skeleton. The classical benchmark is $5$--$6$. The simulated upper bound is common to all rows, $S_{\rm sim}^{\rm upper}=\lambda_{\rm turn}/W\lesssim1.0$--$1.3$ (Section~\ref{sec:sims}); it is an upper bound, not a measured wavelength. Perseus retains only $57$ eligible cores after masking and its value is masking-radius dependent (Table~\ref{tab:maskrad}), so it is bracketed; Taurus retains $13$ and its value is not usable at any masking radius, so it is omitted rather than quoted. With $N=4$ clouds these are descriptive values, not a population sample.}
\label{tab:principal}
\footnotesize
\setlength{\tabcolsep}{4pt}
\begin{tabular}{lccccc}
\toprule
Cloud & $S_{\rm local}$ & $S_{\rm global}$ & $S_{\rm elig}$ & $f_{\rm elig}$ [$N_{\rm elig}$] & $f_{\rm occ}$ \\
\midrule
Orion~B & 1.7 & 3.4 & 2.0    & 0.16 [268] & 0.27 \\
Aquila  & 2.1 & 3.4 & 3.2    & 0.90 [303] & 0.24 \\
Perseus & 3.0 & 5.8 & (4.1)  & 0.08 [57]  & 0.11 \\
Taurus  & 1.7 & 4.5 & ---    & --- [13]   & 0.06 \\
\bottomrule
\end{tabular}
\end{table}

\begin{table}
\caption{Convention decoder for the \emph{conditional} local descriptor: read every observed $s_{\rm ACS,med}/W$ quoted in the text against it. This table tabulates the adjacent-core spacing, which is measured only within core-bearing stretches and is \emph{not} the quantity a one-fragment-per-wavelength theory predicts (that is $L/N$; Section~\ref{sec:linedensity}). Values are given under each width/skeleton convention and correction state for the four primary HGBS regions (revised distances); every combination is conditionally sub-classical (classical $\lambda_{\rm m}\approx5$--$6\times$FWHM). We adopt the directly-measured region-Plummer value as our reference, $\boldsymbol{(s_{\rm ACS,med}/W)_0\approx1.4}$; the fixed-width value ($\approx1.9$, retained for comparability with prior HGBS work) and all others are read against it. The ``periodic'' column applies the periodic injection--recovery correction to the raw value; the data-matched \emph{clustered} column (bracketed) applies $\div1.64$. The DisPerSE conventions are sub-classical by a factor $\approx2.4$--$4.3$; across the full envelope the factor spans $\approx2$--$13$, the upper end set by the densest FilFinder skeleton.}
\label{tab:headline}
\footnotesize
\setlength{\tabcolsep}{4pt}
\begin{tabular}{p{4.6cm}cc}
\toprule
Quantity & raw & periodic \\
\midrule
\multicolumn{3}{@{}l}{\emph{Conditional local descriptor} $s_{\rm ACS,med}/W$ (core-bearing stretches only)}\\
Region Plummer, DisPerSE (adopted, $(s_{\rm ACS,med}/W)_0$) & $\mathbf{1.4}$ & $1.2$ \\
Fixed $0.10$\,pc, DisPerSE & $2.1$ & $1.9$ \\
FilFinder skeleton & $0.4$--$0.8$ & $0.4$--$0.8$ \\
Per-segment local width, $\Delta s/W_{\rm local}$ & \multicolumn{2}{c}{$0.6$--$1.1$} \\
\emph{Deprojected} (adopted row $\div\,\langle\sin i\rangle_{\rm med}=0.87$) & $1.6$ & $1.4$ \\
\emph{Deprojected} (fixed $0.10$\,pc) & $2.4$ & $2.2$ \\
\midrule
\multicolumn{3}{@{}l}{\emph{Unconditional} (what fragmentation theory predicts)}\\
Global length per core $(L/N)/W$, fixed $0.10$\,pc & \multicolumn{2}{c}{$3.4$--$5.8$} \\
Mean adjacent, fixed $0.10$\,pc (branch-jump inflated) & \multicolumn{2}{c}{$2.1$--$4.7$} \\
\midrule
\multicolumn{3}{@{}l}{\emph{Legacy comparator and reference scales}}\\
Pairwise median ($\approx0.293L$; biased, not a spacing) & $2.8$ & \\
Matched simulation (upper limit) & \multicolumn{2}{c}{$\lambda_{\rm turn}/W\lesssim1.0$--$1.3$} \\
Classical $\lambda_{\rm m}$ (F\&M 2012) & \multicolumn{2}{c}{$5$--$6$ ($=4\times$ diam.)} \\
\bottomrule
\end{tabular}
\\[2pt]
\footnotesize
\emph{Notes} The reference point is the region-Plummer ACS $(s_{\rm ACS,med}/W)_0\approx1.4$ (bold), one illustrative point within the convention range $0.8$--$1.9$; it is an adopted reference value, chosen for comparability with prior HGBS work, not a uniquely preferred physical scale (and not a formal confidence interval; Section~\ref{sec:error_budget}). The data-matched \emph{clustered} branch ($\div1.64$ on the raw value, giving $\approx0.8$ region-Plummer, $\approx1.2$ fixed-width and $\approx0.3$--$0.5$ FilFinder) is \emph{not} shown as a column because it is not independent: it applies a correction derived by \emph{assuming} the observed clustering (Section~\ref{sec:injection_recovery_model}), so it is circular and serves only as a lower bound; we carry the raw value throughout. The fixed-width convention gives a \emph{higher} value ($\approx1.9$) because it uses $0.10$\,pc rather than the larger measured widths; FilFinder gives lower values still. The full envelope factor $\approx2$--$13$ (Section~\ref{sec:injection_recovery}) is $\lambda_{\rm m}/(s_{\rm ACS,med}/W)$: its lower endpoint $\approx2$ uses $(\lambda_{\rm m},s_{\rm ACS,med}/W)=(5.5W,2.6)$ (fixed-width DisPerSE, Perseus) and its upper endpoint $\approx13$ uses $(5.5W,0.42)$ (densest FilFinder skeleton). We quote $s_{\rm ACS,med}/W$ to one decimal place; finer digits are not significant. The two \emph{deprojected} rows apply the median isotropic projection factor $\langle\sin i\rangle_{\rm med}=0.87$ (Section~\ref{sec:error_budget}); this correction acts \emph{against} our conclusion, raising the ratio by $15$--$27$ per cent, and we show it in the table rather than only in the text so that the effect is visible in the same place as the adopted numbers. We nevertheless carry the uncorrected value as the headline, because deprojection assumes an isotropic orientation distribution that the filaments in a single cloud need not obey.
\end{table}

\textbf{Primary sample.}\footnote{Throughout, \emph{primary} is a sample-selection label ($N>500$ cores, secure distance, and a skeleton simple enough for unambiguous core-to-filament association; Section~\ref{sec:distance_uncertainties}) for the regions retained for the adjacent-core-spacing measurement, and is \emph{not} a claim of statistical robustness.} We adopt the four regions Orion~B, Aquila, Perseus and Taurus ($N>500$ cores each, well-established distances) as the primary sample.

\textbf{Full sample result.} The core-count-weighted mean across all 8 HGBS regions ($0.279$\,pc) and the primary-only weighted mean ($0.290$\,pc) differ by $\approx3.9\%$, and a leave-one-out analysis confirms no single region dominates ($<6\%$ shifts).

All weighted means use core-count weighting,
\begin{equation}
\bar{\lambda} = \frac{\sum_i N_i\,\lambda_i}{\sum_i N_i},
\label{eq:weighting}
\end{equation}
where $\lambda_i$ and $N_i$ are the per-region spacing and core count.

The sub-classical result is a property of the prestellar population, not of protostellar contamination: the prestellar-only ACS median is $s_{\rm ACS,med}/W=2.10$, indistinguishable from the all-cores $2.08$ (Section~\ref{sec:nnpipeline}). The preferred ACS measurement follows in Section~\ref{sec:injection_recovery}.

\subsection{The dominant observational systematic: filament width}
\label{sec:widths}

Two width conventions recur throughout: the \emph{fixed-width} convention (the canonical HGBS $W_{\rm fil}=0.10$\,pc for every region) and the \emph{region-Plummer} convention (a beam-deconvolved Plummer FWHM fitted per region); since $s_{\rm ACS,med}/W\propto1/W$, the broader region-specific widths give a \emph{lower} $s_{\rm ACS,med}/W$. The fixed $W_{\rm fil}=0.10$\,pc convention is the largest single lever on $s_{\rm ACS,med}/W$, comparable to the whole discrepancy, so we measured region-specific widths directly. Along each spine we sampled the column-density profile perpendicular to the local tangent and fitted a beam-deconvolved Plummer function (the HGBS width definition; \citealt{Arzoumanian2019}) to hundreds of profiles per region ($18.2''$ beam; Table~\ref{tab:widths}). All four regions give $p\simeq2.0$ and $W_{\rm fil}=0.11$--$0.16$\,pc, agreeing with the published HGBS values and slightly broader than $0.10$\,pc, so the region-specific $s_{\rm ACS,med}/W$ ($1.1$--$1.8$, median $\approx1.4$) is \emph{lower} than the raw fixed-width value ($\approx2.1$; Table~\ref{tab:headline}); every region is sub-classical under either convention. A Gaussian-crest estimate returns a much narrower Taurus width ($0.031$\,pc). We reject it on physical grounds rather than because of the ratio it implies: a Gaussian fits only the inner ridge and ignores the power-law wings that carry most of the column density, whereas the Plummer form models the wings and is the profile used to define the HGBS filament width. The Plummer fit ($0.11$\,pc, $s_{\rm ACS,med}/W\approx1.1$) makes Taurus the most sub-classical region. The result is thus not conditional on the width convention. We prefer the Plummer widths on physical grounds: they fit the observed $p\simeq2$ profile of HGBS filaments and are far less sensitive to the fitting radius than either a single fixed $0.10$\,pc value or a Gaussian inner-ridge fit.

\begin{table}
\caption{Region-specific, beam-deconvolved filament widths from a direct Plummer fit to perpendicular column-density profiles along the skeletons (the HGBS width definition; \citealt{Arzoumanian2019}), and the resulting ACS $s_{\rm ACS,med}/W$. $W_{\rm fil}$ is the beam-deconvolved Plummer FWHM ($18.2''$ beam), $p$ the Plummer index (a \emph{free} fit parameter, not fixed; it converges to $p\simeq2.00$ in all four regions (Taurus $2.01$, a difference in the second decimal place only), consistent with the canonical Plummer-2 profile), $N$ the number of fitted profiles. Quoted $W_{\rm fil}$ uncertainties are the standard error on the median of the per-profile fits within each region and do not include the beam-deconvolution or background-subtraction systematics, which are carried separately in the error budget (Section~\ref{sec:error_budget}); the fitted $p$ has a per-region dispersion of $\pm0.1$ across profiles. All four regions give $p\simeq2$ and sub-classical $s_{\rm ACS,med}/W$ on the (primary-sample) \emph{median} spacing quoted here; the \emph{mean} adjacent spacing is a factor $\approx1.3$--$2.2$ larger (Section~\ref{sec:clustering}), so the mean-based ratio is $\approx2.1$--$4.7$, still sub-classical in every region.}
\label{tab:widths}
\footnotesize
\setlength{\tabcolsep}{3pt}
\begin{tabular}{p{1.7cm}ccccc}
\toprule
Region & $W_{\rm fil}$ (pc) & $p$ & $N$ & ACS (pc) & $s_{\rm ACS,med}/W$ \\
\midrule
Taurus  & $0.110$ & $2.01$ & 295 & $0.124$ & $1.1$ \\
Orion B & $0.148$ & $2.00$ & 356 & $0.207$ & $1.4$ \\
Perseus & $0.147$ & $2.00$ & 235 & $0.263$ & $1.8$ \\
Aquila  & $0.161$ & $2.00$ & 446 & $0.210$ & $1.3$ \\
\midrule
median  & $0.148$ & $2.00$ &, &, & $1.4$ \\
\bottomrule
\end{tabular}
\par\vspace{2pt}
\footnotesize The measured widths ($0.11$--$0.16$\,pc) agree with the published HGBS values and exceed the canonical $0.10$\,pc, so region-specific Plummer widths give a \emph{lower} $s_{\rm ACS,med}/W$ ($\approx1.4$, raw) than the raw fixed-width value ($\approx2.1$; Table~\ref{tab:headline}): all four regions are sub-classical under either convention. ACS spacings are the full-coverage medians (Section~\ref{sec:injection_recovery}), measured on the same DisPerSE skeleton as the widths (its deposited crest for the width fit, its bridged form for the spacing); the injection--recovery bias correction lowers each $s_{\rm ACS,med}/W$ by a further $\sim$10--20\%. (Pipeline outputs retained to 2 s.f.; the convention systematic dominates.)
\end{table}

The fit is consistent across the sample: every region returns a Plummer index $p\simeq2.0$, the value found for HGBS filaments generally \citep{Arzoumanian2019}, and a width within $\sim50\%$ of $0.10$\,pc (inter-quartile ranges $\sim0.07$--$0.20$\,pc), so the sub-classical $s_{\rm ACS,med}/W$ is a property of all four primary regions, not of one anomalous cloud. The single earlier outlier, a spuriously narrow $0.031$\,pc Gaussian crest for Taurus, was a profile-selection effect of fitting the inner ridge rather than the broad Plummer wings (not a pixel-scale artefact: the $3''$ pixels subtend $0.002$\,pc at $135$\,pc); it does not survive the Plummer fit, which gives $0.11$\,pc and $s_{\rm ACS,med}/W\approx1.1$.

\subsection{Sample Classification and Distance Robustness}
\label{sec:distance_uncertainties}

We classify the eight regions (Tables~\ref{tab:sample} and \ref{tab:sample_limited}) as \textbf{primary} (Orion~B, Aquila, Perseus, Taurus) or \textbf{limited} by three criteria fixed \emph{a priori} (before the measurement): $N>500$ cores; a well-established distance; and a skeleton simple enough for unambiguous core-to-filament association. This selection is cross-checked by the leave-one-out test below. The limited sample fails at least one criterion (Ophiuchus on skeleton complexity, the others on small samples) and are excluded from the primary ACS measurement. The largest distance revisions (Serpens $+76\%$, Aquila $+68\%$) are independently corroborated by VLBI maser parallaxes (validating Gaia DR3 to $<5\%$; \citealt{Kounkel2018,Xu2021,OrtizLeon2018,OrtizLeon2017}), extinction estimates \citep{Yan2022,Green2024} and co-spatiality with Aquila \citep{Zucker2020}; Serpens is \emph{limited} and excluded in any case. The sub-classical conclusion is independent of these regions: leave-one-out shifts the weighted mean by $<6\%$, and dropping Aquila leaves the region-Plummer median at $1.40$ and the fixed-width characteristic at $2.1$, both sub-classical. The primary/limited split is likewise not tuned to the distance revision: its criteria act on the distance-independent core count and skeleton morphology, the closest any region lies to the $N>500$ boundary is Taurus at $536$ cores ($\sim7\%$ above), and because the Gaia revision rescales spacings but not core counts, no plausible distance change moves a region across the classification.

\subsection{Choice of spacing statistic: the region-extent bias of the pairwise median}
\label{sec:statistics}

We adopt the along-skeleton adjacent-core spacing (ACS), denoted $s_{\rm ACS,med}$ for its per-region median, as the primary measure of the core spacing, reporting the all-pairs median only for comparability with prior HGBS work. This is \emph{not} the Euclidean nearest-neighbour distance, which we use only for the Clark--Evans statistic in Section~\ref{sec:clustering}. In figures whose axis reads $\lambda/W$ for continuity with prior work, the plotted observed quantity is $s_{\rm ACS,med}/W$, a line-density statistic, never a wavelength. The two differ fundamentally: for $N$ points uniformly distributed on a filament of length $L$ the pairwise median converges to $d^\star=L(1-1/\sqrt{2})\approx0.293\,L$ for uniform, extent-filling distributions (i.e.\ $\sim L/3$), the extent of the filament segment ($L_{\rm fil}\sim1$\,pc), for a random or a perfectly periodic distribution alike, so it measures the extent, not the local core spacing (triangular-distribution derivation, finite-$N$ Monte-Carlo and branched-skeleton generalisation in Appendix~\ref{sec:pwmedian}). It tracks $\sim L/3$ for any filament with $N\gtrsim5$ cores and never approaches the local spacing ($\sim30$--$100\times$ smaller). The ACS statistic is unaffected by this bias: the $18.2''$ HGBS beam corresponds to $0.012$--$0.039$\,pc, $\gtrsim6\times$ smaller than the measured ACS spacing, but because source extraction and deblending operate on a scale larger than one nominal beam this beam ratio is not by itself sufficient to establish the absence of resolution bias; the relevant argument is the injection--recovery test (Section~\ref{sec:injection_recovery}), which recovers the input spacings to a few per cent.

We validate this directly with synthetic point processes (Fig.~\ref{fig:nn_validation}). For core distributions built as periodic, clustered (gamma-renewal), hierarchical (three projected fibres) and fractal (power-law-gap) processes with a fixed injected local scale $\lambda_{\rm in}$, the ACS median is \emph{flat} with segment extent $L$ and recovers each process's local scale, whereas the pairwise median rises as $0.29\,L$ for \emph{every} process, independent of $\lambda_{\rm in}$; i.e.\ it estimates the segment extent, not the fragmentation scale. This confirms the analytic result across the full range of plausible core-placement statistics.

\begin{figure}
\centering
\includegraphics[width=0.99\columnwidth]{figures/fig_nn_validation.pdf}
\caption{Synthetic validation that the adjacent-core spacing (ACS) statistic recovers the injected local spacing scale while the pairwise median does not (a clustered process need not possess a fragmentation wavelength). Cores are placed on a segment of extent $L$ by four point processes (periodic, clustered, hierarchical three-fibre, fractal), each with a fixed injected local scale $\lambda_{\rm in}$. \emph{(a)} The ACS median (normalised by $\lambda_{\rm in}$) is flat with $L$ and recovers each process's local/apparent scale (the hierarchical case gives $\lambda_{\rm in}/3$, the projected value). \emph{(b)} The pairwise median instead rises along $0.293\,L$ (dashed; this is the median of the triangular pairwise-distance law, whose \emph{mean} is $L/3$; Appendix~\ref{sec:pwmedian}) for \emph{all four} processes, independent of $\lambda_{\rm in}$; confirming that it measures the segment extent, not the fragmentation wavelength.}
\label{fig:nn_validation}
\end{figure}

We reconstruct the along-skeleton core ordering for the four primary regions from the deposited DisPerSE skeleton maps, bridging fragmentation by morphological closing, and extend the ACS measurement to all four at full coverage in Section~\ref{sec:injection_recovery}; the limited sample is excluded from the ACS measurement owing to branching skeletons or small samples. Figure~\ref{fig:spacing} shows the per-region spacings with revised distances.

\begin{figure*}
\includegraphics[width=0.95\columnwidth]{figures/figure1_spacing_comparison.pdf}
\caption{The conditional local descriptor $S_{\rm local}=s_{\rm ACS,med}/W_{\rm fil}$ per primary region. This is \emph{not} the quantity classical theory predicts (that is $S_{\rm global}$; Fig.~\ref{fig:intermittency}b) and the comparison with the classical band below should be read diagnostically. It is the \emph{conditional} descriptor measured within core-bearing stretches; the unconditional length per core $L/N$, whose reciprocal is the line density predicted by fragmentation theory, is shown in Fig.~\ref{fig:intermittency}, in the two DisPerSE width conventions: directly-measured region-Plummer widths (blue circles, preferred; per-region raw $1.1,1.4,1.8,1.3$, dashed line at the median, which we quote as $1.4$, the rounded value of the exact median of $1.35$) and the fixed $0.10$\,pc width (orange squares; $1.2,2.1,2.6,2.1$), both shown \emph{raw} (before the injection--recovery correction; Section~\ref{sec:injection_recovery}). The red band is the classical fastest-growing wavelength $\lambda_{\rm m}\approx5$--$6\times$FWHM (\citealt{Fischera2012}; $=4\times$ the flat diameter, \citealt{Inutsuka1992}) and the dotted line the classical \emph{critical} wavelength $\lambda_{\rm cr}\approx2.6\times$FWHM; every region lies well below both. The grey band marks the independent FilFinder skeleton range ($\lambda/W\approx0.4$--$0.8$). (Uncertainty convention: Section~\ref{sec:error_budget}.)}
\label{fig:spacing}
\end{figure*}

\subsection{Injection--recovery validation and model dependence of the correction}
\label{sec:injection_recovery}

To test whether the ACS statistic recovers known spacings under realistic HGBS conditions, we placed synthetic cores with known periodic spacings (0.15--0.35\,pc) on HGBS-property-matched DisPerSE skeletons for all 8 regions. On clean skeletons the raw measurement recovers the input to about $3$\%; the larger correction actually applied is the \emph{full-coverage pipeline} bias ($\div1.11$--$1.29$, from junction/branch topology on the real bridged skeletons), which moves raw $2.1\to1.9$ and is a genuine, region-measured correction. We do \emph{not} additionally apply a point-process model factor (periodic $\div1.20$ vs clustered $\div1.64$): if the cores are clustered the measured ACS median already \emph{is} the physical line-density quantity, so correcting it toward a periodic template the data reject would be circular. We therefore carry the raw value as the adopted $(s_{\rm ACS,med}/W)_0$ and quote the clustered branch ($\div1.64$: raw $2.0\to1.2$) only as a lower bound. The sub-classical \emph{sign} holds under every choice.

Repeating the exercise for the pairwise median confirms it does not recover $\lambda_{\rm true}$ (it returns a value $\sim$$20\times$ larger, tracking $L/3$), whereas the ACS statistic recovers $\lambda_{\rm true}$ to $1.8$--$3\%$ on the identical skeletons (Appendix~\ref{sec:pwmedian}).

\textbf{Full-coverage along-skeleton ACS.} We extend the ACS measurement to all four primary regions using the dense deposited DisPerSE skeletons, which (unlike the sparse FilFinder spines) pass near essentially all cores. Bridging the fragmented skeletons by morphological closing plus 1-pixel re-skeletonization, ordering cores by graph arclength and associating within a half-width ($\sim$0.06\,pc), we obtain $N_{\rm pairs}=461,614,563,168$ spacings (Taurus, Orion~B, Perseus, Aquila), median ACS $0.124,0.207,0.263,0.210$\,pc, i.e.\ $s_{\rm ACS,med}/W=1.2,2.1,2.6,2.1$ (characteristic $\approx2.1$; robust to the bridging/association radii). The periodic pipeline factor gives $s_{\rm ACS,med}/W\approx1.0,1.9,2.1,1.8$ (characteristic $\approx1.9$). Applied to every width convention (Section~\ref{sec:widths}), the fixed-width characteristic is $s_{\rm ACS,med}/W\approx1.2$--$1.9$ and the region-Plummer $\approx0.8$--$1.4$; the primary observed result is this DisPerSE envelope $s_{\rm ACS,med}/W\simeq0.8$--$1.9$ (wider under other filament definitions; Fig.~\ref{fig:envelope}), within which the raw region-Plummer $(s_{\rm ACS,med}/W)_0\approx1.4$ is a bookkeeping reference convention (not the observed value) carried with the correction as a uniform systematic. \textbf{The full-coverage ACS spacing ($s_{\rm ACS,med}/W\approx1.2$--$2.1$), which is conditional on being inside a core-bearing stretch, is sub-classical by a factor $\approx2.4$--$4.3$ across the DisPerSE conventions ($\approx3.6$--$4.3$ for the region-Plummer reference alone; more under FilFinder), while the pairwise-median ratio $\approx2.8$ (in $W$ units) overestimates it by $\sim$45--50\% through the $L/3$ bias.} Taurus is the most sub-classical region, not an outlier. The two high-coverage clouds (Perseus, Taurus; $76/89\%$ core--skeleton association) alone give the same median ($s_{\rm ACS,med}/W=1.8$ and $1.1$, median $1.4$), so the result is not driven by the low-coverage clouds.

As a segmentation/ontology \emph{sensitivity} test rather than an independent confirmation, it uses the same \textit{Herschel} maps and source catalogue, so it is not an independent dataset, only an independent way of defining the filament network, running the same along-skeleton BFS-arclength ACS pipeline on independently extracted FilFinder spine maps (all four primary regions) gives $s_{\rm ACS,med}/W\approx0.4/0.6/0.7/0.8$ for Taurus, Orion~B, Perseus, Aquila respectively, systematically $2$--$3\times$ smaller than the DisPerSE values ($1.0/1.9/2.1/1.8$) because the FilFinder mask threshold produces a denser, more branched network. Both methods agree the characteristic spacing is far sub-classical ($s_{\rm ACS,med}/W\ll4$) in every region; FilFinder if anything strengthens the conclusion. The injection--recovery estimator test (Section~\ref{sec:injection_recovery_model}) was run on the DisPerSE skeletons; for FilFinder, its denser, more-branched topology (Fig.~\ref{fig:skeleton_compare}) is precisely why it returns shorter spacings, so we carry the FilFinder value only as the extreme of the convention envelope, not as a load-bearing number; a matched FilFinder injection--recovery, reported in the next paragraph, confirms the shorter values are genuine rather than an estimator artefact.

\smallskip\noindent\textbf{Model dependence of the correction.}\label{sec:injection_recovery_model}
Deriving for each region a correction matched to its own measured clustering (anchored to the periodic $\div1.20$ and Poisson $\div1.64$ endpoints) gives region-matched factors $\div1.40$--$1.86$ and corrected $s_{\rm ACS,med}/W\approx0.7$--$1.6$ (fixed $W=0.10$\,pc). The effect on the characteristic is small (mean $s_{\rm ACS,med}/W$ $1.2\to1.3$): the sub-classical result holds under each region's own matched correction, and $s_{\rm ACS,med}/W\approx1.9$ (which requires the rejected periodic model) remains excluded.

As a matched, independent-segmentation test, we ran the same synthetic injection--recovery on the deposited FilFinder spines: cores placed at known arclength spacings ($0.15$--$0.35$\,pc) are recovered to $<1\%$ (mean $|{\rm err}|=0.2$--$1.0\%$ across the four clouds). The shorter FilFinder $s_{\rm ACS,med}/W$ is therefore \emph{not} an estimator or branching artefact; it reflects that FilFinder traces a denser, more-branched network; a genuinely different definition of ``the filament'' (Section~\ref{sec:whatfilament}), not a measurement bias.

\subsection{Core line density, the local-width test, and the coverage bias}
\label{sec:linedensity}
The median adjacent spacing measures the spacing \emph{where cores are adjacent}; the quantity a one-fragment-per-wavelength theory actually predicts is the core \emph{line density}, cores per unit filament length. We therefore report three complementary statistics (Fig.~\ref{fig:linedensity}), computed in a single independent re-implementation and all normalised by the \emph{same} fixed $0.10$\,pc width so that they are directly comparable (their absolute values differ from the adopted pipeline by less than the convention systematic; Section~\ref{sec:error_budget}). (i)~The median adjacent (along-skeleton adjacent-core) spacing is $s_{\rm ACS,med}/W\approx1.7$--$3.0$, sub-classical in every cloud; these single-re-implementation fixed-width medians reconcile with the adopted-pipeline fixed-width per-region values ($1.2$--$2.6$; Table~\ref{tab:headline}) to within the convention systematic (Taurus differs most, $\sim40\%$). (ii)~The mean adjacent spacing is a factor $1.3$--$2.2$ larger ($\approx2.1$--$4.7$), but on the heavily-bridged skeletons it is inflated by arclength-ordering jumps across branches (the same effect that lifts Perseus's ACS above its pairwise median; Section~\ref{sec:statistics}), so we treat it as an upper descriptor, not the line density. (iii)~The direct core line density is $\rho_{\rm core,1D}=N/L$, the associated cores divided by the total skeleton arclength. For comparison with a wavelength we quote its reciprocal $L/N$, the mean filament length per associated core; this is ordering-independent and is the correct quantity to set against a one-fragment-per-wavelength prediction: $L/N=0.34$--$0.58$\,pc, i.e.\ $(L/N)/W\approx3.4$ (Orion~B), $3.4$ (Aquila), $4.5$ (Taurus), $5.8$ (Perseus). It is a factor $\approx2$ larger than the median adjacent spacing in every cloud. \textbf{On the global length per core the deficit is therefore much milder: Perseus is essentially classical and Taurus near-classical, while Orion~B and Aquila remain sub-classical at $\approx3.4$.} The two contributions to this near-classical global value can be separated, and neither is small-scale clustering. The median$\to$mean rise ($\approx40$--$120\%$) is the \emph{generic mean/median skew} of a broad gap distribution (a pure exponential already gives $1.44$; Section~\ref{sec:clustering}), not clustering; the further mean$\to L/N$ rise is \emph{dilution by core-poor filament}; extensive skeleton stretches with no cores (a coverage effect). The line density restricted to core-bearing stretches (the mean adjacent spacing, $\approx2.1$--$4.7\,W$) is intermediate and still mildly sub-classical, so the globally near-classical $L/N$ is \emph{not} evidence of a uniform classical spacing: it is a spatially inhomogeneous distribution with sub-classical local spacing and an observational small-separation deficit, whose filament-averaged density is lifted toward classical by the skew and by long core-free filament. The median adjacent ratio adopted elsewhere is accordingly a \emph{descriptor} of the local spacing within core-bearing stretches; $L/N$ (core-bearing or full) is the quantity to compare with a one-fragment-per-$\lambda_{\rm m}$ theory.

Two further tests support this reading. First, a \emph{local-width} normalisation (a genuinely like-for-like test): for every adjacent-core interval we fit the beam-deconvolved transverse width at the interval midpoint and form the per-cloud distribution of $\Delta s/W_{\rm local}$ ($146$--$545$ reliable fits per cloud). The medians are $0.61,0.91,1.05,0.99$ (Taurus, Orion~B, Perseus, Aquila), \emph{all} sub-classical, with only a weak spacing--width correlation ($r=-0.06$ to $0.31$). These median-of-ratios values are if anything \emph{more} sub-classical than the corresponding ratio-of-medians ($\approx1.4$--$1.9$ in the skewed clouds), differing by up to a factor $\sim2$ in Taurus and Perseus; the per-segment normalisation therefore \emph{strengthens} the sub-classical conclusion rather than depending on the two agreeing. Second, the along-skeleton \emph{pair-correlation} function $g(s)$ (Fig.~\ref{fig:paircorr}) shows \emph{no} peak at the classical scale $5$--$6\,W$ ($\approx0.5$--$0.8$\,pc): $g\approx1$ there (slightly below unity in Orion~B and Aquila), with a mild small-separation excess in Taurus and weak region-envelope features near $1$\,pc. To decide whether the absence of a feature is significant rather than merely visual, we generated a $95$ per cent envelope for $g(s)$ from the same column-density-conditioned null (specified in full in Appendix~\ref{sec:nullspec}). At the classical scale the observed $g$ lies \emph{inside} that envelope in every cloud (Orion~B $1.02$ within $[0.84,1.26]$; Aquila $0.89$ within $[0.58,1.52]$; Perseus $1.02$ within $[0.91,1.14]$; Taurus $1.00$ within $[0.91,1.10]$), so no classical-scale periodicity is detected at the $95$ per cent level. We phrase the result that way rather than as an absolute absence.

\begin{figure*}
\centering
\includegraphics[width=0.94\textwidth]{figures/fig_scaletests.pdf}
\caption{Scale and convention tests of the inhomogeneity result. \emph{(a)} Occupancy fraction against the adopted half-width $\ell$ (dotted line: the $0.06$\,pc value quoted in the text). Taurus and Perseus saturate, so their low occupancy is not a consequence of the chosen $\ell$; Orion~B and Aquila rise steadily and are less strongly segregated. \emph{(b)} Fano factor against cell length. Over-dispersion ($F>1$, dashed) persists across $0.25$--$4$\,pc in Taurus, Perseus and Orion~B, but Aquila is consistent with Poisson below $1$\,pc. \emph{(c)} Convergence against association half-width: the conditional $s_{\rm ACS,med}/W$ (solid) plateaus, whereas the unconditional $(L/N)/W$ (dashed) falls as more cores are admitted, approaching the classical band (shaded) only in Perseus and Taurus.}
\label{fig:scaletests}
\end{figure*}

\subsection{Is the whole skeleton eligible to fragment?}
\label{sec:eligible}
Identifying $L/N$ with the quantity a one-fragment-per-wavelength theory predicts carries a physical assumption that we have so far left implicit: that every metre of filament counted in $L$ was capable of taking part in the same fragmentation process. In a paper whose central observation is that much of the skeleton carries no cores, that assumption deserves testing rather than asserting. Classical fragmentation does not predict one core per $\lambda_{\rm m}$ along subcritical filament; including such material in $L$ inflates $L/N$ towards, and past, the classical value for a purely geometric reason.

We therefore define eligibility \emph{independently of where the cores are}, using column density rather than proximity to a core. A filament of width $0.10$\,pc reaches the critical line mass $M_{\rm line,crit}=2c_s^2/G\simeq16\,M_\odot\,{\rm pc^{-1}}$ at $N_{\rm H_2}\simeq7.2\times10^{21}\,{\rm cm^{-2}}$ (taking $\mu=2.8$), the familiar HGBS supercriticality threshold, and we take the \emph{column-density-selected} eligible length $L_{\rm elig}$ to be the skeleton length above it. This is not a direct measurement of the local line mass, which would require deprojection, background subtraction and a local width; we use column density as its observable proxy and treat the threshold as uncertain rather than exact. The eligible fractions differ greatly between clouds: $0.19$ (Orion~B), $0.91$ (Aquila), $0.13$ (Perseus) and $0.07$ (Taurus). Orion~B in particular is dominated by filament that could not have fragmented at all.

Restricting to that material changes the conclusion, but the size of the change depends on a subtlety that we must deal with first. Computing both lengths on the same edge-summed convention used in Section~\ref{sec:linedensity}, the raw calculation gives $(L_{\rm elig}/N_{\rm elig})/W=1.6$, $3.2$, $2.0$ and $0.8$ for Orion~B, Aquila, Perseus and Taurus, against $3.4$, $3.4$, $5.8$ and $4.5$ over the whole skeleton.

\smallskip\noindent\textbf{The criterion is partly endogenous, and correcting for it matters.} A core raises the column density of its own surroundings, so a threshold applied to the raw map can be pushed over the line \emph{by the core it is meant to be independent of}. That selects short eligible patches centred on cores, depressing $L_{\rm elig}/N_{\rm elig}$ in exactly the direction we reported. Between $25$ and $34$ per cent of skeleton pixels lie within one association half-width of a core, so this is not a small effect. We therefore rebuilt the criterion on a \emph{core-masked background}: pixels within $0.06$\,pc of an associated core have their column density replaced by interpolation along the spine from untainted pixels, and the threshold is applied to that reconstruction.

The result is a substantial weakening of our earlier claim, and we report it as such. On the core-masked background, $(L_{\rm elig}/N_{\rm elig})/W$ becomes $2.0$, $3.2$, $4.1$ and $3.1$, against the raw $1.6$, $3.2$, $2.0$ and $0.8$. Aquila shows no change larger than the mode-quantisation floor, because almost all of its skeleton is eligible either way, but Perseus doubles and Taurus rises by a factor of nearly four. All four clouds remain below the classical $5$--$6$, so the sign of the result survives; but the margin is much reduced, and the strong statement that eligible filament is decisively sub-classical in every cloud does not survive. We adopt the core-masked values as fiducial and quote the raw values only as the endogeneity-affected comparison.

\smallskip\noindent\textbf{The masking radius is itself a choice, and it matters in two clouds.} The reconstruction above uses the same $0.06$\,pc half-width adopted for core association, so we repeated it over $0.03$--$0.12$\,pc (Table~\ref{tab:maskrad}). In Orion~B and Aquila the result is stable: $S_{\rm elig}$ varies only between $1.89$ and $2.04$ and between $3.14$ and $3.21$ respectively, and $N_{\rm elig}$ stays above $230$. In Perseus and Taurus it is not. Perseus moves over $3.6$--$4.5$ while $N_{\rm elig}$ falls from $85$ to $37$, and Taurus is worse still: $S_{\rm elig}=1.8$, $1.6$ and $3.1$ at masking radii of $0.03$, $0.045$ and $0.06$\,pc, and beyond $0.09$\,pc a single core survives and the statistic ceases to exist. The Taurus value quoted above is therefore an artefact of one particular masking scale and we withdraw it; the Perseus value is radius-dependent at the tens-of-per-cent level but stays sub-classical throughout. This is the sensitivity test a referee is entitled to demand of a correction of this kind, and it changes what we are willing to claim: the core-masked eligibility comparison is supported in Orion~B and Aquila, indicative in Perseus, and not available in Taurus.

\begin{table}
\caption{Sensitivity of the core-masked eligibility statistic to the masking radius $r_{\rm mask}$ (the distance from an associated core within which skeleton column density is replaced by along-spine interpolation). Each cell gives $S_{\rm elig}=(L_{\rm elig}/N_{\rm elig})/W$ with $N_{\rm elig}$ in brackets. The first row, $r_{\rm mask}=0$, is the uncorrected raw-map calculation. The adopted value is $0.06$\,pc. Orion~B and Aquila are stable; Perseus drifts; Taurus is not usable at any radius, its apparent value being set by which handful of cores survives.}
\label{tab:maskrad}
\footnotesize
\setlength{\tabcolsep}{4pt}
\begin{tabular}{lcccc}
\toprule
$r_{\rm mask}$ (pc) & Orion~B & Aquila & Perseus & Taurus \\
\midrule
$0$ (raw)      & 1.56 [417] & 3.14 [306] & 1.96 [199] & 0.81 [186] \\
$0.030$        & 1.89 [317] & 3.14 [306] & 3.68 [85]  & 1.80 [43] \\
$0.045$        & 2.04 [277] & 3.15 [305] & 3.91 [70]  & 1.61 [37] \\
$0.060$        & \textbf{2.02} [268] & \textbf{3.17} [303] & \textbf{4.13} [57] & (3.11) [13] \\
$0.090$        & 1.98 [252] & 3.21 [296] & 4.45 [41]  & --- [1] \\
$0.120$        & 1.95 [238] & 3.20 [297] & 3.62 [37]  & --- [1] \\
\bottomrule
\end{tabular}
\end{table}

Two of the core-masked numbers additionally rest on small samples. Masking reduces $N_{\rm elig}$ from $199$ to $57$ in Perseus and from $186$ to $13$ in Taurus, because little of the \emph{background} filament in those clouds is supercritical once the cores' own contribution is removed. With thirteen cores the Taurus value is not a usable statistic and we do not rely on it. The honest summary is that Orion~B ($2.0$, $268$ cores) and Aquila ($3.2$, $303$ cores) support a sub-classical eligible line density, that Perseus is indicative at $\approx3.6$--$4.5$ depending on the masking radius, and that Taurus cannot be used at all.

\smallskip\noindent\textbf{What distinguishes a fertile stretch from a sterile one?} Our own test answers this only partly. Column density above the critical line mass is clearly a necessary condition, and it is not a sufficient one: in Aquila $90$\,per cent of the crest is eligible on this criterion, yet only $24$\,per cent of it carries a core. Something beyond the thermal line mass therefore decides which eligible stretches actually fragment. Three candidates are available and none is testable with column density alone. Local compression, from an expanding \ion{H}{ii} region or a converging flow, can raise the line mass above threshold on a timescale shorter than the local free-fall time and so seed fragmentation preferentially in one part of a cloud; the strongly one-sided occupancy of Orion~B, which lies beside NGC\,2024 and NGC\,2023, is at least consistent with this. Longitudinal velocity gradients, which are observed along HGBS filaments at $\sim1$\,km\,s$^{-1}$\,pc$^{-1}$ \citep{Kirk2013,Hacar2013}, can either assemble material into hubs or shear a stretch faster than it can fragment. And non-thermal support varies along a filament, so a stretch that is thermally supercritical may still be virially stable; our transonic sensitivity row is a crude, uniform version of this. Deciding among them needs resolved kinematics per stretch, which is exactly the observation we lack, and we therefore report the occupancy as a measured property rather than an explained one. Because the threshold depends on an assumed width and temperature, we tested its sensitivity rather than adopting one value (Table~\ref{tab:eligible}). Raising the threshold from $4\times10^{21}$ to $1.2\times10^{22}\,{\rm cm^{-2}}$ moves $(L_{\rm elig}/N_{\rm elig})/W$ monotonically \emph{downwards}, from $1.9$, $3.4$, $3.5$, $1.9$ to $1.5$, $2.0$, $1.4$, $0.6$, so a stricter eligibility criterion makes the result more sub-classical and our fiducial choice is the conservative one. Using each region's own measured Plummer width instead of $0.10$\,pc gives $1.7$, $3.4$, $2.1$ and $0.8$, and adopting $T=15$\,K with $W=0.13$\,pc gives $1.6$, $2.8$, $1.9$ and $0.8$. Across every variant the eligible length per core lies in the range $0.6$--$3.5$, well below the classical $5$--$6$.

A further objection is that we have used the purely thermal Ostriker line mass while invoking magnetic and turbulent support elsewhere in this paper. A filament can be thermally supercritical yet still supported. We therefore repeated the calculation with a virial line mass, $M_{\rm line,vir}=2(c_s^2+\sigma_{\rm nt}^2)/G$, which raises the threshold to $1.5\times10^{22}\,{\rm cm^{-2}}$ for transonic non-thermal support ($\sigma_{\rm nt}=c_s$) and $3.7\times10^{22}$ for $\sigma_{\rm nt}=2c_s$. Treating $\sigma_{\rm nt}$ as isotropic pressure support is a prescription, not a measurement: observed non-thermal linewidths may contain infall, longitudinal flows or unresolved components, some of which promote rather than resist fragmentation, and we have no local velocity dispersions of our own. As an illustrative sensitivity test the effect runs in the direction that strengthens our conclusion: the eligible length shrinks and the length per core falls further, to $1.4$, $1.8$, $1.3$ and $0.5$ in the transonic case against $1.6$, $3.1$, $2.0$ and $0.8$ for thermal support alone. The thermal-only criterion is therefore the \emph{conservative} choice, which is why we adopt it. We do not push beyond the transonic case: at $\sigma_{\rm nt}=2c_s$ only $0$--$5$\,pc of filament remains eligible, with $0$--$47$ cores, and Taurus retains none at all, so the statistic degenerates. Even the transonic row of Table~\ref{tab:eligible} rests on $49$--$205$ cores and, in Perseus and Taurus, on a few per cent of the skeleton; those two entries should be read as indicative only. We have no per-region velocity dispersions of our own and treat the non-thermal correction as a sensitivity band rather than a measurement.

Taken together, these tests qualify rather than overturn the global result. Part of the near-classical global line density does come from counting filament that was never eligible to fragment, most severely in Orion~B where four fifths of the skeleton is subcritical; but part of the apparent sub-classical eligible spacing came from the cores themselves raising the column density used to select the eligible length. With the endogeneity removed, the eligible statistic lies between the local and the global one rather than coinciding with the local one, and in the two clouds where it is well sampled it remains below the classical band. We therefore treat Section~\ref{sec:eligible} as showing that the local/global tension is partly, not wholly, a definition effect, and as a demonstration of how sensitive the comparison is to what one is willing to call a filament. Because the eligible length depends on an assumed width and temperature and on region-specific kinematics we do not possess, we regard these numbers as provisional pending resolved velocity dispersions per region.

\begin{table}
\caption{Sensitivity of the fragmentation-eligible length per core, $(L_{\rm elig}/N_{\rm elig})/W$, to the adopted eligibility criterion. Each cell gives that ratio followed, in brackets, by $f_{\rm elig}$, the fraction of total skeleton length passing the criterion, and $N_{\rm elig}$, the number of cores lying on the eligible length. Rows labelled by a column density apply that threshold in $N_{\rm H_2}$ (cm$^{-2}$) to the skeleton; both lengths are edge-summed (Section~\ref{sec:linedensity}) and $W$ is the fixed $0.10$\,pc width unless stated. The fiducial row applies the threshold to a \emph{core-masked} reconstruction of the column density (pixels within $0.06$\,pc of an associated core replaced by along-spine interpolation), which removes the cores' own contribution to the criterion; the row above it shows the same threshold on the raw map for comparison. Perseus and Taurus retain only $57$ and $13$ eligible cores after masking and are not usable. Among raw-map thresholds the trend is monotonic, a stricter threshold giving a more sub-classical result. ``Local $W$'' uses each region's measured Plummer width at $T=10$\,K; the virial row uses $M_{\rm line,vir}=2(c_s^2+\sigma_{\rm nt}^2)/G$ as an illustrative effective-support prescription. All rows other than the fiducial use the raw map.}
\label{tab:eligible}
\footnotesize
\setlength{\tabcolsep}{3pt}
\begin{tabular}{lcccc}
\toprule
Criterion & Orion B & Aquila & Perseus & Taurus \\
 & \multicolumn{4}{c}{$(L_{\rm elig}/N_{\rm elig})/W$ \ [$f_{\rm elig}$, $N_{\rm elig}$]} \\
\midrule
$4\times10^{21}$      & 1.9 [0.43, 761] & 3.4 [1.00, 308] & 3.5 [0.52, 431] & 1.9 [0.36, 400] \\
$6\times10^{21}$      & 1.7 [0.26, 516] & 3.4 [1.00, 308] & 2.3 [0.21, 259] & 1.0 [0.12, 254] \\
$7.15\times10^{21}$ (raw map)  & 1.6 [0.19, 417] & 3.2 [0.91, 306] & 2.0 [0.13, 199] & 0.8 [0.07, 186] \\
$7.15\times10^{21}$, core-masked (fiducial) & \textbf{2.0} [0.16, 268] & \textbf{3.2} [0.90, 303] & \textbf{4.1} [0.08, 57] & \textbf{3.1} [0.02, 13] \\
$9\times10^{21}$      & 1.5 [0.15, 326] & 2.7 [0.75, 295] & 1.8 [0.09, 136] & 0.7 [0.04, 132] \\
$1.2\times10^{22}$    & 1.5 [0.10, 214] & 2.0 [0.48, 256] & 1.4 [0.04, 89] & 0.6 [0.02, 79] \\
\midrule
Local $W$, $T=10$\,K  & 1.7 [0.33, 640] & 3.4 [0.99, 308] & 2.1 [0.16, 222] & 0.8 [0.05, 141] \\
$T=15$\,K, $W=0.13$\,pc & 1.6 [0.16, 339] & 2.8 [0.81, 302] & 1.9 [0.10, 152] & 0.8 [0.05, 141] \\
\midrule
Virial, transonic $\sigma_{\rm nt}=c_s$ & 1.4 [0.07, 166] & 1.8 [0.34, 205] & 1.3 [0.03, 67] & 0.5 [0.01, 49] \\
\midrule
Whole skeleton (no cut) & 3.4 & 3.4 & 5.8 & 4.5 \\
\bottomrule
\end{tabular}
\end{table}

\begin{figure}
\centering
\includegraphics[width=0.99\columnwidth]{figures/fig_skelzoom.pdf}
\caption{Close-up of the highest-column hub in Orion~B, showing where the two algorithms actually differ. DisPerSE (blue) tracks the visible crests. FilFinder (red) reproduces the main ridge but adds long, nearly straight spokes that radiate from the hub and do not follow the greyscale emission: these are medial-axis artefacts of a large connected mask rather than resolved sub-filaments. The extra FilFinder length is therefore not obviously tracing real hierarchical structure, which is why we carry the FilFinder branch as the extreme of the convention envelope rather than as an equally-weighted alternative. We stop short of calling FilFinder wrong: a medial axis of a column-density mask is a legitimate and physically motivated definition of a filament network, and deciding which network is real requires velocity information (position--position--velocity data resolving coherent sub-structure), which column density alone cannot supply.}
\label{fig:skelzoom}
\end{figure}

\subsection{Robustness to the persistence threshold}
\label{sec:persistence}

The DisPerSE skeletons we use are the published HGBS products, extracted at a $3\sigma$ persistence cut. Since Section~\ref{sec:error_budget} identifies the skeleton algorithm as the largest single systematic, the persistence threshold of our primary ontology cannot be left untested. We cannot regenerate the published skeletons at a \emph{lower} threshold without the original DisPerSE run, so we approach the question from both sides.

\smallskip\noindent\emph{(i) Pruning the published skeleton upwards.} We decompose each skeleton into branches, terminating at junction pixels, and compute for each branch the persistence of its column-density maximum above the junction at which it attaches, in units of the local map noise $\sigma$ (estimated as $1.4826\times$MAD of the map minus a $5$-pixel median filter). Branches below a chosen threshold are removed and the procedure iterated to convergence. Raising the effective threshold from the published $3\sigma$ to $4\sigma$ removes $16$, $21$, $18$ and $15$ per cent of the skeleton length in Orion~B, Aquila, Perseus and Taurus, but changes the median adjacent spacing by only $-8.2$, $-2.7$, $+0.3$ and $-5.0$ per cent and the length per core by $-1.9$, $-3.2$, $-3.1$ and $-3.7$ per cent. Going further, to $6\sigma$ and $8\sigma$, changes nothing further at the per-cent level: the pruning saturates once the low-persistence spurs are gone. The occupancy fraction rises slightly, as expected when bare low-persistence branches are removed.

\smallskip\noindent\emph{(ii) An independent extraction with a tunable threshold.} To probe the $2\sigma$ direction we re-implemented the persistent-crest construction from scratch: a descending-sweep merge tree of the smoothed column-density field pairs every maximum with the saddle at which it merges, maxima whose persistence falls below the chosen multiple of $\sigma$ are cancelled, and the surviving crests are traced by steepest-ascent walks from the retained saddles. This is an independent code, not DisPerSE, and it does not reproduce the published skeleton pixel-for-pixel; what it provides is a skeleton whose persistence threshold is a free parameter. Run at $2\sigma$, $3\sigma$ and $4\sigma$ on Orion~B it gives $S_{\rm local}=1.63$, $1.63$ and $1.63$ and $S_{\rm global}=2.53$, $2.50$ and $2.48$, with the associated core count falling from $635$ to $597$ and the occupancy fraction rising from $0.41$ to $0.42$. The local statistic is thus insensitive to the persistence threshold at the third significant figure over a factor of two in threshold, and the global statistic moves by two per cent.

\smallskip\noindent\emph{What this does and does not settle.} Both tests say the same thing: the persistence threshold is not what makes the DisPerSE and FilFinder answers differ. Between $2\sigma$ and $8\sigma$ the local spacing moves by less than ten per cent, whereas changing the ontology from persistent crests to the medial axis of a mask moves it by a factor of two. The dominant systematic is therefore the \emph{definition} of a filament, not the tuning of a parameter within one definition, which is the conclusion Section~\ref{sec:whatfilament} draws and which the tests here support rather than assume. What remains untested is the manual editing of the published skeletons, which we did not perform and cannot undo.

\subsection{Does the local/global contrast survive changing the definition of a filament?}
\label{sec:ontology}
Because a skeleton defines $L$ itself, changing the extraction algorithm alters not only the spacing but $L$ itself, and hence $N/L$, the core association and the topology. DisPerSE and FilFinder are therefore not two noisy estimates of one observable; they are partly different definitions of the object. We accordingly recomputed \emph{all} of the principal statistics under both (Table~\ref{tab:ontology}), rather than the adjacent spacing alone.

Inspecting where the two disagree is instructive (Fig.~\ref{fig:skelzoom}). At a high-column hub, FilFinder recovers the main ridge but adds long, almost straight spokes that leave the emission entirely; these are most likely medial-axis artefacts of a large connected mask rather than resolved sub-filaments. We have not quantified what fraction of the total FilFinder length behaves this way, so this remains a characterisation from one close-up rather than a measured classification. Some of the extra FilFinder length is therefore not real structure, and we treat that branch as the extreme of the convention envelope rather than as an equally plausible ontology. The FilFinder networks are shorter ($67$--$121$\,pc against $106$--$337$\,pc) and associate far fewer catalogue cores, so both statistics move, but they move in \emph{opposite} directions. The conditional spacing becomes shorter still ($s_{\rm ACS,med}/W=0.54$--$1.54$, against $1.66$--$2.97$ for DisPerSE), while the unconditional length per core becomes much larger ($(L/N)/W=9.1$--$16.7$, against $3.4$--$5.8$). The contrast between the local and global statistics is thus not an artefact of the DisPerSE ontology: it survives the change and in fact widens, from a factor $\approx2$ under DisPerSE to $\approx6$--$17$ under FilFinder. The occupancy fraction stays low under both ($0.05$--$0.09$ for FilFinder). We regard this as one of the more secure results in the paper, precisely because the two algorithms disagree so strongly about what a filament is.

\begin{table}
\caption{Principal statistics under both filament definitions, at a $0.06$\,pc association half-width and a $0.06$\,pc occupancy half-width. $F(1\,{\rm pc})$ is omitted where the FilFinder components are mostly shorter than the cell. The local and global statistics move in opposite senses between the two algorithms, so their contrast is not an artefact of either ontology.}
\label{tab:ontology}
\footnotesize
\setlength{\tabcolsep}{3pt}
\begin{tabular}{llccccc}
\toprule
Region & Skeleton & $L$ & $N_{\rm assoc}$ & $(L/N)/W$ & $s_{\rm ACS,med}/W$ & $f_{\rm occ}$ \\
 & & (pc) & & & & \\
\midrule
Orion B & DisPerSE  & 337 & 995 & 3.4 & 1.77 & 0.27 \\
        & FilFinder & 121 &  73 & 16.6 & 1.54 & 0.06 \\
Aquila  & DisPerSE  & 106 & 308 & 3.4 & 1.66 & 0.27 \\
        & FilFinder &  67 &  40 & 16.7 & 0.97 & 0.05 \\
Perseus & DisPerSE  & 290 & 501 & 5.8 & 2.97 & 0.12 \\
        & FilFinder &  78 &  85 &  9.1 & 1.05 & 0.09 \\
Taurus  & DisPerSE  & 209 & 463 & 4.5 & 1.81 & 0.06 \\
        & FilFinder & 111 & 120 &  9.3 & 0.54 & 0.05 \\
\bottomrule
\end{tabular}
\end{table}

Finally we quantify the skeleton-coverage bias. Only a fraction of catalogue cores associate with the primary skeleton ($24$--$89\%$ across the four clouds; Table~\ref{tab:nnpipeline}), and the omission is not neutral: the rejected cores have a $2$-D projected core-to-core spacing $1.5$--$3.0\times$ \emph{larger} than the associated cores in every cloud, so including them would lengthen the measured spacing. Coverage selection therefore biases the adjacent spacing \emph{short}, in the direction of the sub-classical result. We now \emph{bound} it: relaxing the association half-width from $0.06$ to $0.30$\,pc adds the rejected cores back, raising the coverage to $70$--$97\%$, and lengthens the median adjacent spacing by only $+3\%$ (Taurus), $+8\%$ (Perseus), $+11\%$ (Orion~B) and $+16\%$ (Aquila); it stays sub-classical throughout. Rather than leave this as a set of separate robustness statements, we ran it as a convergence experiment, tracking every principal statistic against the association half-width from $0.04$ to $0.30$\,pc (Fig.~\ref{fig:scaletests}c). The conditional spacing is remarkably stable: $s_{\rm ACS,med}/W$ varies only over $1.7$--$2.0$ (Orion~B), $1.6$--$1.9$ (Aquila), $3.0$--$3.2$ (Perseus) and $1.8$--$1.9$ (Taurus) while the associated fraction rises from $0.39$--$0.77$ to $0.70$--$0.97$. It plateaus, and the coverage bias on it is one-sided and bounded at $\lesssim16\%$, which we carry as an explicit systematic (Table~\ref{tab:errorbudget}).

The unconditional statistic behaves differently, and we flag this as a genuine limitation rather than a robustness success. Because adding cores at fixed $L$ necessarily raises the line density, $(L/N)/W$ falls monotonically with association radius: from $3.6$ to $2.1$ in Orion~B, $3.6$ to $2.0$ in Aquila, $6.6$ to $4.2$ in Perseus and $5.1$ to $4.0$ in Taurus. The near-classical global line density is therefore robust in Perseus and Taurus, which remain at $\approx4$ even at $85$--$97\%$ coverage, but in Orion~B and Aquila it holds only near the adopted radius and becomes sub-classical ($\approx2$) once coverage is pushed to $70$--$85\%$. The local/global contrast persists in all four clouds, but its \emph{size} in the two low-coverage clouds depends on how generously cores are assigned to a skeleton, and we do not claim a precise value for it there. (The coverage fractions, three-tier statistics and this bound come from the independent re-implementation at its own association radius; they differ from the adopted-pipeline values by less than the convention systematic, and the $\lesssim16\%$ bound is insensitive to the radius choice.) A cleaner per-segment analysis (splitting each skeleton at junctions and measuring unbranched segments) was attempted but the morphologically-bridged skeletons over-fragment at spurious junctions, so a physically meaningful segmentation would itself require collinear-segment merging, the same ``what is a filament'' ambiguity discussed in Section~\ref{sec:whatfilament}. We defer it; the coverage bound above and the three-tier decomposition address the underlying concern that the graph representation drives the short spacing.

\subsection{Spatial inhomogeneity measured: filament occupancy and scale-dependent over-dispersion}
\label{sec:intermittency}

The local-versus-global contrast invites an obvious question: do cores actually occupy preferential stretches of filament? This can be measured rather than inferred, and doing so turns uneven occurrence from an interpretation into an observable. We assign each associated core an occupancy interval of $\pm0.06$\,pc along its spine (one association half-width), merge overlapping intervals, and define the \emph{occupancy fraction} $f_{\rm occ}$ as the fraction of total skeleton length that is core-bearing. We also count cores in $1$\,pc cells along the spines, from which the Fano factor $F=\sigma^2/\mu$ measures over-dispersion relative to a Poisson process ($F=1$), and compute the Gini coefficient of the same counts.

Table~\ref{tab:intermittency} collects the results, and Fig.~\ref{fig:intermittency} displays them. The occupancy fraction is small in every cloud: between $6$ and $27$\,per cent of the skeleton carries cores, so most filament length is empty at HGBS sensitivity. The Fano factor exceeds unity everywhere ($1.6$, $2.8$, $5.7$, $25.2$ for Aquila, Orion~B, Perseus and Taurus), which is significant over-dispersion on parsec scales, and the Gini coefficients ($0.36$--$0.57$) confirm that cores are distributed very unevenly along the spines. Within the core-bearing components the core-free runs are short (median $0.05$--$0.11$\,pc, $90$th percentile $0.24$--$0.51$\,pc), but they account for only part of the empty length: in Orion~B, for example, $92$\,pc is occupied and a further $155$\,pc lies in core-free runs inside core-bearing components, leaving roughly $90$\,pc of the $337$\,pc total in components containing no catalogued core at all.

\smallskip\noindent\textbf{Is this an artefact of the adopted $\ell=0.06$\,pc interval?} The occupancy fraction necessarily grows with the assumed half-width, so we report it as a function of $\ell$ rather than at one value (Fig.~\ref{fig:scaletests}a) and refer to the number quoted above as the \emph{$0.06$\,pc occupancy fraction} rather than as a duty cycle. The curves separate the clouds cleanly. In Taurus the occupancy rises only from $0.052$ to $0.078$ as $\ell$ increases by a factor of ten, and in Perseus from $0.078$ to $0.202$: widening the interval adds almost nothing because the cores are already tightly grouped. In Orion~B and Aquila the occupancy climbs steadily to $\approx0.6$, so those two clouds are much less strongly segregated. 

We also test the over-dispersion itself against scale, since a scale-independent statement is stronger than one tied to a single cell size. The Fano factor $F(\ell)=\mathrm{Var}[N(\ell)]/\langle N(\ell)\rangle$ (Fig.~\ref{fig:scaletests}b) exceeds the Poisson value of unity across the whole accessible range $\ell=0.25$--$4$\,pc in Taurus ($9.4$--$25$), Perseus ($2.3$--$8.4$) and Orion~B ($1.2$--$4.9$), and rises monotonically with $\ell$ in each. Aquila is the exception: $F=0.92$, $1.06$ and $1.59$ at $0.25$, $0.5$ and $1$\,pc, so it is consistent with Poisson on sub-parsec scales and only mildly over-dispersed at $1$\,pc. We therefore claim scale-persistent over-dispersion for three of the four clouds and not for Aquila. Because $F(\ell)$ needs no occupancy convention at all, it is this statistic, not the occupancy fraction, that carries the inhomogeneity result.

\smallskip\noindent\textbf{Is the over-dispersion more than the filament's own structure implies?} A Fano factor above unity compares the cores with a \emph{homogeneous} Poisson process, and a filament is not homogeneous: column density, line mass and detectability all vary systematically along a skeleton. Cores drawn independently with a probability that simply follows those variations would already give $F>1$. The physically meaningful question is therefore not whether the cores depart from a uniform process but whether they are more clumped than the parent filament structure alone predicts, and we test that directly.

We build an \emph{inhomogeneous} null in which the core intensity is conditioned on the local column density. Binning the skeleton by $N_{\rm H_2}$ we measure the empirical intensity $\lambda(N_{\rm H_2})$, cores per unit length in each bin, from the data themselves; we then scatter synthetic cores along the \emph{actual} skeletons with that intensity and pass them through the same association and statistics pipeline, in $10^4$ realisations per cloud (Appendix~\ref{sec:nullspec}). The result is sobering for a naive reading of $F$. In Orion~B, Aquila and Perseus the observed Fano factors ($2.8$, $1.6$, $5.7$) lie \emph{below} the null medians ($3.8$, $2.3$, $7.3$), with $p=0.98$, $0.88$ and $0.96$: the column-density structure alone accounts for all of the observed over-dispersion, and slightly more than all of it. Only Taurus exceeds its null ($F=25.2$ against a null median of $13.4$ and a $95$th percentile of $17.1$; $p=2\times10^{-4}$, now resolved rather than sitting at the Monte-Carlo floor). We quote empirical $p$-values as $(1+\#\{F_{\rm null}\geq F_{\rm obs}\})/(1+N_{\rm MC})$.

A null whose intensity is fitted to the same cores it is then tested against could be too easily satisfied, so we repeated the test out of sample: for each cloud we built $\lambda(N_{\rm H_2})$ from the \emph{other three} clouds and used it to predict the held-out one. The conclusion is unchanged. Orion~B and Perseus remain below their predicted nulls ($p=1.00$, $0.995$), Aquila is consistent with its null ($F=1.6$ against a predicted median of $1.3$; $p=0.27$), and Taurus alone exceeds it ($F=25.2$ against $8.8$, $95$th percentile $12.7$; $p=0.005$). The single-cloud excess is therefore not an artefact of fitting the intensity to the same data.

One bias remains, and it works against us. If cores raise the column density of their own surroundings, then the pixels carrying the highest $N_{\rm H_2}$, to which the null assigns the highest intensity, are the core sites themselves, and the null then reproduces the observed clumping too easily. We quantified this by masking a $0.06$\,pc radius around every associated core and rebuilding the exposure: the intensity in the highest column-density bins changes by factors of $4$--$10$, so that material is indeed largely core-adjacent and $\lambda(N_{\rm H_2})$ is poorly constrained there by core-free filament. The direction of the bias is towards \emph{under}-detecting excess clustering. Our conclusion that three clouds show no excess should therefore be read as conservative rather than secure, whereas Taurus, which exceeds even this easily-satisfied null, is correspondingly strengthened.

We therefore do not claim that these data demonstrate a distinct physical intermittency of the fragmentation process. What is established is that core occurrence along the adopted skeletons is strongly \emph{inhomogeneous}, and that in three of the four clouds this inhomogeneity is inherited from the filament's own column-density structure rather than added on top of it. Taurus is the one cloud where the cores are clumped beyond that expectation. This is a weaker statement than the one the raw Fano factors invite, and it is the one the evidence supports.

Two length scales therefore coexist. On sub-parsec scales the gap distribution shows a deficit of the smallest separations (Section~\ref{sec:clustering}), which is mild regularity rather than clustering. On parsec scales the counts are strongly over-dispersed. That combination is a two-scale spatial occupancy pattern: cores form in groups with short internal separations, and the groups are separated by long stretches that have produced none. It also explains the factor of about two between the local and global statistics without appealing to small-scale clustering, since the filament-averaged line density is diluted by filament that has not fragmented. The quantity a fragmentation model must reproduce is thus not only a spacing but a spatial occupancy pattern: where along a filament fragmentation happens, as well as how far apart the fragments lie.

\begin{table}
\caption{Spatial-inhomogeneity diagnostics. $L_{\rm tot}$ is the total skeleton length, $L_{\rm core}$ the core-bearing length (union of $\pm0.06$\,pc intervals about each associated core), and $f_{\rm occ}$ is their ratio. $f_{\rm assoc}$ is the fraction of catalogue cores associated with the skeleton, $F$ the Fano factor of counts in $1$\,pc cells and $G$ their Gini coefficient. Both association half-widths are listed. Values come from the independent re-implementation (Section~\ref{sec:linedensity}).}
\label{tab:intermittency}
\footnotesize
\setlength{\tabcolsep}{3pt}
\begin{tabular}{lccccccccc}
\toprule
Region & half & $L_{\rm tot}$ & $L_{\rm core}$ & $f_{\rm occ}$ & $f_{\rm assoc}$ & $\frac{L/N}{W}$ & $\frac{L_{\rm core}/N}{W}$ & $F$ & $G$ \\
       & (pc) & (pc) & (pc) & & & & & & \\
\midrule
Orion B & 0.06 & 337 & 92 & 0.27 & 0.53 & 3.4 & 0.9 & 2.8 & 0.49 \\
        & 0.30 & 337 & 117 & 0.35 & 0.85 & 2.1 & 0.7 & 4.5 & 0.50 \\
Aquila  & 0.06 & 106 & 29 & 0.27 & 0.41 & 3.4 & 0.9 & 1.6 & 0.40 \\
        & 0.30 & 106 & 38 & 0.36 & 0.70 & 2.0 & 0.7 & 2.1 & 0.36 \\
Perseus & 0.06 & 290 & 35 & 0.12 & 0.61 & 5.8 & 0.7 & 5.7 & 0.43 \\
        & 0.30 & 290 & 40 & 0.14 & 0.85 & 4.2 & 0.6 & 10.5 & 0.47 \\
Taurus  & 0.06 & 210 & 13 & 0.06 & 0.86 & 4.5 & 0.3 & 25.2 & 0.57 \\
        & 0.30 & 210 & 14 & 0.07 & 0.97 & 4.0 & 0.3 & 27.3 & 0.57 \\
\bottomrule
\end{tabular}
\end{table}

\begin{figure}
\centering
\includegraphics[width=0.99\columnwidth]{figures/fig_intermittency.pdf}
\caption{Spatial inhomogeneity of core occurrence along the filaments. \emph{(a)} The $0.06$\,pc occupancy fraction (bars, left axis) is $6$--$27$\,per cent, while the Fano factor of counts in $1$\,pc cells (red, right axis, logarithmic) exceeds the Poisson value of unity in every cloud. \emph{(b)} The three normalised spacing statistics: the local adjacent-core spacing is well below the classical band, the global $(L/N)/W$ approaches it, and the core-bearing $(L_{\rm core}/N)/W$ lies below both because the occupied stretches are densely populated. The separation between the blue and green curves measures how unevenly the cores occupy the crest.}
\label{fig:intermittency}
\end{figure}

\begin{figure}
\centering
\includegraphics[width=0.99\columnwidth]{figures/fig_linedensity.pdf}
\caption{The three spacing statistics per cloud, $S_{\rm local}$ (blue), the mean adjacent spacing (orange) and $S_{\rm global}$ (green), from the single re-implementation and all normalised by the fixed $0.10$\,pc width so they are directly comparable. $S_{\rm global}$ is the quantity a one-fragment-per-wavelength theory predicts; against the classical band ($5$--$6$, red; critical $\approx2.6$, dotted). The median along-skeleton ACS spacing (blue, $\approx1.7$--$3.0$) is sub-classical everywhere; the global length per core $(L/N)/W$ (green, $\approx3.4$--$5.8$) is classical in Perseus, near-classical in Taurus, and sub-classical in Orion~B and Aquila; the mean adjacent spacing (orange) lies between but is inflated by arclength branch-jumps. The factor-$\approx2$ gap between the blue and green points reflects the generic mean/median skew plus dilution by core-poor filament (a coverage effect), not small-scale clustering (Section~\ref{sec:linedensity}): short adjacent spacing within core-bearing stretches, diluted by extensive core-free filament.}
\label{fig:linedensity}
\end{figure}

\begin{figure}
\centering
\includegraphics[width=0.99\columnwidth]{figures/fig_paircorr.pdf}
\caption{Along-skeleton pair-correlation function $g(s)$ for the four primary clouds (uniform-random expectation on the matched components divided out; $g=1$ is no correlation). Dotted vertical lines mark the classical scale $5.5\,W_{\rm fil}$ per region. No cloud shows a peak at the classical scale ($g\approx1$ there, slightly below unity in Orion~B and Aquila); the only departures are a mild small-separation excess (Taurus) and weak region-envelope features near $1$\,pc. The classical periodic comb is therefore absent, consistent with a non-periodic, spatially inhomogeneous distribution bearing an observational small-separation deficit (Section~\ref{sec:clustering}).}
\label{fig:paircorr}
\end{figure}

\subsection{Regularity of the core spacing: periodic, random, or clustered?}
\label{sec:clustering}
Classical fragmentation predicts a \emph{quasi-periodic} chain of cores; fibre projection or a turbulent/accretion-driven cascade predicts a broad distribution with no single comb. Ordering associated cores by arclength along each spine and measuring adjacent-core gaps, we find a coefficient of variation $\mathrm{CoV}=\sigma/\mathrm{mean}$ of order unity (Taurus $1.36$, Orion~B $0.83$, Perseus $0.96$, Aquila $0.65$; pooled $1.16$; Fig.~\ref{fig:clustering}).

We calibrated the goodness-of-fit tests by parametric bootstrap, resampling from the fitted distribution to build the exact null distribution of the Kolmogorov--Smirnov statistic and so removing the estimated-parameter bias of the standard test. The resampled gaps are drawn as \emph{independent renewal intervals} from the fitted exponential, not as gaps re-derived from a resampled point pattern. The result is two-sided and unambiguous: the regular ``periodic$+$jitter'' model is rejected ($p_{\rm boot}<10^{-5}$ in every region, at the parametric-bootstrap floor $1/N_{\rm boot}$ with $N_{\rm boot}=10^5$), \emph{and} the pure-exponential (Poisson) model is also rejected ($p_{\rm boot}<10^{-5}$ in Orion~B, Aquila and Taurus, and $1.7\times10^{-3}$ in Perseus). Neither a periodic nor a pure-Poisson description is therefore supported. The exponential fails for different reasons in different clouds, and it is worth being precise about this rather than summarising all four together (Table~\ref{tab:percloud}). Writing $R$ for the ratio of the observed median gap to the median expected for an exponential of the same mean, $R>1$ indicates a deficit of the smallest separations and $R<1$ an excess. Only Orion~B ($R=1.03$) and Aquila ($R=1.14$) show the deficit, and in Aquila alone is it appreciable; Perseus ($R=0.92$) and Taurus ($R=0.64$) show the opposite, an \emph{excess} of small gaps. The beam-scale exclusion we discuss below is therefore a feature of the two clouds with $R>1$, not a universal property of the sample, and Taurus in particular is clumped at small separations as well as over-dispersed on parsec scales. These are four separate per-cloud statements and we do not combine them into a claim about Gould Belt filaments as a class. Comparing five renewal gap models by the corrected Akaike information criterion (AICc; Table~\ref{tab:aic}) makes the two-sided result quantitative and robust to the model set: the pure-Poisson (exponential) model is disfavoured ($\Delta$AICc$=37$--$246$) and the periodic (Gaussian-comb) model strongly rejected ($\Delta$AICc$=81$--$742$) in every region. The specific hard-core (shifted-exponential) model is \emph{not} uniquely preferred once broader alternatives are admitted: a positively-skewed lognormal fits best in three of the four regions (the shifted-exponential is best only in Perseus, where lognormal is $\Delta$AICc$=2.5$ behind). The data therefore favour a broad, skewed distribution truncated at small gaps rather than a sharp physical hard core. The inferred exclusion scale is $0.011$--$0.044$\,pc, of order the beam/core diameter; we therefore treat it as an \emph{observational exclusion scale} rather than a demonstrated physical hard core. The two contributions can be bounded even without a full map-level injection experiment. The $18.2''$ beam corresponds to $0.012$--$0.039$\,pc across our distance range, and \textsc{getsources} does not reliably separate two sources closer than about one beam, so the instrumental floor spans essentially the whole inferred exclusion range $0.011$--$0.044$\,pc: the exclusion is consistent with being entirely instrumental, and we can place no useful lower limit on a physical contribution. Any thermal Jeans-scale suppression would in any case act at $\lambda_{\rm J}\sim0.05$--$0.1$\,pc, above the measured exclusion, so the data do not require a physical hard core at all. A definitive separation needs synthetic source pairs injected into the \textit{Herschel} maps at a range of separations and re-extracted through the HGBS pipeline to measure the completeness as a function of pair separation; that is a map-level experiment we have not performed, and until it is done we treat the exclusion scale as instrumental by default. That the many gaps in a structured cloud may not be independent draws is testable, so we tested it. The lag-one autocorrelation of the gap sequence is weak in every cloud ($|r|<0.14$), but a runs test about the median gap shows significant non-independence in Perseus ($z=-4.2$) and Taurus ($z=-5.4$), with Orion~B ($z=+0.4$) and Aquila ($z=+0.4$) consistent with independence. The deficit of runs means similarly-sized gaps occur together, which is the same core-poor/core-rich alternation documented in Section~\ref{sec:intermittency}. The consequence is that for Perseus and Taurus the effective number of independent gaps is smaller than $N$ and the quoted $p$-values and AICc differences are optimistic; they should not be read at face value for those two clouds. We have therefore implemented a moving-block bootstrap that resamples contiguous runs of five adjacent gaps drawn from within single skeleton components, so that the local dependence structure is preserved in each resample, and recomputed $\Delta$AICc over $400$ resamples per cloud. The effective sample sizes implied by the lag-one autocorrelation are $651$ of $699$ gaps in Orion~B, $225$ of $168$ in Aquila (the negative $r_1=-0.15$ formally raises it), $344$ of $438$ in Perseus and $122$ of $440$ in Taurus, so Taurus loses roughly three quarters of its nominal information. The rejections nevertheless survive with room to spare. The $5$th percentile of the block-bootstrap $\Delta$AICc for the periodic model is $415$ (Orion~B), $45$ (Aquila), $249$ (Perseus) and $587$ (Taurus), and for the pure-Poisson model $190$, $71$, $31$ and $55$; both models are therefore rejected in all four clouds even under the most pessimistic resampling, and the two-sided statement stands without qualification. What the block bootstrap does remove is any claim to discriminate among the surviving broad models: the shifted-exponential is within the lognormal's bootstrap interval in Orion~B, Aquila and Perseus, so ``a broad, positively-skewed gap distribution truncated at small separations'' is as far as the data go. The rejections are large enough ($\Delta$AICc $=37$--$742$) but we treat these as \emph{conditional} goodness-of-fit statistics at fixed skeleton rather than population-level significance tests; the four clouds remain the independent astrophysical units (Section~\ref{sec:limitations}).

\begin{table}
\caption{Gap-model comparison by the corrected Akaike information criterion (AICc; relative to the best model, $\Delta$AICc$=0$), with the number of free parameters $k$. Five renewal models are compared: Poisson (pure exponential), shifted-exponential (hard-core), periodic (Gaussian comb), gamma and lognormal. Both the periodic ($\Delta$AICc$=81$--$742$) and the pure-Poisson ($\Delta$AICc$=37$--$246$) models are strongly disfavoured in every region, so the two-sided ``reject both periodic and Poisson'' result is robust to the model set; but the specific hard-core shifted-exponential is \emph{not} uniquely preferred, a broad, positively-skewed lognormal fits best in three of the four regions, so the observed small-gap deficit is better read as an exclusion-limited (largely instrumental) truncation of a broad distribution than as a demonstrated physical hard core. AICc corrections are negligible here ($n\gg k$). This comparison is \emph{conditional} on the fixed deposited skeleton and is a per-region, not a population-level, statement; with only $N=4$ clouds it does not establish a population-level gap law.}
\label{tab:aic}
\footnotesize
\setlength{\tabcolsep}{3pt}
\begin{tabular}{lccccc}
\toprule
Region & Poisson & shifted-exp & periodic & gamma & lognormal \\
       & ($k{=}1$) & ($k{=}2$) & ($k{=}2$) & ($k{=}2$) & ($k{=}2$) \\
\midrule
Orion B & $245.8$ & $46.4$ & $580.0$ & $87.1$ & $\mathbf{0}$ \\
Aquila  & $82.8$  & $5.8$  & $81.1$  & $11.0$ & $\mathbf{0}$ \\
Perseus & $37.1$  & $\mathbf{0}$ & $367.3$ & $22.4$ & $2.5$ \\
Taurus  & $98.3$  & $76.1$ & $741.8$ & $89.9$ & $\mathbf{0}$ \\
\bottomrule
\end{tabular}
\end{table}

The cores are consistent with a point process bearing a beam-scale exclusion (largely instrumental, and driving the Poisson rejection) plus region-dependent tail clustering (super-Poisson in Taurus, sub-Poisson in Aquila and Orion~B, so ``clustered'' is a population-level description). There is no single characteristic comb, the observational counterpart of the \citet{Clarke2020} result that fragmentation through sub-filaments leaves no single characteristic length.
\begin{figure}
\centering
\includegraphics[width=0.99\columnwidth]{figures/fig_clustering.pdf}
\caption{Adjacent-core gap distributions along the skeletons of the four primary regions (histograms), with exponential/Poisson (solid) and regular/periodic (dashed) reference curves; each panel gives the coefficient of variation CoV, the 1-D Clark--Evans ratio $R$, and the exponential goodness-of-fit $p$-value \emph{calibrated by parametric bootstrap} ($p_{\rm boot}$; the null distribution of the KS statistic is built by resampling from the fitted exponential, removing the estimated-parameter bias of the standard test). The result is two-sided: the regular/periodic model is rejected ($p_{\rm boot}<10^{-5}$, the parametric-bootstrap floor, $1/N_{\rm boot}$ with $N_{\rm boot}=10^5$) in every region, \emph{and} the pure-exponential (Poisson) model is also rejected ($p_{\rm boot}<10^{-5}$ (Perseus $1.7\times10^{-3}$) in all four), so in each of the four clouds the cores are neither periodic nor pure-Poisson; consistent with a point process that is limited at small separations by the extraction resolution but over-dispersed on parsec scales (the deficit of small gaps, $R>1$; largely instrumental).}
\label{fig:clustering}
\end{figure}

Two consequences follow. First, because the gap distribution is broad and positively skewed the mean adjacent spacing exceeds the median by a factor $\approx1.3$--$2.2$ (Orion~B $1.4$, Aquila $1.3$, Perseus $1.6$, Taurus $2.2$; $1.44$ for a pure exponential); but neither the mean nor the median gap is the inverse line density, which is the ordering-independent $N/L$ over the whole skeleton (Section~\ref{sec:linedensity}), and it is $N/L$ that should be set against a one-core-per-wavelength model. We retain the median as the robust descriptor of the observed distribution (Table~\ref{tab:widths}). Second, this fixes the logical status of the comparison: a fragmentation process still sets a characteristic \emph{mean} separation even when the point pattern is not a sharp comb, which clustering, the beam-scale exclusion and projection broaden into the observed near-exponential distribution. The \emph{intra-group} adjacent spacing (median ACS) is sub-classical whereas the filament-averaged line density $N/L$ is near-classical (Taurus, Perseus) and sub-classical only in Orion~B and Aquila; there is thus no globally large excess of cores per unit length, and we read all spacing-vs-wavelength comparisons below as order-of-magnitude consistency of a characteristic scale, not a match of a resonant wavelength.

\begin{table}
\caption{Cloud-by-cloud summary of the point-process diagnostics, to avoid summarising four different behaviours as one. $R$ is the observed median gap divided by the median expected for an exponential of the same mean ($R>1$ deficit of small gaps, $R<1$ excess); $F(1\,{\rm pc})$ is the Fano factor; ``null'' compares $F$ with the inhomogeneous, column-density-conditioned null of Section~\ref{sec:clustering}. The clouds do not behave alike, and only Taurus exceeds its inhomogeneous null. With $N=4$ this is a single case, not a demonstrated population split; and because the null is built partly from column density the cores themselves raise, the three non-detections are conservative.}
\label{tab:percloud}
\footnotesize
\setlength{\tabcolsep}{3.5pt}
\begin{tabular}{lccccc}
\toprule
Region & CoV & $R$ & $F(1\,{\rm pc})$ & vs null & preferred gap model \\
\midrule
Orion B & 0.83 & 1.03 & 2.8 & below ($p=0.99$) & lognormal \\
Aquila  & 0.65 & 1.14 & 1.6 & below ($p=0.95$) & lognormal \\
Perseus & 0.96 & 0.92 & 5.7 & below ($p=0.97$) & shifted-exp. \\
Taurus  & 1.36 & 0.64 & 25.2 & \textbf{above} ($p=0.005$) & lognormal \\
\bottomrule
\end{tabular}
\end{table}

\subsection{Combined Systematic Error Budget}
\label{sec:error_budget}

We separate three kinds of term (Table~\ref{tab:errorbudget}): \emph{statistical/sensitivity} terms (the cloud-to-cloud and choice-dependent scatter that would change under resampling), \emph{systematic} terms (deterministic offsets from a fixed methodological choice, e.g.\ the pipeline bias correction), and \emph{modelling} terms (arising from the simulation set-up, e.g.\ the line-mass dependence). The uncorrelated statistical/sensitivity terms are combined in quadrature,
\begin{equation}
\sigma_{\rm sys}^2=\sum_i\sigma_i^2,
\label{eq:quadrature}
\end{equation}
and the multiplicative width- and skeleton-convention factors are then applied \emph{afterwards} as separate \emph{analysis-choice envelopes}, not added into this quadrature sum. The adopted region-Plummer ACS result $(s_{\rm ACS,med}/W)_0\approx1.4$ (raw fixed-width analogue $\approx2.1$) then carries
\begin{itemize}
\setlength{\itemsep}{0pt}
\item \emph{in quadrature}: residual pipeline bias ($\sim$10\,per cent), protostellar migration ($\sim$7.5\,per cent), between-region scatter ($\sim$4\,per cent), and a $\lesssim5$\,per cent manual-spur proxy (Section~\ref{sec:methodology});
\item \emph{as multiplicative analysis-choice envelopes applied afterwards}: the width convention ($\sim$25\,per cent; \citealt{Panopoulou2022}, counted once) and the skeleton algorithm (a factor $\sim$2).
\end{itemize} The skeleton algorithm is the single largest observational term (Fig.~\ref{fig:skeleton_compare}; discussed below). Because the dominant term is a deterministic offset we treat the aggregate as a min--max envelope, $s_{\rm ACS,med}/W\approx0.8$--$1.9$ (width-convention min--max, DisPerSE; Table~\ref{tab:headline}), rather than a Gaussian $\sigma$ (the skeleton factor of two is carried separately). The adopted headline is $(s_{\rm ACS,med}/W)_0\approx1.4$; the fixed-width $s_{\rm ACS,med}/W\approx1.9$ is a comparator anchored to that convention, not a competing headline.

\emph{Projection.} The observed separations are plane-of-sky projections of intrinsically 3-D spacings. For an isotropic distribution of filament orientations the projection factor $s_{\rm obs}/s_{\rm true}=\sin i$ has median $0.87$ and mean $0.79$ (Monte-Carlo, $2\times10^6$ orientations; 16--84\% range $0.54$--$0.99$), so deprojection would \emph{raise} the intrinsic spacing by $\approx15$--$27\%$, moving $s_{\rm ACS,med}/W$ towards the classical value (e.g.\ the adopted $1.4$ becomes $1.6$--$1.8$) and reducing, but not removing, the sub-classical deficit. We do not apply this correction, since it depends on the unknown orientation distribution, but carry it as a one-sided comparison systematic that acts \emph{against} the sub-classical result rather than for it.

The legacy pairwise-median value ($\approx2.8$ in $W$ units) is an $L/3$ scale, not a local core spacing (Section~\ref{sec:statistics}); its bootstrap scatter is $\pm4.2\%$, and we quote no $\sigma$-level tension with theory for it.

\emph{Skeleton-algorithm dependence.} The DisPerSE and FilFinder values ($s_{\rm ACS,med}/W\approx1.9$ and $\approx0.4$--$0.8$; Fig.~\ref{fig:skeleton_compare}) bracket a genuine ambiguity in what constitutes ``the filament''; both lie far below $\lambda_{\rm m}$, so the factor of two moves the magnitude but not the sign. A partial matched-threshold FilFinder extraction shows the inferred spacing is itself strongly sensitive to the mask threshold, so the factor is real and, if anything, a lower bound; an agreed, algorithm-independent segment definition is the important future work. The most promising route is not a better image-plane algorithm but kinematic information: velocity-coherent structures identified in molecular-line cubes (for example N$_2$H$^+$ or NH$_3$) define filaments dynamically rather than morphologically, and would break the DisPerSE/FilFinder degeneracy by testing which extraction corresponds to a coherent physical object. The skeletoniser also selects the interpretation: in matched Plummer units the simulated single-filament upper bound ($\lesssim1.0$--$1.3$) overlaps the region-Plummer values within the methodological envelope, so a denser FilFinder extraction would need the extra compression of hierarchical-fibre projection (Section~\ref{sec:fibre_projection}) while DisPerSE would not.

\noindent\emph{Note on uncertainties: with only four clouds, all $\sigma$ and $\pm$ values in this paper (including in Table~\ref{tab:errorbudget} and Fig.~\ref{fig:spacing}) are descriptive half-widths of an $N=4$ range or min--max envelopes, not formal confidence intervals.}

\begin{table}
\caption{Summary of the terms on the $\lambda/W$ comparison, grouped as \emph{observational} (on $s_{\rm ACS,med}$), \emph{simulation} (on $\lambda_{\rm turn}$) and \emph{comparison} (terms that enter only when the two are matched). Statistical/sensitivity terms are combined in quadrature (Eq.~\ref{eq:quadrature}); the width- and skeleton-convention factors are applied afterwards as a separate multiplicative \emph{analysis-choice envelope}, not as Gaussian uncertainties. All $\pm$ values and the $\sigma\approx0.4$ cloud-to-cloud scatter are descriptive $N=4$ envelopes (half-widths or min--max ranges), \emph{not} Gaussian confidence intervals (Section~\ref{sec:error_budget}).}
\label{tab:errorbudget}
\footnotesize
\setlength{\tabcolsep}{3pt}
\begin{tabular}{@{}p{3.1cm}p{1.9cm}p{2.3cm}@{}}
\toprule
Source & Type & Magnitude \\
\midrule
\multicolumn{3}{@{}l}{\emph{Observational (on $s_{\rm ACS,med}$)}}\\
Bootstrap / cloud-to-cloud & sensitivity & $\sigma\approx0.4$ ($\sim$20\%) \\
ACS pipeline bias & systematic & $\sim$10\% (corrected) \\
Protostellar migration & sensitivity & $\sim$7.5\% \\
Projection (deprojection) & systematic (one-sided, \emph{raises} $s$) & $+15$--$27\%$ \\
LOS blending (Aquila) & systematic (one-sided, \emph{shortens} $s$) & $\lesssim50\%$ worst case \\
Skeleton-coverage selection & systematic (one-sided, \emph{shortens} $s$) & $\lesssim16\%$ (bounded) \\
Width convention, fixed vs region & analysis-choice env. & $\sim$25\% \\
Skeleton algorithm, DisPerSE/FilFinder & analysis-choice env. & factor $\sim$2 \\
Manual skeleton editing & systematic (proxy) & $\lesssim$5\% \\
\midrule
\multicolumn{3}{@{}l}{\emph{Simulation (on $\lambda_{\rm turn}$)}}\\
Growth-rate resolution/epoch & modelling & $\lesssim$5\% (converged) \\
Turbulence sensitivity & sensitivity & $<$5\% (single realisation) \\
Line-mass ($f$) dependence & modelling & $\sim$25\% \\
\midrule
\multicolumn{3}{@{}l}{\emph{Comparison (matching the two)}}\\
Gaussian vs Plummer FWHM & systematic & factor $\sim$1.5 \\
$T_1$ width normalisation ($\eta_W$) & systematic & $\sim$30\% \\
\bottomrule
\end{tabular}
\end{table}

The comparison-only terms, the $T_1$ width normalisation ($\eta_W=0.59$--$0.78$; Section~\ref{sec:width_correction}) and the Gaussian-vs-Plummer FWHM factor, enter only when the simulated $\lambda/W_{\rm core}$ is converted to observable units; they do not affect the ACS observational measurement, which is already in $W_{\rm fil}$ units. The $T_1=0.61$ (band $0.59$--$0.78$; $\sim30\%$) normalisation is propagated as a multiplicative systematic on the \emph{simulated} $s_{\rm ACS,med}/W$ alone (via the $W_{\rm core}/[T_1 W_{\rm form}]$ conversion of Eq.~\ref{eq:t1_apply}): the $\pm0.10$ maps to a $\sim30\%$ band on the simulated value when it is set against the region-Plummer baseline, and is carried in the comparison group here rather than on the observed spacing.

The empirical \emph{per-region range} of the bias-corrected fixed-width values is $s_{\rm ACS,med}/W=1.0$--$2.1$ (the real cloud-to-cloud spread), and the \emph{methodological envelope} $s_{\rm ACS,med}/W\approx1.1$--$2.4$ (per-region sweep incl.\ FilFinder; $\sigma_{\rm sys}\approx28\%$) is obtained by varying the width, skeleton and association-radius choices; the descriptor ``$\sigma\approx0.4$'' is only the half-width of the per-region range.

\subsection{ACS pipeline reproducibility, population splits and hierarchical statistics}
\label{sec:nnpipeline}

The along-skeleton ACS pipeline thresholds the deposited DisPerSE skeleton, bridges fragmentation by morphological closing and re-skeletonisation, associates each core to the nearest spine within $0.06$\,pc, orders cores along each component by graph arclength (double-BFS diameter), and takes adjacent-core separations, rejecting cores not on a primary segment. It reproduces the deposited catalogue counts and raw medians exactly (Table~\ref{tab:nnpipeline}). The estimator was also re-implemented from scratch as an internal check against coding error, reproducing the same counts and medians (see Data Availability).

\begin{table}
\caption{ACS pipeline counts: skeleton components, catalogue cores, cores associated to a skeleton, and adjacent ACS pairs, per primary region.}
\label{tab:nnpipeline}
\footnotesize
\setlength{\tabcolsep}{4pt}
\begin{tabular}{lcccc}
\toprule
Region & comps & catalogue & associated & ACS pairs \\
\midrule
Orion B & 67 & 1844 & 680 & 614 \\
Aquila  & 14 & 749  & 182 & 168 \\
Perseus & 54 & 816  & 619 & 563 \\
Taurus  & 32 & 536  & 479 & 461 \\
\bottomrule
\end{tabular}
\end{table}

The fraction of catalogue cores associating to a primary skeleton segment varies across regions, Orion~B $37\%$, Aquila $24\%$, Perseus $76\%$, Taurus $89\%$, so in Orion~B and Aquila the ACS spacing is measured on a minority of cores, the distributed clouds having many cores off any primary crest. Any bias runs \emph{toward} the sub-classical result (rejected cores lie on fainter segments with \emph{longer} spacings; Section~\ref{sec:linedensity}), and the result does not depend on the low-coverage regions: restricting to the high-coverage clouds (Perseus, Taurus) gives $s_{\rm ACS,med}/W=1.8$ and $1.1$, median $1.4$, identical to the four-region median. It is also stable to the association radius: sweeping the half-width over $0.04$--$0.08$\,pc changes the per-region ACS median by only $4$--$7\%$ (Taurus $6.6\%$, Orion~B $7.3\%$, Perseus $7.4\%$, Aquila $4.2\%$), so the effect is not load-bearing. The complementary \emph{relaxed}-radius coverage test reinforces this: widening the association half-width from $0.06$ to $0.30$\,pc raises the coverage to $70$--$97\%$ yet lengthens the median adjacent spacing by only $+3$ to $+16\%$ (Section~\ref{sec:linedensity}), so the omitted cores bias the result short by a bounded, one-sided $\lesssim16\%$ and the spacing stays sub-classical.

Splitting the sample by the published classification flags, the prestellar-only median $s_{\rm ACS,med}/W=2.10$ ($717$ pairs) is indistinguishable from the all-cores $2.08$ ($1806$ pairs; both $\approx1.9$ after bias correction). Prestellar cores dominate the on-skeleton budget in every region and set the result; unbound starless cores give a larger, noisier $s_{\rm ACS,med}/W\approx3.3$ and protostellar cores are too few ($6$--$80$ pairs) to constrain it. The sub-classical result is therefore not an artefact of population mixing.

\smallskip\noindent\textbf{Hierarchical statistics and angular-resolution robustness.}

Cores within a filament are not independent, so we assessed the cloud$\to$filament$\to$core nesting with a three-level nonparametric bootstrap (2000 iterations, carrying spacings that share a central core together; the second level samples the $67,14,54,32$ filament components per cloud, Table~\ref{tab:nnpipeline}). On a 2-D projected ACS metric (which underestimates the arc length, $s_{\rm ACS,med}/W\approx1.1$--$1.5$) the bootstrap median is stable, the dominant uncertainty being the cloud-to-cloud scatter $\sigma\approx0.4$ (range $\sim0.8$ in Taurus to $\sim1.7$ in Aquila); core-count and equal-cloud weighting agree to $\sim0.1$. We therefore quote this scatter, not a per-pair error, as the relevant uncertainty (Section~\ref{sec:error_budget}).

Angular-resolution blending is negligible: merging cores within one beam shifts the characteristic $s_{\rm ACS,med}/W$ by $\le0.004$, so distance-dependent blending does not bias the adopted value (Table~\ref{tab:validation_summary}).

\begin{table}
\caption{Summary of validation tests performed in this paper.}
\label{tab:validation_summary}
\footnotesize
\setlength{\tabcolsep}{3pt}
\begin{tabular}{p{2.3cm}p{2.9cm}p{2.0cm}}
\toprule
Test & What it checks & Result \\
\midrule
Injection--recovery (ACS) & ACS bias on real skeletons & $3\%$; pipeline $15$--$20\%$ \\
Injection--recovery (pairwise) & $L/3$ convergence & tracks $L/3$, not $\lambda$ \\
Resolution convergence & $t_{\rm frag}(128^3)$ vs $256^3$ & ratio $0.915\pm0.012$ \\
Truelove criterion & cells per $\lambda_{\rm J}$ at measurement & long.\ $5$--$7$; perp.\ $1$--$2$ \\
IC sensitivity & Gaussian vs uniform profile & identical \\
EOS sensitivity & $\gamma=0.8$--$5/3$ & $\gamma<1$ accelerates; $5/3$ suppresses \\
Region resampling (descriptive) & spread over 4 clouds & $0.261$--$0.298$\,pc \\
Leave-one-out & single-region dominance & $<6\%$ shift \\
Population splits & prestellar vs all cores & $s_{\rm ACS,med}/W\!=\!2.10$ vs $2.08$ \\
Hierarchical bootstrap & cloud$\to$filament$\to$core & $\sigma_{\rm cloud}\!\approx\!0.4$ \\
Mock-blending & distance-dependent beam & $\Delta s_{\rm ACS,med}/W\!\le\!0.004$ \\
Region-specific Plummer widths & beam-deconvolved $W_{\rm fil}$ & $0.11$--$0.16$\,pc; $s_{\rm ACS,med}/W\!=\!1.1$--$1.8$ \\
\bottomrule
\end{tabular}
\end{table}

\subsection{Comparison with Literature}

Our revised-distance ACS measurements (Table~\ref{tab:literature}) differ from the original HGBS values by the region-dependent distance revisions; the core-count-weighted primary-region mean rises only $\sim10\%$, dominated by Aquila.

\begin{table*}
\caption{Comparison with published HGBS spacing measurements. Perseus and Taurus are the two regions in which our own gap sequence shows significant non-independence (runs test; Section~\ref{sec:clustering}) and in which the post-masking eligible-core samples collapse (Section~\ref{sec:eligible}); their entries should be compared with the published values with those caveats in mind.}
\label{tab:literature}
\setlength{\tabcolsep}{4pt}
\footnotesize
\setlength{\tabcolsep}{2.5pt}
\begin{tabular}{lcccccc}
\toprule
Region & ACS (pc) & Orig.\ HGBS & pairwise (pc) & Lit.\ (pc) & Lit.\ statistic & Reference \\
\midrule
\textbf{Primary sample} &  &  &  &  &  &  \\
Orion B & \textbf{0.207} & 0.212 & 0.313 & 0.20--0.30 & local & \citet{Konyves2020} \\
Aquila  & \textbf{0.210} & 0.206 & 0.346 & 0.22--0.26 & local & \citet{Konyves2015} \\
Perseus & \textbf{0.263} & 0.199 & 0.248 & 0.22 & local & \citet{Pezzuto2021} \\
Taurus  & \textbf{0.124} & 0.206 & 0.198 & 0.19 & along-fibre & \citet{Hacar2013} \\
\midrule
\textbf{Limited sample} &  &  &  &  &  &  \\
Ophiuchus &, & 0.197 & 0.206 & 0.18 & local & \citet{Konyves2020} \\
Serpens  &, & 0.188 & 0.331 & 0.17--0.19 & local & \citet{Konyves2020} \\
TMC1     &, & 0.203 & 0.195 & $\sim$0.20 & along-fibre & \citet{Hacar2013} \\
CRA      &, & 0.212 & 0.248 & $\sim$0.21 & local & \citet{Konyves2020} \\
\bottomrule
\end{tabular}

\vspace{0.1in}
\footnotesize
\emph{Notes} $^{a}$ The ``Lit.\ statistic'' column records what kind of statistic each cited value represents. All are \emph{local} separation measures, either nearest-neighbour or along-fibre adjacent separations, rather than all-pairs medians, and that is precisely why they agree with our ACS column rather than with our pairwise column; the exact definitions differ in detail between papers. The region-extent bias identified here therefore applies wherever an all-pairs or pooled separation statistic is adopted, including the comparability statistic carried in this paper and in cross-region compilations, and not to these particular published values. $^{b}$ The primary statistic is the \emph{adjacent-core} (ACS) spacing (bold); the ``pairwise'' column is a legacy comparator only, inflated towards $L/3$ (Section~\ref{sec:statistics}) and not a physical spacing. ACS spacings for the limited sample are omitted. Perseus is the one region where the ACS ($0.263$) slightly exceeds the pairwise median, because arclength ordering on its heavily bridged skeleton inflates path distances; we verified this is not a bridging artefact, an independent minimum-spanning-tree ordering of the same cores gives a consistent value, insensitive to the morphological bridging, and both remain sub-classical. $^{c}$ ``Original HGBS'' values use pre-Gaia distances; literature ranges reflect different methods and distance assumptions. $^{d}$ For Orion~B the pre-Gaia distance was already $\sim$400\,pc \citep{Konyves2020}, so the pairwise $0.212\to0.313$\,pc difference is a statistic-definition effect, not a distance revision.
\end{table*}

We decomposed this shift on the real Orion~B data by toggling each methodological choice at fixed distance (Table~\ref{tab:decomp}). The \emph{statistic definition} dominates: switching from ACS to the pairwise median inflates the spacing by $\approx1.5$ ($0.207\to0.313$\,pc), accounting for essentially the entire $0.212\to0.313$\,pc difference in Table~\ref{tab:literature}; the remaining choices (skeleton re-extraction, association half-width, 2-D$\to$arclength ordering, and the $-3.5\%$ distance revision) are each $\le15\%$ and largely cancel (Table~\ref{tab:decomp}). On our \emph{primary} ACS statistic Orion~B is $0.207$\,pc, within $3\%$ of \citet{Konyves2020}, so the apparent shift is a property of the pairwise-median comparability statistic, not the physical spacing; the ACS re-analysis \emph{shortens} Taurus and \emph{lengthens} Perseus, if anything strengthening the conclusion.

\begin{table}
\caption{Decomposition of the apparent Orion~B spacing shift ($0.212\to0.313$\,pc in Table~\ref{tab:literature}) into its methodological contributions, toggled one at a time at fixed distance. The statistic definition (ACS$\to$pairwise median) dominates; the remaining choices are small and largely cancel.}
\label{tab:decomp}
\footnotesize
\setlength{\tabcolsep}{4pt}
\begin{tabular}{@{}p{4.3cm}l@{}}
\toprule
Contribution & Effect on spacing \\
\midrule
Statistic definition (ACS $\to$ pairwise median) & $\times1.51$ ($0.207\!\to\!0.313$\,pc) \\
Skeleton re-extraction & $\le15\%$ (partly cancels) \\
Association half-width & $\le15\%$ (partly cancels) \\
2-D $\to$ arclength ordering & $\le15\%$ (partly cancels) \\
Distance revision & $-3.5\%$ \\
\midrule
Net (all toggled) & $\times1.48$ ($0.212\!\to\!0.313$\,pc) \\
\bottomrule
\end{tabular}
\end{table}

\subsection{What constitutes the theoretical filament?}
\label{sec:whatfilament}
Before comparing with theory we note that the factor-of-two DisPerSE-versus-FilFinder difference (Fig.~\ref{fig:skeleton_compare}) is not measurement noise but reflects that the two algorithms \emph{define different objects} (persistent crests versus the medial axis of a masked skeleton), so there is no single, algorithm-independent ``filament'' whose spacing theory should reproduce. Our robustness claim applies to the \emph{sign} of the deficit, every convention gives a sub-classical $s_{\rm ACS,med}/W$, while its \emph{magnitude} cannot be assigned a unique value: the same data yield $s_{\rm ACS,med}/W\approx1.0$ (region-Plummer FilFinder), $\approx1.4$ (region-Plummer DisPerSE) or $\approx1.9$ (fixed-width DisPerSE). Because segmentation alone moves the answer by as much as the differences between the physical models compared in Section~\ref{sec:validation_endtoend}, the theory should not be read as discriminating between these values. Figure~\ref{fig:envelope} makes the resulting envelope explicit.

\begin{figure}
\centering
\includegraphics[width=0.99\columnwidth]{figures/fig_envelope.pdf}
\caption{The full methodological envelope of $\lambda/W$, organised by \emph{which quantity is being measured} rather than by convention alone. The four groups are the conditional adjacent-core spacing $S_{\rm local}$ (blue), the unconditional length per core $S_{\rm global}=(L/N)/W$ (green), the same restricted to fragmentation-eligible filament $S_{\rm elig}$ (purple; the core-masked row is fiducial, the raw-map row shows the endogeneity-affected comparison), and the simulated upper bound $S_{\rm sim}^{\rm upper}=\lambda_{\rm turn}/W$ together with the hierarchical-fibre toy model (red). Each bar spans the range across clouds and conventions; the grey band is the classical $\lambda_{\rm m}\approx5$--$6$ (\citealt{Inutsuka1992,Fischera2012}) and the dotted line the critical $\lambda_{\rm cr}\approx2.6$. Grouping this way makes the central point of the paper visible directly: $S_{\rm local}$ lies well below the classical band under every convention, $S_{\rm global}$ reaches it, $S_{\rm elig}$ falls between the two, and $S_{\rm sim}^{\rm upper}$ is an \emph{upper bound} that should not be read as a measured wavelength. Comparisons between groups are comparisons between different quantities and are only meaningful once that is stated (Section~\ref{sec:whatfilament}). The horizontal axis is symmetric-log below $\lambda/W=1$. The $S_{\rm sim}^{\rm upper}$ bar is drawn with a left-pointing arrow because it is an upper limit, not a measured interval; the open marker on the $S_{\rm local}$ group shows where the adopted value moves if the isotropic deprojection factor of Table~\ref{tab:headline} is applied.}
\label{fig:envelope}
\end{figure}

\section{THEORETICAL FRAMEWORK}
\label{sec:theory}

\subsection{Classical Fragmentation Theory}

\begin{table}
\caption{The benchmark this paper is measured against, stated once and in full, because the ``four filament widths'' often quoted in the literature is \emph{four times the flat diameter} of the Ostriker equilibrium and not four times the observed FWHM. Ratios are exact for the isothermal equilibrium ($p=4$) cylinder; the last is the one that is \emph{not} invariant when the observed profile has $p\simeq2$ (Table~\ref{tab:benchmark}).}
\label{tab:benchmarkbox}
\footnotesize
\begin{tabular}{ll}
\toprule
Scale & Value \\
\midrule
Scale height & $H=c_s/\sqrt{4\pi G\rho_c}$ \\
Flat radius & $R_{\rm flat}=\sqrt{8}\,H$ \\
Flat diameter & $2R_{\rm flat}=2\sqrt{8}\,H\simeq5.66\,H$ \\
Fastest-growing mode & $\lambda_{\rm m}=22H=4\times(2R_{\rm flat})$ \\
Critical wavelength & $\lambda_{\rm cr}\simeq2.6\times{\rm FWHM}$ \\
\midrule
\textbf{Benchmark used here} & $\boldsymbol{\lambda_{\rm m}\simeq5}$--$\boldsymbol{6\times{\rm FWHM}}$ \\
\bottomrule
\end{tabular}
\end{table}


For an isothermal, self-gravitating cylinder of density $\rho_0$ and sound speed $c_s$, the critical line mass for gravitational instability is \citep{Ostriker1964}:
\begin{equation}
\mu_{\rm crit} = \frac{2c_s^2}{G}
\label{eq:mucrit}
\end{equation}
The most unstable fragmentation wavelength is, as an \emph{exact} result for the isothermal equilibrium cylinder \citep{Inutsuka1992},
\begin{equation}
\lambda_{\rm IM92} = 22\,H = 4\times(2R_{\rm flat}),
\label{eq:im92}
\end{equation}
where $H = c_s\,(4\pi G\rho_c)^{-1/2}$ is the thermal scale height, $R_{\rm flat}=\sqrt8\,H$ the flattening radius, and $W_{\rm fil}\approx0.1$\,pc the characteristic filament FWHM. Solving the dispersion relation exactly (Appendix~\ref{sec:val_appendix}; \citealt{Fischera2012}) gives a fastest-growing wavelength $\lambda_{\rm m}\approx5\times W_{\rm fil}$ near-critical, rising to $\approx6\times W_{\rm fil}$ at low line mass, and a critical wavelength $\lambda_{\rm cr}\approx2.6\times W_{\rm fil}$. The often-quoted ``$\lambda\approx4\times$ the width'' is $4\times$ the flat \emph{diameter} $2R_{\rm flat}$ \citep{Inutsuka1992}, which equals $\approx5$--$6\times$ the \emph{FWHM}; taking it as $4\times$ the FWHM understates the classical scale. We therefore adopt $\lambda_{\rm m}\approx5$--$6\,W_{\rm fil}$ (Fischera \& Martin 2012) as the benchmark. Because we compare in FWHM units and forward-model the simulations to a beam-convolved Plummer FWHM (Section~\ref{sec:width_correction}), the sub-classical deficit does not depend on the equilibrium-vs-observed profile-index difference ($p=4$ vs $p\simeq2$), and we compare as the density-independent ratio $\lambda/W_{\rm fil}$. The deficit is not new in sign: \citet{Fischera2012} noted observed cores sit ``about the FWHM'' apart, and \citet{Zhang2023cal} find a spacing $\approx1.3\times$ the width in California.

\begin{table}
\caption{Where the $5$--$6\times$FWHM benchmark comes from, for the isothermal equilibrium (Ostriker, $p=4$) cylinder. The ratios in the last column are exact for that equilibrium; the final row is the one that is \emph{not} invariant when the observed profile has $p\simeq2$.}
\label{tab:benchmark}
\footnotesize
\setlength{\tabcolsep}{4pt}
\begin{tabular}{llc}
\toprule
Quantity & Definition & Value \\
\midrule
$H$ & $c_s(4\pi G\rho_c)^{-1/2}$, scale height & $1$ \\
$R_{\rm flat}$ & flattening radius, $\sqrt{8}\,H$ & $2.83\,H$ \\
$2R_{\rm flat}$ & flat diameter & $5.66\,H$ \\
$W_{\rm FWHM}$ & FWHM of the $p=4$ profile & $\approx3.9\,H$ \\
$\lambda_{\rm cr}$ & cutoff wavelength & $\approx10\,H\;(\approx2.6\,W)$ \\
$\lambda_{\rm m}$ & fastest-growing wavelength, $22H$ & $22\,H\;(\approx5.6\,W)$ \\
\midrule
$\lambda_{\rm m}/(2R_{\rm flat})$ & the ``$4\times$ diameter'' form & $3.9$ \\
$\lambda_{\rm m}/W_{\rm FWHM}$ & the benchmark used here & $5$--$6$ \\
\bottomrule
\end{tabular}
\end{table}

Two things follow. The ratio $\lambda_{\rm m}/H=22$ and the conversion $\lambda_{\rm m}\approx4\times(2R_{\rm flat})$ are properties of the $p=4$ equilibrium and are exact within it. What is \emph{not} invariant is $\lambda_{\rm m}/W_{\rm FWHM}$, because the FWHM of a $p\simeq2$ Plummer profile is a different fraction of $H$ from that of a $p=4$ profile, and beam convolution broadens it further. The observational benchmark is therefore convention-dependent at the tens-of-per-cent level even before any measurement error, which is one reason we treat the simulated comparison through forward modelling (Section~\ref{sec:width_correction}) rather than through this ratio.

With $f=\mu_{\rm line}/\mu_{\rm crit}$ and $t_{\rm frag}\propto(f-1)^{-1/2}$, marginally supercritical filaments fragment slowly. At magnetically supercritical field strength ($\beta\gtrsim0.2$) all configurations fragment for $\mathcal{M}\geq1$ at every $f$ tested once integrated adequately: above the critical line mass there is no \emph{thermal} stability boundary. This is distinct from \emph{magnetic} stabilisation, where a strong longitudinal field ($\beta\lesssim0.15$) suppresses collapse even for $f>1$ (the ``strong-field, magnetically supported'' regime; Appendix~\ref{sec:supplementary_campaigns_app}); the thermal ($f$) and magnetic (mass-to-flux) criticalities are physically distinct.

We adopt a single reference $W_{\rm fil}=0.10$\,pc for comparability with prior HGBS work, though its universality is debated \citep{Panopoulou2017,Panopoulou2022}; the directly-measured region-Plummer widths ($0.11$--$0.16$\,pc, $\lambda/W\approx1.4$; Section~\ref{sec:widths}) show the sub-classical conclusion is robust to the convention, carried as the $\sim25\%$ width systematic (Section~\ref{sec:error_budget}).

\subsection{Width normalisation: $W_{\rm core}$ to $W_{\rm fil}$}
\label{sec:width_correction}

\begin{figure}
\centering
\includegraphics[width=0.99\columnwidth]{figures/fig_schematic.pdf}
\caption{Schematic of the length-scale definitions used throughout (Table~\ref{tab:symbols}). \emph{(a)} Transverse column-density profile, marking the initial half-width $W_{\rm core}=\sigma_0$, the evolved spine FWHM $W_{\rm form}$, and the observed beam-convolved Plummer FWHM $W_{\rm fil}=W_{\rm form}/T_1$ ($T_1\simeq0.61$; Eq.~\ref{eq:t1_apply}). \emph{(b)} Axial view contrasting the classical widely-spaced comb ($\lambda_{\rm m}\approx5$--$6\,W$) with the observed unevenly-distributed cores: a short intra-group adjacent spacing $s_{\rm ACS,med}$ within core-rich stretches separated by core-poor filament, whose filament-averaged line density $L/N$ is near-classical, while the simulated fragmentation scale is an upper bound $\lambda_{\rm turn}\lesssim1\,\lambda_{\rm J}$.}
\label{fig:schematic}
\end{figure}
Three widths enter the simulation--observation comparison and must be kept separate (Fig.~\ref{fig:schematic}; Table~\ref{tab:widthdefs}): the initial Gaussian half-width $W_{\rm core}=0.3\,\lambda_{\rm J}$; the formation FWHM $W_{\rm form}=0.683\,\lambda_{\rm J}$ of the evolved projected profile ($\simeq2.28\,W_{\rm core}$, slightly below the initial $2.355\,W_{\rm core}$ owing to contraction; resolution- and seed-converged to $<1\%$); and the observed HGBS width $W_{\rm fil}$, the beam-convolved Plummer FWHM \citep{Arzoumanian2019}, which we obtain by forward-modelling the simulated filaments through the full synthetic-HGBS pipeline (projection, beam convolution, Plummer fit). This gives $\eta_W=W_{\rm form}/W_{\rm fil}=0.59$--$0.78$ and
\begin{equation}
T_1 \equiv \frac{W_{\rm form}}{W_{\rm fil}} = 0.61 \quad (\text{forward-model band } 0.59\text{--}0.78),
\label{eq:t1_correction}
\end{equation}
so that $W_{\rm fil}=W_{\rm form}/T_1\approx1.12\,\lambda_{\rm J}$. We quote the forward-model best-fit $T_1=0.61$; because it lies near the lower edge of the band, a mid-band $T_1\approx0.68$ would give a slightly \emph{larger} (less sub-classical) $\lambda/W_{\rm fil}$, so this normalisation is deliberately not the headline; the primary comparison uses the width-convention-independent growth-rate upper bound (Section~\ref{sec:validation_endtoend}), and this $T_1$ conversion is retained only as a cross-check. The observable is $\lambda/W_{\rm fil}$; because the fragmentation spacing $\lambda$ is a physical length independent of the width definition, the conversion from the simulation-internal ratio $\lambda/W_{\rm core}$ (normalised by the initial $\sigma=0.3\,\lambda_{\rm J}$) is
\begin{equation}
\Big(\tfrac{\lambda}{W}\Big)_{\rm fil} = \frac{W_{\rm core}}{W_{\rm fil}}\Big(\tfrac{\lambda}{W}\Big)_{\rm core} \approx 0.27\,\Big(\tfrac{\lambda}{W}\Big)_{\rm core},
\label{eq:t1_apply}
\end{equation}
i.e.\ $\lambda/W_{\rm fil}=\lambda\,[\lambda_{\rm J}]/1.12$; the prefactor is $0.27=W_{\rm core}/W_{\rm fil}=0.30/1.12$ (equivalently $T_1$ times the $\sigma\!\to\!$FWHM factor, $0.61\times0.44=0.27$, using $T_1\,W_{\rm core}/W_{\rm form}$), and using $T_1$ alone would overstate $\lambda/W_{\rm fil}$ by $\approx2.3\times$. This prefactor is \emph{not} a fixed constant: it carries the full $\approx30\%$ $T_1$ band ($T_1=0.59$--$0.78$), so $0.27$ is the best-fit value of a $\approx0.26$--$0.34$ range. Because the observable is sensitive to how $W_{\rm fil}$ is defined, the primary comparison uses the width-convention-independent growth-rate wavelength (an upper limit; Section~\ref{sec:validation_endtoend}), and this normalisation is retained only as a cross-check.

\begin{table}
\caption{The three filament widths used in the simulation--observation comparison (all in Jeans lengths, $\lambda_{\rm J}=1$).}
\label{tab:widthdefs}
\footnotesize
\setlength{\tabcolsep}{3pt}
\begin{tabular}{p{1.5cm}p{1.4cm}p{4.0cm}}
\toprule
Symbol & Value ($\lambda_{\rm J}$) & Definition \\
\midrule
$W_{\rm core}$ & $0.30$ & initial Gaussian half-width $\sigma$ (input) \\
$W_{\rm form}$ & $0.683$ & Gaussian FWHM of evolved projected profile, supercritical ensemble ($\simeq2.28\,W_{\rm core}$) \\
$W_{\rm FWHM}$ & $0.58$ & Gaussian FWHM at the growth-rate epoch, near-critical runs (narrower, already contracting); normalises Table~\ref{tab:wave} \\
$W_{\rm fil}$ & $\approx1.12$ & beam-convolved Plummer FWHM ($=W_{\rm form}/T_1$); the observable normalisation \\
\bottomrule
\end{tabular}
\end{table}

For the ambient-confined supercritical ensemble ($\lambda/W_{\rm core}=3.27$) this gives $\lambda/W_{\rm fil}\approx0.9$.

\subsection{Magnetic Tension Mechanism}
\label{sec:magnetic}

For a filament threaded by a magnetic field at an angle $\theta$ to the axis (Alfv\'en speed $v_A$), a schematic oblique dispersion relation, valid only for a static background field, is \citep{Nakamura1993}:
\begin{equation}
\omega^2 = c_s^2 k^2 + v_A^2 k^2 \sin^2\theta - 4\pi G\rho_0 F(kR)
\label{eq:dispersion}
\end{equation}
where $F(kR)$ encodes the cylindrical geometry and $\theta$ is the field--axis angle. This is a \emph{schematic static-field} relation that treats $\mathbf{B}$ as a fixed background and \emph{cannot be used quantitatively} for a magnetised cylindrical equilibrium; in particular it does not describe a current-carrying (helical) equilibrium. We use it only as an analytic guide, so this paper contains \emph{no quantitative model of magnetic tension}; it nonetheless fixes the \emph{sign} of each idealised case, the magnetic term $v_A^2k^2\sin^2\theta$ vanishing for a \emph{longitudinal} field (reducing to the field-free thermal relation) and being maximal for a \emph{perpendicular} one ($\theta=90^\circ$, which lengthens the axial mode). Consistent with this, our longitudinal runs keep the growth-rate wavelength at order $\lambda_{\rm J}$ across $\beta=0.3$--$1.0$ ($f=1.1$; the $\beta=1.0$ ladder gives $\lambda_{\rm turn}=1.14\,\lambda_{\rm J}$, Table~\ref{tab:wave}) with no resolved $\beta$-trend, so the scale is set by the thermal Jeans length, not the field, and the shortening we measure is the dynamical/contraction shift, not magnetic tension. A genuinely shorter magnetic wavelength would require a force-balanced \emph{helical}/toroidal field, whose stability is a cylindrical magnetostatic-equilibrium problem that Eq.~(\ref{eq:dispersion}) cannot address. The full case-by-case discussion of the longitudinal, perpendicular and helical geometries, and the sense in which the simulations \emph{constrain only the specific initial field geometries simulated} rather than dispose of magnetic alternatives, is deferred to Section~\ref{sec:why_mhd_fails}. The seeded helical run tested here is an \emph{out-of-equilibrium pinch}, not a force-balanced equilibrium, and does \emph{not} rule out the helical mechanism (Section~\ref{sec:optional_tests}).

\subsection{Hierarchical Fragmentation}

Filament fragmentation is hierarchical: filaments fragment into velocity-coherent fibres, and fibres into cores \citep{Hacar2013,Hacar2018,Yang2024}. Whether the fibre-to-core spacing is quasi-periodic at the classical scale is not established; the ALMA study of \citet{Yang2024} finds a clear filament--fibre--core hierarchy but only a weak quasi-periodic imprint at the core level. If HGBS filaments are bundles of fibres each fragmenting near the classical scale, measuring core spacings across the whole filament compresses the apparent spacing, in the observed direction. We test this quantitatively in Section~\ref{sec:fibre_projection}, where a \emph{toy geometrical model} of $N$-fibre bundles run through our ACS pipeline reproduces the observed $s_{\rm ACS,med}/W$; it emerges as a viable but not unique explanation (single-filament dynamical fragmentation is also consistent), and confirming it requires fibre-resolved core-spacing analysis beyond the single Orion~B region of \citet{Yang2024}.

\section{MHD SIMULATIONS}
\label{sec:sims}

\noindent\emph{Status of this and the following section.} The calculations below are exploratory. They are not a survey of filament fragmentation and they do not deliver a measured fragmentation wavelength: the growth-rate spectrum of a contracting filament has no resolved turnover, so every simulated number quoted here is an \emph{upper bound} $\lambda_{\rm turn}$ on where a turnover could lie. Their role is to provide an existence proof, that fragmentation from a contracting rather than an equilibrium state is not constrained to the classical wavelength, calibrated by the equilibrium-recovery control of Section~\ref{sec:validation_endtoend}. They should be read as a counterexample to a theoretical expectation, not as a prediction to be matched number-for-number against the observations.

The full campaign comprises $2251$ Athena++ runs across 18 grids (Table~\ref{tab:campaigns}), consuming of order $10^5$ CPU-hours on Google Cloud E2 infrastructure; controlled numerical experiments, not models of individual HGBS filaments. The large majority are \emph{exploratory} runs that map outcomes and timescales; the headline \emph{wavelength} result rests instead on the load-bearing \emph{production} runs (the 25-run near-critical growth-rate campaign, Table~\ref{tab:wave}, and its 8-run $128^3$--$512^3$ convergence ladder; Table~\ref{tab:production}), so the $2251$-run total is \emph{not} the statistical weight behind the wavelength. Two results emerge (Section~\ref{sec:why_mhd_fails}): the \emph{outcome} of a run (radial collapse or longitudinal beading) depends on regime and boundary conditions, so a single filament has no setup-independent fate, whereas the scale on which it fragments is bounded above at a sub-classical value ($\approx0.5$--$1\times$ the classical value; Section~\ref{sec:validation_endtoend}) for near-critical, longitudinal-field filaments (the perpendicular geometry is not numerically resolved). This section provides the evidence for both.

\subsection{Simulation Methodology}

We performed self-gravitating simulations with Athena++ v21.0 \citep{Stone2020}, solving the standard ideal-MHD equations, continuity, momentum with the Lorentz force, ideal induction, and total energy, coupled to self-gravity through the Poisson equation $\nabla^2\phi=4\pi G\rho$.

\smallskip\noindent\textbf{The production initial condition, stated in full.} An earlier version of this paper described the initial profile incorrectly, and we therefore give it explicitly. The production runs use a \emph{Gaussian} filament embedded in a uniform ambient medium,
\begin{equation}
\rho(r)=\rho_{\rm bg}+\rho_{\rm amp}\exp\!\left(-\frac{r^{2}}{2W_{\rm core}^{2}}\right),\qquad
\rho_{\rm amp}=\frac{f\,M_{\rm line,crit}}{2\pi W_{\rm core}^{2}},
\label{eq:ic_gauss}
\end{equation}
with $r^{2}=y^{2}+z^{2}$, $\rho_{\rm bg}=1$ in code units, $M_{\rm line,crit}=2c_s^{2}/G$ and $W_{\rm core}=0.3\,\lambda_{\rm J}$ the initial Gaussian half-width $\sigma_0$ of Table~\ref{tab:symbols}. The normalisation of $\rho_{\rm amp}$ is chosen so that $\int_0^\infty(\rho-\rho_{\rm bg})\,2\pi r\,{\rm d}r=f\,M_{\rm line,crit}$ exactly; the line mass therefore converges without an imposed truncation radius, and $f$ is imposed directly rather than inferred. There is no separate truncation: the filament merges into the uniform ambient medium, which also sets the confining pressure $P_{\rm bg}=\rho_{\rm bg}c_s^{2}$. The equilibrium-recovery control instead uses the Ostriker profile available in the same problem generator,
\begin{equation}
\rho(r)=\rho_{\rm bg}+\frac{f\rho_{c,\rm ost}}{\left[1+(r/W_{\rm core})^{2}\right]^{2}},\qquad
\rho_{c,\rm ost}=\frac{2c_s^{2}}{\pi G W_{\rm core}^{2}},
\label{eq:ic_ost}
\end{equation}
that is a Plummer profile of index $p=4$, whose line mass also converges. Neither production nor control run uses a $p=2$ profile; the observed filaments have $p\simeq2$, which is why the width conversion of Section~\ref{sec:width_correction} is needed and why it is not invariant.

The remaining initial fields are $P(r)=\rho(r)c_s^{2}$ (isothermal, $c_s=1$), $v_r(r)=0$ and $v_\phi(r)=0$, a uniform field $\mathbf{B}_0$ of magnitude $B_0=c_s\sqrt{2\rho_{\rm bg}/\beta}$ oriented at $\theta$ to the axis, and a longitudinal velocity seed built from twelve Kolmogorov modes, $v_x=\sum_{n=1}^{12}A_n\cos(k_nx+\phi_n)$ with $k_n=2\pi n/L_x$, $A_n\propto k_n^{-11/6}$ and random phases, normalised so that the total $\delta v/c_s$ equals the quoted seed amplitude. Domains are $L_x\times L_\perp^{2}$ with $L_x=16\,\lambda_{\rm J}$ for the growth-rate ladder, periodic in $x$ and, for the ambient runs, with the transverse boundary at $d\geq1\,\lambda_{\rm J}$ from the spine (a distance chosen after we found that closer transverse boundaries pin the collapse time). Every load-bearing run in Table~\ref{tab:production} uses Eq.~\ref{eq:ic_gauss}; the equilibrium control uses Eq.~\ref{eq:ic_ost} after explicit drag relaxation.

The filament is characterised by the thermal line-mass fraction $f=M_{\rm line}/M_{\rm line,crit}$ (Eq.~\ref{eq:mucrit}), the plasma $\beta=8\pi P/B^2$, and the seed-amplitude parameter $M_{\rm seed}=\delta v/c_s\times10^{4}$. At filament conditions ($n\sim10^4\,{\rm cm^{-3}}$, $T\sim10$\,K) $\beta=1$--$0.3$ corresponds to $B\sim15$--$30\,\mu$G, the range of dense-filament estimates \citep{Crutcher2012}. $M_{\rm seed}$ is a perturbation-spectrum label, not a physical Mach number: the base grid uses a single-mode seed $\delta v/c_s=10^{-4}$, whereas the growth-rate ladder of Section~\ref{sec:validation_endtoend} uses a $10^{-3}$ twelve-mode broadband seed. The magnetic mass-to-flux ratio scales as $\mu_\Phi/\mu_{\Phi,\rm crit}\propto f\sqrt{\beta}$ and exceeds unity for every configuration (none strictly magnetically subcritical), so the low-$\beta$ runs are magnetically \emph{supported}.

We explored a base grid ($\mathcal{M}=0.5$--$5$, $\beta=0.05$--$5$, near-critical) and a transition grid ($f=1.4$--$2.0$, $\beta=0.3$--$1.3$, $\mathcal{M}=1$--$5$) to resolve the stable--unstable boundary; full ranges, run counts and the resulting three outcome regimes (magnetically supported quasi-stability at $\beta\lesssim0.15$, radial collapse, and longitudinal fragmentation) are deferred to Appendices~\ref{sec:campaigns} and \ref{sec:supplementary_campaigns_app}. The near-critical longitudinal case that is our focus beads at $t_{\rm frag}\approx1.1$--$1.2\,t_{\rm J}$, turbulence changing the \emph{wavelength} by $<5\%$ while shortening the timescale $3$--$5\times$. The load-bearing \emph{production} runs on which the quoted wavelength rests, the near-critical growth-rate ladder and the direct supercritical ensemble, are collected in Table~\ref{tab:production}. These runs satisfy the Truelove criterion at the spacing-setting epoch ($N_J\approx5$--$7$ cells per local Jeans length, in both the axial and transverse directions on the isotropic cubic grids) but sit close to the numerical floor, only $1.25$--$1.75\times$ the $N_J=4$ minimum, which is one reason we quote only an order-$\lambda_{\rm J}$ upper bound (Appendix~\ref{sec:validation}).

\begin{table}
\caption{Load-bearing \emph{production} simulations underlying the fragmentation-wavelength result: the near-critical growth-rate ladder (campaign~15) and the direct supercritical ensemble (campaign~11). Shared parameters ($f$, $\mathcal{M}$, $\beta$, $\theta$, boundary conditions and collapse/bead time) are given in each block sub-header; per-row entries reuse the growth-rate measurements of Table~\ref{tab:wave}. These production runs, not the 2251-run exploratory total, which includes $144$ runs (the oblique-field and ambipolar-perpendicular surveys of Appendix~\ref{sec:supplementary_campaigns_app}) that are explicitly unresolved at $N_J\approx1$--$2$ and contribute no constraint at all, carry the quoted $\lambda_{\rm turn}$.}
\label{tab:production}
\footnotesize
\setlength{\tabcolsep}{4pt}
\begin{tabular}{ccccc}
\toprule
$f$ & seed & resolution & $\lambda_{\rm turn}/\lambda_{\rm J}$ & $\lambda_{\rm turn}/W$ \\
\midrule
\multicolumn{5}{@{}p{8cm}}{\emph{Near-critical growth-rate ladder} (campaign~15): $\beta{=}1.0$, $\theta{=}0^\circ$, $\mathcal{M}{=}10^{-3}$ broadband seed, ambient BC (periodic/reflecting), bead $t_{\rm frag}{\approx}1.1$--$1.2\,t_{\rm J}$}\\
1.1 & 42  & $128^3$ & 1.14 & 1.95 \\
1.1 & 42  & $256^3$ & 1.14 & 1.99 \\
1.1 & 42  & $512^3$ & 1.14 & 2.00 \\
1.2 & 42  & $256^3$ & 1.00 & 1.81 \\
1.2 & 137 & $256^3$ & 0.89 & 1.60 \\
1.5 & 42  & $256^3$ & 0.89 & 1.47 \\
1.5 & 137 & $256^3$ & 0.89 & 1.47 \\
1.5 & 42  & $512^3$ & 0.89 & 1.47 \\
1.8 & 42  & $256^3$ & 0.80 & 1.19 \\
2.0 & 42  & $256^3$ & 0.80 & 1.20 \\
\midrule
\multicolumn{5}{@{}p{8cm}}{\emph{Direct supercritical ensemble} (campaign~11, 39 runs): $f{=}1.5$--$2.0$, $\beta{=}0.5$--$2.0$, $\theta{=}0^\circ$, seeds 42/137, periodic$+$reflecting ambient BC, bead $t_{\rm bead}{\approx}0.6$--$0.75\,t_{\rm J}$}\\
\multicolumn{2}{@{}c}{ensemble} & $256^3$ & $1.0$--$1.14$ & $0.89$--$1.01^{\dagger}$ \\
\bottomrule
\end{tabular}
\\[2pt]
\footnotesize $^{\dagger}$~In beam-convolved Plummer units, $\lambda/W_{\rm fil}$ (KS $p=0.17$ between periodic and reflecting BCs); the ladder's $\lambda_{\rm turn}/W$ column is in the Gaussian formation FWHM of Table~\ref{tab:wave}.
\end{table}

\subsection{Why a contracting filament gives only an upper bound}
\label{sec:whyupper}

The numerical result rests on three physical ingredients, and because they are what makes $\lambda_{\rm turn}$ a bound rather than a measurement we set them out here rather than leaving them to the appendix.

\smallskip\noindent\emph{(i) Contraction competes with mode growth.} The classical calculation assumes a static, marginally-stable cylinder, so every unstable axial mode has unlimited time and the mode that wins is simply the one with the largest growth rate, $\lambda_{\rm m}\approx22H$. Our production filaments are not in force balance: they begin with $f\gtrsim1$ and no supporting radial velocity, so the spine contracts radially while the axial modes are still growing. The relevant comparison is between the mode growth time $\tau_{\rm grow}=\Gamma(\lambda)^{-1}$ and the time on which the background itself changes, $\tau_{\rm contract}=|{\rm d}\ln\rho_c/{\rm d}t|^{-1}$. Long-wavelength modes have the longest growth times and are the ones that fail to complete before the state they were computed for has gone; the surviving modes are those at the short end of the converged band. Global axial infall reinforces this, because material is delivered into the growing density maxima faster than an individual perturbation can detach itself, so no single mode is allowed to run away in isolation. The full argument, including why the effect should scale with the initial contraction rate, is given in Appendix~\ref{sec:nullspec}.

\smallskip\noindent\emph{(ii) That is why there is no resolved maximum.} Because the selection is set by the timescale competition rather than by the shape of $\Gamma(\lambda)$, the measured spectrum does not turn over: $\Gamma$ rises monotonically until it meets the grid scale, where the rise becomes numerical (Fig.~\ref{fig:growthrate}a). We therefore quote the converged-band edge $\lambda_{\rm turn}$, which is an \emph{upper limit} on where a physical turnover could lie, and we do not quote a peak. The equilibrium-recovery control shows that this is a property of the initial state and not of our measurement: the identical pipeline applied to a drag-relaxed Ostriker cylinder recovers the classical maximum at $22H$ to within $10$\,per cent.

\smallskip\noindent\emph{(iii) The role of the magnetic variables.} The field enters through $\beta$ and, equivalently, through the mass-to-flux ratio $\mu_\Phi/\mu_{\Phi,\rm crit}\propto f\sqrt{\beta}$, which exceeds unity for every configuration run, so none of our filaments is magnetically subcritical. For a longitudinal field the axial modes compress gas \emph{along} the field lines, so magnetic tension does no work against them and the growth-rate spectrum is close to the unmagnetised one; this is what we find, the axial spectrum being unchanged to within the fit noise between $\beta=0.3$ and $\beta=2$. The field's effect is instead radial: it resists the transverse contraction, and at $\beta\lesssim0.15$ it halts it altogether, which is the magnetically-supported quasi-stable regime of Appendix~\ref{sec:campaigns}. A field with a toroidal component would in principle act on the axial mode, but only if it is in radial force balance; a toroidal component seeded onto an already-contracting filament merely pinches it, narrowing $W$ by $\approx25$\,per cent and thereby \emph{raising} $\lambda/W$ without shortening $\lambda$ (Fig.~\ref{fig:helical_scan}). The one configuration that could genuinely shorten the axial mode, the force-balanced Fiege--Pudritz equilibrium, is not simulated here, and we are explicit in Section~\ref{sec:why_mhd_fails} that our campaign therefore does not constrain the magnetic channel.

\smallskip\noindent\emph{(iv) Bridging to observable units.} Converting a simulated $\lambda_{\rm turn}$ into something comparable with $S_{\rm local}$ or $S_{\rm global}$ requires the width to be measured the way an observer measures it. Three widths are involved (Table~\ref{tab:widthdefs}): the initial Gaussian half-width $W_{\rm core}=0.3\,\lambda_{\rm J}$, the evolved formation FWHM $W_{\rm form}=0.683\,\lambda_{\rm J}$ of the projected spine profile, and the beam-convolved Plummer FWHM $W_{\rm fil}$ that an observer fits. The bridge is a single multiplicative factor,
\begin{equation}
\frac{\lambda_{\rm turn}}{W_{\rm fil}}=\eta_W\,\frac{\lambda_{\rm turn}}{W_{\rm form}},
\qquad \eta_W\equiv\frac{W_{\rm form}}{W_{\rm fil}}=T_1,
\label{eq:etaW}
\end{equation}
with $T_1=\eta_W=0.59$--$0.78$ measured by forward-modelling the simulation output through a synthetic \textit{Herschel} beam and the same Plummer fitting used on the data (Section~\ref{sec:width_correction}). It is Eq.~\ref{eq:etaW}, and not the raw simulation number, that turns the Gaussian-normalised bound $\lambda_{\rm turn}/W_{\rm form}\approx2$ into the matched bound $S_{\rm sim}^{\rm upper}\lesssim1.0$--$1.3$ quoted throughout. The $\eta_W$ range is the dominant modelling term in Table~\ref{tab:errorbudget}, and it widens the matched bound by about a quarter.

\subsection{Ambient confinement and the direct supercritical measurement}
\label{sec:modecompetition}

The outcome for a supercritical filament is set by whether it is confined by an ambient medium, as real HGBS filaments are ($d$ denotes the axis-to-boundary distance in Jeans lengths). The main 654-run grid has no ambient medium, so its Gaussian skirt feeds an accretion channel through the periodic boundary that suppresses beading and drives radial collapse (Appendix~\ref{sec:regimes}); with ambient confinement (\texttt{filament\_ambient} generator) the same configurations fragment longitudinally, a 12-run reconciliation suite beading at physical, $\beta$-dependent collapse times ($0.55$--$0.70\,t_{\rm J}$) that converge across boundary conditions.

Ambient-confined supercritical filaments ($f=1.5$--$2.0$, $\theta=0^\circ$) fragment longitudinally: an ensemble of 39 runs (periodic and reflecting BCs, ambient confinement) develops 3--10 interior axial density maxima at a spacing of order the Jeans length. We take the wavelength from the linear growth-rate spectrum, whose low-$n$ modes are resolution-stable: no physical turnover is detected above $\approx1\,\lambda_{\rm J}$, so any turnover is constrained only to $\lambda_{\rm turn}\lesssim1.0$--$1.14\,\lambda_{\rm J}$ (an upper bound; no resolved peak), $\lesssim0.35$--$0.5$ times the classical value (not the median peak count, which is epoch- and resolution-dependent). The ambient medium confines the filament (central-to-ambient density ratio $\sim10^2$, consistent with HGBS crest-to-background contrasts; \citealt{Arzoumanian2019}), so it is representative rather than tuned. Across the 39 runs periodic and reflecting boundaries are statistically indistinguishable (KS $p=0.17$; medians $\lambda/W_{\rm fil}\approx0.89$ and $1.01$), and this ensemble median and the growth-rate result are the same physical wavelength $\lambda\approx1.1\,\lambda_{\rm J}$ in different normalisations. The ensemble beads early ($t_{\rm bead}\approx0.6$--$0.75\,t_{\rm J}$), before runaway and at moderate contrast, so with $\lambda\approx\lambda_{\rm J}$ resolved by ${\sim}32$ cells it satisfies Truelove at measurement (Appendix~\ref{sec:validation}).

Perpendicular-field filaments ($\theta=90^\circ$) do not universally spindle in the raw low-resolution grids: an 18-run ideal-MHD grid at $d=1.0$ shows $7/18$ bead (at $\beta=0.5$ and $2.0$) with an apparently short spacing ($\lambda/W_{\rm fil}\approx0.4$), but this conflicts with static linear theory (which predicts a \emph{longer} mode) and a transverse-refinement ladder to $N_J\approx170$ resolves the near-critical case into a monolithic radial spindle, identifying the short beading as a Truelove artefact (Section~\ref{sec:optional_tests}). An exploratory ambipolar-diffusion survey was also run, but it fragments only at $N_J\approx1$--$2$, below the Truelove limit, so it cannot support any conclusion and we exclude it from the argument here; it is described, with an explicit warning that it is numerically unresolved, in Appendix~\ref{sec:supplementary_campaigns}.

\subsection{Testing the resolution and magnetic-field alternatives}
\label{sec:optional_tests}
Three targeted campaigns test the concerns that most affect the interpretation: whether the longitudinal wavelength is resolution-limited, whether the perpendicular geometry produces genuine short beading, and whether a helical field, the one configuration expected to shorten the axial mode, actually does.

\emph{Resolution stability of $\lambda_{\rm turn}$.} Two refinement directions behave oppositely. Under \emph{transverse} refinement (the isotropic $128^3$--$512^3$ ladder) the short-wavelength spectrum ($\lambda\lesssim0.9\,\lambda_{\rm J}$, $n\gtrsim9$) grows with resolution and is numerical (unresolved transverse sub-fragmentation; Fig.~\ref{fig:growthrate}a), whereas the physical band ($\lambda\gtrsim1\,\lambda_{\rm J}$, $n\leq8$) is converged. Under \emph{axial-only} refinement ($512\times128^2\to1024\times128^2$) the physical band reproduces to better than $0.5\%$ ($\Gamma$ at $n=6,7,8$, i.e.\ $\lambda=1.33,1.14,1.00\,\lambda_{\rm J}$, is $17.6,18.7,18.7$ at both resolutions, plateauing without turning over), \emph{and} the short-$\lambda$ rise is unchanged, confirming it is fixed by the axial Nyquist scale, not a physical mode migrating to shorter $\lambda$. The $\lambda_{\rm turn}\lesssim1\,\lambda_{\rm J}$ upper bound is therefore stable: axial refinement does not move the physical-band plateau, and the only shorter-$\lambda$ structure is the transverse-resolution artefact already identified.

\emph{The perpendicular geometry: a resolved spindle, not short beading.} On a transverse-refinement ladder ($128^3$--$512^3$) the near-critical \emph{ideal} perpendicular filament collapses monolithically in the radial direction at every resolution, the axial mode staying at the box scale and the axial perturbation at seed level throughout the well-resolved phase (at $512^3$ the spine stays resolved with $N_J\approx170$ cells at $C\approx3$, falling only to $N_J\approx40$ at $C\approx30$; still $\approx10\times$ above Truelove; Appendix~\ref{sec:validation}). Short-wavelength axial structure ($\lambda/W\approx0.4$) appears only once the spine drops below $N_J\approx1$; refinement delays its onset to deeper collapse (Fig.~\ref{fig:perp_ladder}), the signature of a grid artefact. The perpendicular field therefore \emph{suppresses} axial fragmentation in favour of radial collapse, as static linear theory predicts (a perpendicular field lengthens the axial mode), and the short low-resolution beading is numerical. This is for \emph{ideal} MHD; the ambipolar-perpendicular case (Section~\ref{sec:modecompetition}) remains separately unconverged.

\emph{A seeded helical field pinches the width but does not shorten $\lambda_{\rm turn}$; the force-balanced case remains untested.} A divergence-free helical-field scan ($B_\phi/B_z=0$--$2$; construction and scan detail in Appendix~\ref{sec:supplementary_campaigns}) leaves the axial growth-rate plateau at $\lambda\approx1\,\lambda_{\rm J}$ for every winding and produces only a radial \emph{pinch} that narrows $W_{\rm FWHM}$ by $\approx25\%$ (raising $\lambda/W$ by reducing $W$, not by shortening $\lambda$; Fig.~\ref{fig:helical_scan}), so, being seeded rather than force-balanced, it cannot test the confining Fiege--Pudritz sausage ($m=0$) mode \citep{Fiege2000}; the one geometry that could shorten $\lambda$, which therefore remains the single relevant magnetic configuration left unconstrained (Section~\ref{sec:why_mhd_fails}).

\begin{figure}
\includegraphics[width=0.48\textwidth]{figures/fig_perp_ladder.pdf}
\caption{Ideal perpendicular-field refinement ladder ($128^3$--$512^3$; Section~\ref{sec:optional_tests}). \emph{(a)} Cells per local Jeans length $N_J$ at the beading-relevant spine contrast $C\approx3$: the collapsing spine is resolved far above the Truelove threshold ($N_J=4$, dashed) at every resolution. \emph{(b)} The spurious short-$\lambda$ axial ``beading'' emerges only once the spine drops below $N_J\approx1$, and refinement pushes that onset to ever-deeper collapse (at $512^3$ the filament stays monolithic to $N_J\approx40$ at $C\approx30$); the short-$\lambda$ onset deepening with resolution being the textbook Truelove under-resolution artefact signature, not a physical fragmentation mode.}
\label{fig:perp_ladder}
\end{figure}

\begin{figure}
\includegraphics[width=0.48\textwidth]{figures/fig_helical_scan.pdf}
\caption{Divergence-free helical-field scan ($B_\phi/B_z=0$--$2$, $512\times128^2$; Section~\ref{sec:optional_tests}). \emph{(a)} The axial growth-rate spectrum $\Gamma(\lambda)$ in the physical band is invariant to the toroidal fraction ($<6\%$, fit noise), with the converged-band edge fixed at $\lambda\approx1\,\lambda_{\rm J}$; $B_\phi/B_z=0$ (black) reproduces the pure-longitudinal spectrum, a code check. No excess growth appears at shorter $\lambda$ for any winding. \emph{(b)} The one systematic effect is a radial \emph{pinch}: the formation FWHM narrows by $\approx25\%$, with onset at $B_\phi\gtrsim B_z$. The seeded helical field therefore raises $\lambda/W$ by shrinking $W$, not by shortening $\lambda$; the force-balanced \citet{Fiege2000} pinch mode is not tested here.}
\label{fig:helical_scan}
\end{figure}

\section{The fragmentation scale and its comparison with observation}
\label{sec:validation_endtoend}

\noindent\textbf{The controlled experiment first.} The single most informative simulation result is not the size of the campaign but a two-configuration control, so we state it before anything else. We prepared a relaxed, force-balanced Ostriker cylinder and a dynamically contracting filament, and passed both through the \emph{identical} measurement pipeline. For the equilibrium the pipeline recovers the classical maximum at $22H$ (Fig.~\ref{fig:eqrecovery}); for the contracting filament the same pipeline finds no maximum at all within the trustworthy band. The difference is therefore a property of the prepared physical state and not of the estimator. Everything else in this section, including the $2251$-run campaign, is supporting exploration around that comparison.

The measured core spacing is sub-classical only when expressed against the filament FWHM, and that comparison hinges on how the width is defined. We therefore anchor it on the spacing in units of the filament FWHM measured the same way on simulations and data, and verify the whole chain end to end (Appendix~\ref{sec:val_appendix}).

First, the measurement recovers injected spacings: passing synthetic filaments with a known spacing and width through the identical pipeline returns $\lambda/W$ to better than $0.5\%$ over $\lambda/W=1.4$--$5.7$ (Appendix~\ref{sec:val_appendix}), so the observed$\to$simulated comparison depends on the width convention, not on residual measurement bias. This validates the measurement machinery given a synthetic filament definition; it does not establish that a skeleton extracted from the real sky corresponds to the theoretical filament, which is a separate and unresolved question (Section~\ref{sec:whatfilament}). Second, forward-modelling the simulations end to end through the HGBS pipeline (projection, beam convolution, Plummer fitting, and the same ACS measurement) gives the observable $\lambda/W_{\rm fil}$ directly. Two configurations bracket the result: the \emph{near-critical} growth-rate ladder forward-models to $\lambda_{\rm turn}/W_{\rm fil}\lesssim1.0$--$1.3$, while the \emph{supercritical} 39-run bead ensemble gives $\approx0.89$--$1.01$ (Table~\ref{tab:production}); both are order-$\lambda_{\rm J}$ and both exceed the analytic shortcut ($\approx0.9$, Eq.~\ref{eq:t1_apply}, which applies only a single $\sigma\!\to\!$FWHM factor). We adopt the near-critical forward-model band ($1.0$--$1.3$) as the headline like-for-like value and carry the $\approx30\%$ spread between it and the analytic/ensemble estimates as the width-conversion systematic.

The fragmentation wavelength must be measured with care: counting density peaks late in the collapse is epoch- and resolution-dependent, so we instead work from the growth-rate spectrum. Seeding the filament with a small broadband perturbation ($\delta v/c_s=10^{-3}$, twelve axial modes), we track the power $P(k,t)$ in each mode and fit its growth over its approximately exponential interval (Fig.~\ref{fig:growthcurves}); the growth accelerates as the filament collapses, so the rates describe a \emph{dynamically collapsing} filament. The growth rate $\Gamma(\lambda)$ (Fig.~\ref{fig:growthrate}, left) rises as $\lambda$ decreases and is resolution-stable for $\lambda\gtrsim1\,\lambda_{\rm J}$ (low-$n$ modes agree to $\sim5\%$ across $128^3$--$512^3$; $\Gamma$ at $n=7$ is $17.8,18.5,18.7$), rising to a \emph{plateau} at the converged-band edge ($\Gamma$ at $n=6,7,8$, i.e.\ $\lambda=1.33,1.14,1.00\,\lambda_{\rm J}$, is $17.6,18.7,18.7$), while below $\lambda\lesssim0.9\,\lambda_{\rm J}$ it rises with \emph{transverse} resolution and is numerical. Because $\Gamma$ plateaus rather than turning over, we quote no wavelength at all for these runs, only the \emph{upper bound} $\lambda_{\rm turn}\lesssim1\,\lambda_{\rm J}$ on where a turnover could lie, and do \emph{not} claim a precisely-located maximum (Appendix~\ref{sec:val_appendix}); the true peak could lie at shorter, transverse-unresolved wavelength ($\lambda_{\rm turn}\lesssim1\,\lambda_{\rm J}$), so we use $\lambda_{\rm turn}$ below only as shorthand for this upper bound. It is an order-$\lambda_{\rm J}$, sub-classical value: $\lambda_{\rm turn}/W_{\rm FWHM}\approx1.2$--$2.0$ (near-critical $1.5$--$2.0$, falling to $\approx1.2$ in the supercritical runs; Table~\ref{tab:wave}), a factor $\approx2.5$--$5.0$ below the classical $\lambda_{\rm m}\approx5$--$6\,W_{\rm FWHM}$ (in $W$ units; \citealt{Fischera2012}) and at or below the critical $2.6\,W_{\rm FWHM}$ across $f=1.1$--$2.0$. The like-for-like beam-convolved Plummer bound (below) is $\lambda_{\rm turn}/W\lesssim1.0$--$1.3$, and it is this that should be set against the data (Table~\ref{tab:wave}).\footnote{The near-critical formation FWHM is $W_{\rm FWHM}=0.58\,\lambda_{\rm J}$, narrower than the supercritical $W_{\rm form}=0.683\,\lambda_{\rm J}$ because a near-critical filament has already begun to contract when its mode is set (both tabulated in Table~\ref{tab:widthdefs}).}

\begin{figure}
\includegraphics[width=0.47\textwidth]{figures/fig_growthcurves.pdf}
\caption{Power in three axial modes versus time for a near-critical filament ($f=1.1$). Each grows with an approximately exponential interval (the fit window, shaded) whose slope gives the growth rate; the growth steepens as the filament collapses, so the extracted rate characterises a dynamically collapsing filament. The representative mode at the converged-band edge, $n=7$, grows faster than the long-wavelength $n=3$, giving the sub-classical $\lambda_{\rm turn}$; the grid-scale $n=12$ is shown for reference.}
\label{fig:growthcurves}
\end{figure}

\begin{figure*}
\includegraphics[width=0.70\textwidth]{figures/fig_growthrate.pdf}
\caption{\emph{(a)} Growth-rate spectrum $\Gamma(\lambda)$ of the axial modes for a near-critical filament ($f=1.1$) at three resolutions. For $\lambda\gtrsim1\,\lambda_{\rm J}$ the curves overlap (resolution-converged); at the grid scale ($\lambda\lesssim0.9\,\lambda_{\rm J}$, shaded) they separate and rise with resolution, the signature of unresolved sub-fragmentation. $\Gamma$ plateaus (does not turn over) at the converged-band edge, so $\lambda_{\rm turn}\lesssim1\,\lambda_{\rm J}$ is an \emph{upper bound}, not a resolved maximum. \emph{(b)} $\lambda_{\rm turn}/W_{\rm FWHM}$ versus line-mass fraction (Gaussian FWHM), well below the classical $5$--$6$ at every $f$. The like-for-like comparison uses a beam-convolved Plummer width $\sim1.5\times$ broader, giving a matched bound $\lambda_{\rm turn}/W\lesssim1.0$--$1.3$ (the $T_1$ conversion is itself uncertain over $0.59$--$0.78$, which widens this band by about a quarter) (grey band): this, not the Gaussian $\approx2$, is what should be set against the observed $s_{\rm ACS,med}/W\approx0.4$--$1.9$.}
\label{fig:growthrate}
\end{figure*}

The result is one-sided. The simulations rule out the assertion that a contracting filament \emph{must} fragment only at the equilibrium-cylinder wavelength, and they establish the \emph{absence} of the required classical-scale ($5$--$6\,W$) maximum. They do not determine the preferred wavelength. In width-free form $\lambda_{\rm turn}\lesssim1.0$--$1.14\,\lambda_{\rm J}$ against the classical $\lambda_{\rm m}=22H=2.3\,\lambda_{\rm J}$, a ratio of $\approx0.35$--$0.5$ across $f=1.1$--$2.0$, but $1.0\,\lambda_{\rm J}$ is the converged-band edge rather than a resolved physical maximum, so we quote it only as an order-$\lambda_{\rm J}$ bound. The load-bearing near-critical runs have $N_J\approx5$--$7$ cells per local Jeans length when the mode is set, $1.25$--$1.75$ times the Truelove minimum, and the fit window ($\delta_{\rm rms}<0.05$) keeps $N_J>4$ at all three resolutions (Appendix~\ref{sec:validation}). That criterion guards against \emph{artificial} fragmentation, and it is a different requirement from convergence of $\Gamma(k)$ itself, which we therefore test separately. For the mode carrying the bound ($n=7$, $\lambda=1.14\,\lambda_{\rm J}$) the fitted rate is $\Gamma=17.8$, $18.5$ and $18.7\,t_{\rm J}^{-1}$ at $128^3$, $256^3$ and $512^3$, so $\Delta\Gamma/\Gamma=3.9$ per cent between the first pair and $1.1$ per cent between the second. We accordingly define the trustworthy band by a numerical criterion, $\Delta\Gamma/\Gamma<5$ per cent between successive resolutions, rather than by the Truelove condition alone. That criterion is met for $\lambda\gtrsim1\,\lambda_{\rm J}$ and fails below it, which is the same boundary we identify from the resolution-divergent small-scale rise, so the two diagnostics agree on where the spectrum can be trusted. Satisfying it does not make an unresolved turnover appear, which is why we still claim only a bound. To be quantitative about what could still be hidden: the physical band is converged to $\Delta\Gamma/\Gamma<5$ per cent down to $\approx1\,\lambda_{\rm J}$, and $\Gamma$ is still rising there at $\approx15$ per cent per mode step, so a true maximum could lie anywhere between the opacity-limited floor and $\approx1\,\lambda_{\rm J}$ without contradicting these data. The bound should be read with that floor in mind. Two readings of the missing turnover should be distinguished. It may be a grid-resolution limit, in which case a genuine maximum lies below the converged band and would need transverse resolution beyond $512^3$ to expose; scaling from the present ladder, resolving a turnover at $0.5\,\lambda_{\rm J}$ at the same cells-per-Jeans-length would require of order $1024^3$. Or it may be physical: a filament that accretes and contracts while it fragments has no single quasi-static background, so the notion of one sharply preferred wavelength may not apply, and a broad plateau is then the expected outcome rather than a failure to resolve a peak. Our calculations cannot separate these, and we quote a bound for that reason. The no-turnover result is not a mode-quantisation artefact: the $128^3$--$512^3$ ladder resolves a different maximum axial wavenumber at each resolution, yet the plateau and its edge do not move.

Two caveats temper the causal attribution. First, the fitted rates ($\Gamma\approx18$--$26\,t_{\rm J}^{-1}$) are the accelerating growth of a collapsing background, not quasi-static eigenvalues, so we use only the mode \emph{ordering}. Second, a $p=4$ (isothermal-equilibrium) profile gives the same order-$\lambda_{\rm J}$ wavelength ($\approx0.8\,\lambda_{\rm J}$; Appendix~\ref{sec:val_appendix}), so the deficit is not a profile-index artefact but obtained for the contracting configurations tested here (Section~\ref{sec:limitations}).

A plausible physical interpretation follows from evaluating the classical scale at the \emph{instantaneous central density}, $22H(\rho_{\rm eff})\approx0.8\,\lambda_{\rm J}(\rho_0)$: the instability remains approximately classical with respect to the instantaneous central density. The sub-classical ratio $\lambda/W$ arises because the observable FWHM does not contract in step with $\rho_c$, since the extended envelope lags the collapsing spine. Consistent with this, $\lambda_{\rm turn}\lesssim1.1\,\lambda_{\rm J}$ sits at or below the analytic \emph{cutoff} $\lambda_{\rm cr}\approx1.2\,\lambda_{\rm J}$ of our own eigensolver (Appendix~\ref{sec:val_appendix}), not the classical \emph{peak} $2.3\,\lambda_{\rm J}$: this is the expected signature of a contracting filament sliding towards marginal stability, not a re-measurement of the equilibrium-cylinder cutoff (the growing modes are near the marginal scale of the \emph{contracting}, elevated-density state). The equilibrium-recovery control confirms this reading, returning the classical $5$--$6$ when the filament instead begins from equilibrium.

\emph{A time-dependent mode-selection estimate.} This picture can be sharpened from a counterexample into an approximate \emph{consistency check} (not an independent prediction, since the freeze-out density is read from the measured growth rates) by a WKB/self-similar treatment of fragmentation during contraction, in the spirit of \citet{Larson1985} and \citet{Clarke2016,Clarke2017}. As the filament contracts, its central density $\rho_c(t)$ rises, so the instantaneous classical fastest-growing wavelength shrinks as $\lambda_{\rm m}(t)=22H[\rho_c(t)]=2.3\,\lambda_{\rm J}(\rho_0)\,(\rho_c/\rho_0)^{-1/2}$, while the instantaneous growth rate rises as $\Gamma\propto(4\pi G\rho_c)^{1/2}$. A perturbation is \emph{selected}, imprinted non-linearly, not at the marginally-stable equilibrium scale but at the epoch $t_\star$ when the fastest-growing mode \emph{of that instantaneous classical relation} completes its $\sim\!N_e$ e-foldings from seed to non-linearity within the residual contraction time (the standard collapse-fragmentation condition $\tau_{\rm grow}\lesssim\tau_{\rm contract}$; \citealt{Larson1985}). This fixes a \emph{freeze-out} density $\rho_{\rm eff}$, which we infer from the measured rates: $\Gamma_{\rm meas}\approx18$--$26\,t_{\rm J}^{-1}$ versus the base-density value $\Gamma(\rho_0)\approx6.3\,t_{\rm J}^{-1}$ (Appendix~\ref{sec:val_appendix}) gives $\rho_{\rm eff}/\rho_0=(\Gamma_{\rm meas}/\Gamma(\rho_0))^2\approx8$--$17$; this is corroborated \emph{independently} by the directly-measured central-density contrast at the mode-setting epoch ($\rho_c/\rho_0\approx8$--$17$; Appendix~\ref{sec:val_appendix}), which de-circularises the estimate. The instantaneous classical relation then gives $\lambda_{\rm select}=2.3\,\lambda_{\rm J}(\rho_0)/\sqrt{\rho_{\rm eff}/\rho_0}\approx0.6$--$0.8\,\lambda_{\rm J}$, which \emph{abuts} (and slightly under-predicts) the measured upper bound $\lambda_{\rm turn}\approx0.8$--$1.1\,\lambda_{\rm J}$, consistent in sign, and in size only where the ranges overlap ($\approx0.8\,\lambda_{\rm J}$). The same argument predicts $\lambda_{\rm select}$ to \emph{decrease} with contraction rate, more supercritical filaments contract faster, reach higher $\rho_{\rm eff}$, and freeze at shorter $\lambda$; consistent with the $f$-ladder ($f=1.1$: $\Gamma\approx18$, $\lambda_{\rm turn}=1.14$; $f=1.8$--$2.0$: $\Gamma\approx25$, $\lambda_{\rm turn}=0.80\,\lambda_{\rm J}$; Table~\ref{tab:wave}), although the $f=1.8$--$2.0$ values lie below the converged band edge and are themselves unconverged upper bounds, so this is an upper-bound trend rather than a measured dispersion relation. The sub-classical \emph{ratio} follows because the spine density rises by $\rho_{\rm eff}/\rho_0\approx8$--$17$ (so $H$ falls by $\sqrt{8\text{--}17}\approx3$--$4$) while the observable envelope FWHM contracts only modestly ($W_{\rm form}\approx0.58$--$0.68\,\lambda_{\rm J}$, a factor $\lesssim1.3$): the deficit is precisely this spine--envelope contraction lag. We label this a hypothesis to be tested by the contraction ladder proposed below rather than an established scaling. It is falsifiable, $\lambda_{\rm select}\propto\rho_{\rm eff}^{-1/2}$ with $\rho_{\rm eff}$ set by the contraction rate, that a future relaxation/contraction ladder can test directly.

The bound sits within the methodological envelope of the observed conditional spacing (Fig.~\ref{fig:envelope}), and the robust deliverable is the density-independent ratio $\lambda/W$ with its equilibrium-recovery calibration. That control is the load-bearing evidence (Fig.~\ref{fig:eqrecovery}): the \emph{identical} pipeline recovers the classical turnover at $22H$ for a drag-relaxed equilibrium cylinder ($\lambda_{\rm peak}/22H=1.07$--$1.10$ at both resolutions, a genuinely resolved maximum) but finds no turnover for the contracting runs. We state the limitation here rather than in an appendix, because this control is what licenses the whole interpretation: it is a two-point comparison, one relaxed equilibrium against one contracting run, and it establishes the \emph{sign} of the effect, not its functional dependence on contraction rate. A minimum adequate ladder would be five otherwise identical filaments initialised with radial velocity $v_r=A\,v_{r,0}(r)$ for $A=0,0.25,0.5,0.75,1$ at $256^3$ with a $512^3$ check at the endpoints, measuring $\Gamma(k)$ and $\lambda_{\rm turn}/22H[\rho_c]$ for each. Until that exists the result is a counterexample, not a calibrated law. The underlying competition between the mode growth time and the contraction time is set out explicitly in Appendix~\ref{sec:nullspec}. This is a numerical instance of the generic expectation that a still-contracting filament fragments below the marginally-stable hydrostatic mode \citep{Clarke2016,Clarke2017,Clarke2020}, not a new mechanism; our distinct contributions are the longitudinal near-critical demonstration, that control, and the homogeneous Gaia-DR3 reanalysis. We note the limits of this correspondence: the single-filament simulations produce a \emph{quasi-periodic} axial comb (3--10 near-uniformly spaced beads), whereas the observed cores are non-periodic and spatially inhomogeneous (Section~\ref{sec:clustering}), rejecting both periodic and Poisson models. The simulations therefore reproduce a \emph{scale} (sub-classical, order $\lambda_{\rm J}$) but neither the observed point process nor its non-periodicity. Stated at its strongest, what the numerical work establishes is that sub-equilibrium fragmentation scales are dynamically possible; it does not reproduce the observed core point process and therefore does not identify the origin of the observed HGBS spacing distribution. We would put the point more strongly: a successful theory of filament fragmentation must predict a point process, not merely a wavelength, and reproducing $s_{\rm ACS}/W$, $L/N$, $F(\ell)$ and the gap distribution simultaneously is the outstanding challenge. No calculation presented here does so. The simulations establish only that non-equilibrium gravitational fragmentation \emph{can} generate structure below the equilibrium-cylinder wavelength, not that it produces the observed spatial statistics.

\begin{figure}
\centering
\includegraphics[width=0.99\columnwidth]{figures/fig_eqrecovery.png}
\caption{Equilibrium-recovery control (the primary validation that the sub-classical spectrum is not generated by the wavelength-extraction pipeline and is associated with the non-equilibrium initial state). The identical growth-rate \emph{measurement} pipeline is applied to a \emph{relaxed, force-balanced} pressure-confined Ostriker equilibrium (blue; prepared by explicit drag relaxation, then evolved drag-free) and to a \emph{dynamically contracting} supercritical filament (orange); the two differ in their prepared initial state, not in the measurement, isolating the effect of starting from equilibrium versus contraction. \emph{Left:} onset fragmentation power spectrum; the equilibrium peaks at the classical $\lambda_{\rm m}\approx22H$ ($\approx5.4\,W_{\rm FWHM}$) whereas the contracting filament places its shortest converged scale at $\lesssim1.4\,W_{\rm FWHM}$, without a resolved maximum. \emph{Right:} the dispersion relation $\Gamma(\lambda)$ versus $\lambda/22H$; the equilibrium shows the classical peak-and-decline with its maximum at $22H$, whereas the contracting filament's $\Gamma$ rises monotonically. The pipeline thus recovers the classical wavelength when the physics is classical.}
\label{fig:eqrecovery}
\end{figure}

The physically appropriate comparison is between the simulated fragmentation scale, which produces \emph{one bead per wavelength}, and the observed \emph{global length-per-core scale} $(L/N)/W\approx3.4$--$5.8$ (Section~\ref{sec:linedensity}), the quantity a one-fragment-per-wavelength theory actually predicts, rather than the intra-group $s_{\rm ACS,med}/W$. On $L/N$ the simulated upper bound ($\lambda_{\rm turn}/W_{\rm fil}\lesssim1.0$--$1.3$) is in fact \emph{more} discrepant: the simulation makes about one bead per $\lambda_{\rm turn}$, i.e.\ a denser core production than the observed filament-averaged $L/N$, so single-filament fragmentation \emph{over}-produces cores per unit length unless most of the filament is core-poor; precisely the uneven-occupancy picture (Section~\ref{sec:clustering}). The shorter intra-group $s_{\rm ACS,med}/W\approx1.4$ is then the \emph{intra-group adjacent-spacing scale}, not the global fragment-production comparison, so the near-coincidence of the simulated $1.0$--$1.3$ with it is not the key agreement. Single-filament dynamical fragmentation is thus at best marginally sufficient: it establishes consistency, not sufficiency, and hierarchical-fibre projection (Section~\ref{sec:fibre_projection}) and magnetic tension remain viable within the factor-$\sim$2 systematics. The claim is falsifiable: it is refuted if resolved single filaments or individual fibres show a core line-density no higher than the classical one-core-per-$\lambda_{\rm m}$, and disfavoured if the true observed $s_{\rm ACS,med}/W$ fell well below $\sim1$ (single-filament fragmentation would then \emph{over}-produce cores), so it is operationally viable only between these limits.

\emph{Opacity-limited fragmentation.} Because $\Gamma(\lambda)$ does not turn over within the converged band, an isothermal calculation contains no physical small-scale thermodynamic cutoff and therefore cannot predict a characteristic core mass or a unique CMF peak; a physical turnover, which would convert the isothermal plateau into a physical peak, requires a non-isothermal EOS ($\gamma_{\rm eff}>1$ at $n\gtrsim10^{5}$--$10^{6}\,{\rm cm^{-3}}$) or radiative physics (Section~\ref{sec:cmf}), which is a change in thermal behaviour distinct from, and at much lower density than, the classical opacity limit at which the gas becomes optically thick to its own cooling radiation \citep{Rees1976,LowLyndenBell1976}. Consistent with this, our $\gamma=5/3$ tests suppress fragmentation (Appendix~\ref{sec:validation}), confirming the true short-wavelength cutoff lies below our isothermal grid. Such a departure from isothermality occurs once dust--gas coupling and line cooling can no longer hold $T$ constant against compressional heating; an effective $\gamma_{\rm eff}\gtrsim1$ regime reached at $n\gtrsim10^{5}$--$10^{6}\,{\rm cm^{-3}}$ for these clouds \citep{Larson1985,Glover2012}, above the $n\sim10^{4}\,{\rm cm^{-3}}$ of the fragmenting filaments modelled here. Radiation-hydrodynamics or a piecewise polytropic EOS with $\gamma>1$ at high density is therefore needed to set a physical fragmentation scale and a CMF peak; our isothermal suite deliberately does not include it.

\section{DISCUSSION}
\label{sec:discussion}

\subsection{Why the spacing matters: fragment masses and the CMF}
\label{sec:cmf}
The spacing matters because, with the line mass, it sets the characteristic fragment mass that seeds the core mass function (CMF) and ultimately the IMF: for a filament of line mass $M_{\rm line}$ fragmenting at spacing $\lambda$ the mean fragment mass is $M_{\rm frag}\sim\epsilon\,M_{\rm line}\,\lambda$, and a near-critical $M_{\rm line}\approx16\,M_\odot\,{\rm pc^{-1}}$ with $\lambda\approx0.1$--$0.2$\,pc gives $\sim0.5$--$1\,M_\odot$ for $\epsilon\sim0.3$, consistent with the HGBS CMF peak while the classical $5$--$6\times$FWHM would overshoot it. This is only an order-of-magnitude \emph{consistency check}: it is degenerate in $\epsilon M_{\rm line}$ (only the product enters) and yields a \emph{mean}, not a peaked distribution. An \emph{isothermal} EOS lacks a small-scale thermodynamic cutoff and therefore cannot predict a characteristic core mass or a unique CMF peak; only a non-isothermal EOS ($\gamma_{\rm eff}>1$ at $n\gtrsim10^{5}$--$10^{6}\,{\rm cm^{-3}}$) would convert the isothermal plateau into a physical peak (Section~\ref{sec:validation_endtoend}).

\subsection{Hierarchical-fibre projection is consistent with the sub-classical spacing}
\label{sec:fibre_projection}

A second, purely geometric route reaches the same sub-classical value without modified physics, and rests on independently observed substructure: HGBS filaments are bundles of $N$ velocity-coherent \emph{fibres} \citep{Hacar2013,Hacar2018,Yang2024}. If each fibre fragments near the classical wavelength but core spacing is measured \emph{across} the bundle, as any filament-level pipeline does, the interleaving of cores from $N$ fibres compresses the apparent spacing. The premise, quasi-periodic fibre-to-core fragmentation at the classical scale, is a hypothesis, not an established measurement (\citealt{Yang2024} find only a weak imprint), so we present projection as viable but unconfirmed.

We tested this with a \emph{toy geometrical model} (not a physical forward model) that reuses our exact ACS measurement. Each fibre is a cylinder of width $W_{\rm fib}$ fragmenting at its \emph{own} classical wavelength $\lambda_{\rm fib}\approx5$--$6\,W_{\rm fib}$ (\citealt{Fischera2012}); we populate a filament with $N$ such fibres (random phase, $\sim30\%$ lognormal scatter), project onto the shared spine, and measure in units of the \emph{filament} width $W_{\rm fil}$. This assumes quasi-hydrostatic fibres, which sits uneasily with our own result (a near-critical fibre would itself fragment sub-classically, doubling the required $r/N$), so we treat the classical-fibre $(5$--$6)r/N$ result as an upper bound and note the fibre and single-filament pictures are not independent. Two effects compress the spacing, narrower fibres ($r\equiv W_{\rm fib}/W_{\rm fil}<1$) and several of them, giving in the interleaving limit
\begin{equation}
\frac{\lambda}{W_{\rm fil}} \;\approx\; \frac{(5\text{--}6)\,r}{N},
\label{eq:fourrn}
\end{equation}
verified across $r$ and $N$ in Fig.~\ref{fig:fibre_degen}. Either effect alone can produce a sub-classical value: a single fibre of half the filament width ($r=0.5$, $N=1$) gives $\lambda/W_{\rm fil}\approx2.5$--$3.0$, as do two full-width fibres ($r=1$, $N=2$).

The observation constrains only the combination $r/N$, not $N$ individually (Fig.~\ref{fig:fibre_degen}): the DisPerSE and FilFinder conventions map to $r/N\approx0.07$--$0.35$, so the relation cannot be inverted for $N$ even approximately, and we do not attempt it: any $r/N$ in that range will do. A further limitation is that we quote only the central recovered value, whereas the projected spacing for a given $(N,r)$ is a broad distribution whose width depends on fibre phase, length, orientation and overlap; since the observed gap distributions are themselves broad, that scatter matters and we have not modelled it. The robust statement is only that unresolved fibre multiplicity can reduce the filament-level adjacent spacing below the classical prediction with no modified physics.

\begin{figure}
\includegraphics[width=0.48\textwidth]{figures/fig_fibre_degeneracy.pdf}
\caption{Hierarchical-fibre geometrical toy model, showing the filament-level $\lambda/W_{\rm fil}$ recovered by our ACS pipeline as a function of the fibre-to-filament width ratio $r=W_{\rm fib}/W_{\rm fil}$ and the number of fibres $N$; the model follows $\lambda/W_{\rm fil}\approx(5$--$6)r/N$ (Eq.~\ref{eq:fourrn}). White contours mark the observed fixed-width and region-specific values, the orange contour the FilFinder value, and the red band the observed fibre-width range ($r\approx0.3$--$0.5$). The observation constrains the combination $r/N$, not $N$ alone: the sub-classical value is reproduced by a single narrow fibre, by a few wider fibres, or by intermediate combinations.}
\label{fig:fibre_degen}
\end{figure}

Hierarchical projection is therefore \emph{complementary} to, not exclusive of, single-filament dynamical fragmentation (Section~\ref{sec:validation_endtoend}): the two are multi-scale processes that can operate together, and unresolved fibre multiplicity is observationally \emph{degenerate} with dynamical contraction in the filament-averaged spacing. Breaking that degeneracy requires the tracer that resolves the fibres directly; high-resolution kinematics (ALMA/NOEMA) that decompose the cores into their parent velocity-coherent fibres across the primary regions (Section~\ref{sec:limitations}).

A coherence point remains: our single-filament simulations fragment into a \emph{quasi-periodic} comb at $\lambda_{\rm turn}$, whereas the observed cores show no comb and are over-dispersed on parsec scales (Section~\ref{sec:clustering}). The two reconcile only if sub-fragmentation, multi-fibre interleaving, and projection erase the periodicity into the observed near-exponential gap distribution, as \citet{Clarke2020} find. We do not simulate this stage, so the overlap between a single-filament $\lambda_{\rm turn}$ and the observed line density is an order-of-magnitude consistency within the methodological envelope, not a morphological match.

A third route to the same reconciliation is post-formation dynamical evolution: cores accrete, migrate along the filament and interact over several free-fall times (the migration displacements of $0.01$--$0.05$\,pc are already in our error budget; \citealt{Kirk2016,Mattern2018}), so dynamical relaxation of an initially quasi-periodic chain could erase the comb into the observed inhomogeneous distribution with no modified physics; a candidate on the same footing as fibre projection, testable by evolving the near-critical runs further past first fragmentation, which we have not done.

\emph{A candidate mechanism for the two-scale (local sub-classical / global near-classical) pattern.} The paper's most novel empirical result, short intra-group spacing coexisting with a near-classical filament-averaged line density, calls for a physical origin, which we can only propose here. The most economical is \emph{spatially non-uniform seeding}: real filaments inherit an inhomogeneous spectrum of density and velocity perturbations (from turbulence, accretion streamers, and fibre substructure), so gravitational fragmentation runs to non-linearity preferentially in the perturbation-rich, locally supercritical stretches, which then bead at the (sub-classical) contracting-filament scale, while low-amplitude stretches stay barren over the same interval. The result is exactly the observed picture: core-rich groups with sub-classical intra-group spacing, separated by core-poor filament, and a filament-averaged $L/N$ close to one core per classical wavelength. Sequential/triggered fragmentation and multi-fibre interleaving (Section~\ref{sec:fibre_projection}) produce the same two-scale signature. This is testable with the existing runs: our \emph{uniform-seed single-filament} simulations produce a comb, not this two-scale pattern, so a direct test is to seed the near-critical runs with an inhomogeneous (e.g. turbulent) perturbation field and measure whether the simulated core chains reproduce both the sub-classical adjacent spacing \emph{and} the near-classical $L/N$; a comparison we identify as the natural next step (Section~\ref{sec:limitations}).

\subsection{Field geometry, boundary conditions and the perpendicular case}
\label{sec:why_mhd_fails}

The common argument that most filaments carry perpendicular fields (from the Planck tendency for dense filaments to lie perpendicular to the large-scale field; \citealt{Planck2016}) uses a projected orientation that need not equal the internal geometry, and hourglass morphologies with the field turning longitudinal inside the crest are observed \citep{Jadhav2026}. A resolved perpendicular filament spindles rather than fragmenting axially (Section~\ref{sec:optional_tests}), so it does not produce the observed axial core chains and is not a source of tension.

Where perpendicular geometry does apply, three non-exclusive resolutions of the residual deficit remain: hierarchical fibre fragmentation (bundles of near-critical fibres projected across the bundle; \citealt{Hacar2013,Hacar2018,Yang2024}); non-ideal MHD at later evolution, which narrows the factor-$\sim3$ perpendicular deficit; and an observational lag in which the cores are relics of an earlier episode. These are separated by the future tests listed in Section~\ref{sec:limitations}.

\emph{The one untested magnetic configuration.} Our scan constrains only the specific initial field geometries simulated (Table~\ref{tab:magcases}), and the status of the four cases is simply stated. A uniform longitudinal field was tested in the adopted non-equilibrium configuration and leaves the axial mode essentially unchanged. For a static perpendicular field the apparent short fragmentation is resolution-sensitive and is not established. A seeded toroidal component pinches the filament but is not a stability calculation for a helical equilibrium. The force-balanced Fiege--Pudritz helical/toroidal equilibrium, which is the one configuration that could genuinely shorten the axial mode, is untested here. We therefore describe magnetic tension as neither favoured nor disfavoured by these calculations, and we state plainly that the $2251$-run campaign does not constrain the magnetic channel at all: none of those runs contains a force-balanced helical equilibrium. Readers should not take the size of the campaign as evidence bearing on this mechanism. For scale only, and read off from the published \citet{Fiege2000} results rather than computed here, a strongly toroidally dominated equilibrium shortens the $m=0$ wavelength by up to a factor $\approx2$; our $\beta\approx0.3$--$2$ places these filaments at the weak end of that sequence, so such a field could be astrophysically significant, though its magnitude cannot be assessed without a force-balanced calculation matched to the observed filament parameters. Settling it needs a validated cylindrical MHD eigensolver rather than the algebraic relation of Section~\ref{sec:magnetic}: one would construct a Fiege--Pudritz self-similar equilibrium with poloidal $B_z(r)$ and toroidal $B_\phi(r)$ in radial force balance, confirm it is stationary by drag-relaxing and then evolving drag-free exactly as in our equilibrium control, and only then seed the $m=0$ mode. That is a paper in itself.

\begin{table}
\caption{Schematic summary of the four magnetic-field configurations and their tested status.}
\label{tab:magcases}
\footnotesize
\setlength{\tabcolsep}{4pt}
\begin{tabular}{@{}p{2.9cm}lp{3.0cm}@{}}
\toprule
Configuration & Tested? & Effect on axial $\lambda$ \\
\midrule
Uniform longitudinal & tested & unchanged (term vanishes) \\
Static perpendicular (resolved) & tested & suppresses axial fragmentation (spindle) \\
Seeded/dynamic helical & tested & pinches $W$; does not shorten $\lambda$ \\
\midrule
\multicolumn{3}{l}{\emph{Not simulated in this work; listed for completeness}}\\
Force-balanced helical (Fiege--Pudritz) & not tested & could shorten by $\lesssim2\times$ (order-of-magnitude inference from Fiege--Pudritz, not a result of this work; Section~\ref{sec:why_mhd_fails}) \\
\bottomrule
\end{tabular}
\end{table}

\emph{The infall connection.} A dynamically fragmenting filament fragments below the equilibrium wavelength only while it is still contracting, so the interpretation ultimately rests on the filaments' accretion state. Per-cloud infall-velocity or mass-growth-rate constraints that would place Orion~B, Aquila, Perseus and Taurus filaments inside or outside the simulated non-equilibrium range are not established for these specific clouds, so the simulation--cloud correspondence is a plausibility argument, not a demonstrated match.

Sub-isothermal equations of state ($\gamma<1$) accelerate fragmentation by $30$--$40\%$ without changing the wavelength, while adiabatic heating ($\gamma=5/3$) suppresses it entirely (Appendix~\ref{sec:validation}). The isothermal suite is therefore the conservative case in a globally averaged sense; though not locally, where a $\sim12\%$ dust-temperature variation shifts the local $\lambda/W$ by $\sim10$\,per-cent and relocates where cores form (Section~\ref{sec:mt_variation}).

\subsubsection{Non-Ideal MHD Effects}
\label{sec:nonideal}

At the densities probed here ($n_0\sim10^4$\,cm$^{-3}$), the ambipolar-diffusion time ($\sim10$--$20\,t_{\rm ff}$; \citealt{Tassis2004}) is long compared with the collapse time, so ideal MHD is a good approximation for the longitudinal fragmentation stage: four non-ideal runs at $\eta_{\rm AD}=\{0.01,0.05\}$ still collapse radially, $\sim12$--$15\%$ \textit{earlier} than the ideal control (flux leakage erodes magnetic support, the conservative direction). Non-ideal effects nonetheless change the outcome for perpendicular fields, where ambipolar diffusion promotes longitudinal fragmentation (Section~\ref{sec:modecompetition}).

\subsection{Limitations and future work}
\label{sec:limitations}

Two limitations bound our conclusion. First, the characteristic $\lambda/W$ and its scatter ($\sigma\approx0.4$) rest on only four nearby, low-mass Gould Belt regions of broadly solar metallicity and similar radiation field, so the result is established only for that regime; whether it persists in high-mass complexes, at lower metallicity or under stronger irradiation is untested, and extending the pipeline to further HGBS regions and independent surveys (e.g.\ JCMT/SCUBA-2) is needed. Two validation checks remain \emph{outstanding}: the $\lesssim5\%$ manual-editing bound is a \emph{proxy} pending a blind, multi-editor re-skeletonisation, and a dedicated faint-segment injection--recovery test is not yet done. To date the adjacent-core-spacing and skeleton pipelines have \emph{not} been independently re-run by anyone other than the author; the deposited code and catalogues are released to enable that external check. Two analyses would further sharpen the comparison: (i) a \emph{per-segment hierarchical normalisation}; the interval-midpoint local-width test is already done (Section~\ref{sec:linedensity}; medians $\Delta s/W_{\rm local}\approx0.6$--$1.1$, all sub-classical), but the cleaner unbranched-segment version is deferred because the bridged skeletons over-fragment at spurious junctions (Section~\ref{sec:linedensity}); and (ii) a \emph{line-mass dependence} test; splitting the sample into sub-/near-/super-critical bins would test whether the most supercritical filaments show the smallest $s/W$, as the dynamical picture predicts, but reliable segment-level line masses are not available at the required fidelity, so it is left open (the absence of such an $f$-trend would count against the dynamical interpretation). Second, the simulations are idealised controlled experiments, not models of specific filaments: the periodic-BC main grid does not represent the turbulent, magnetically coupled boundaries of a real filament, and our ambient-confined configuration is an idealisation of that confinement.

\smallskip\noindent\emph{A recommendation for observers.} Because the factor-$\sim$2 skeleton systematic is the dominant limit on any comparison with theory (Section~\ref{sec:whatfilament}), we recommend that, to make cross-algorithm and cross-resolution spacing comparisons reproducible, filament skeletons be extracted to a \emph{fixed column-density contrast threshold relative to the local background}, rather than to an absolute level, and that the spacing be quoted for \emph{both} a persistent-crest (DisPerSE-like) and a medial-axis (FilFinder-like) definition so that the factor-$\sim$2 ontological span is stated explicitly rather than hidden in a single number. A hybrid, matched-threshold pipeline combining the two definitions is the natural community step toward an algorithm-independent ``filament''.

On the causal attribution: our seeded filaments contract from the start, so what we demonstrate is generic ``fragmentation from a non-equilibrium, contracting state'' rather than the specifically accretion-driven case of \citet{Clarke2016,Clarke2017}; the equilibrium-recovery control is only a two-point comparison and a relaxation ladder is needed to establish causation (Appendix~\ref{sec:val_appendix}). Further idealisations are: (i) an \emph{isothermal} EOS (locally uncertain at $\sim10$--$20\%$; Section~\ref{sec:mt_variation}); (ii) \emph{periodic transverse boundaries}, whose accretion channel suppresses fragmentation unless ambient confinement is imposed (Section~\ref{sec:modecompetition}); (iii) \emph{simplified initial profiles}; and (iv) a \emph{prescribed Kolmogorov turbulence spectrum} rather than driven transonic turbulence, leaving the wavelength unchanged ($<5\%$) but shortening the timescale $3$--$5\times$ (Appendix~\ref{sec:supplementary_campaigns}).

The fragmentation \emph{classification} is resolution-independent (the collapse \emph{time} itself carries an $8.5\%$ offset between $128^3$ and $256^3$ and is not claimed converged; Appendix~\ref{sec:validation}) and domain-doubling changes the collapse time by $<2\%$, but the transition zone is sampled with only 2 seeds per point. Non-linear core merging could in principle inflate the apparent supercritical spacing by up to $\sim50\%$ over a few Myr, but we have verified it does \emph{not} bias the quoted value: $\lambda/W_{\rm fil}$ is measured at maximum bead count (fragmentation onset), and in $38/39$ runs the peak count is still rising at the final snapshot while no run shows the peak-count decline that would signal merging. The ambipolar wavelengths are likewise fragmentation-onset measurements (the runs reach gravitational runaway before a converged plateau is established; Section~\ref{sec:modecompetition}).

\subsection{Non-isothermal effects and the temperature-gradient test}
\label{sec:mt_variation}

The MHD suite is isothermal, whereas HGBS dust-temperature maps vary by $\sim12$\,per-cent ($\sim10$--$13$\,K), a local perturbation since $\lambda_{\rm J}\propto\sqrt{T}$. A single illustrative run with an imposed $\sim12$\,per-cent longitudinal temperature gradient fragments \emph{differentially} (the cool, locally supercritical half beads while the warm half is suppressed), shifting the local wavelength by $\sim10$\,per-cent. This is a single illustrative run, not a converged parameter study, and we quote it as an indication of the direction and rough size of the effect only; temperature sets both the Jeans scale and the critical line mass ($\mu_{\rm crit}\propto T$). The isothermal assumption is thus conservative only in a globally averaged sense; a self-consistent treatment is deferred to future work.

\section{CONCLUSIONS}

We have investigated how cores are spaced along the selected filament crests in HGBS clouds, combining a Gaia DR3 observational re-analysis, an MHD campaign and a hierarchical-fibre geometrical toy model. Our primary result, in each of the four clouds examined, is that the core distribution is not a simple periodic fragmentation comb, and that the local and global spacing statistics differ strongly. The along-skeleton pair-correlation shows no comb at $5$--$6\,W$ or at any scale, and core occurrence is strongly inhomogeneous. That inhomogeneity, however, is largely inherited from the parent filament: against a null conditioned on local column density, only Taurus is clumped in excess of its own structure. Because that null is built partly from column density that the cores themselves raise, it is biased against detecting excess clumping, so the three non-detections are conservative. Measured within core-bearing stretches the adjacent separation is short, a factor $\approx2.4$--$4.3$ below $\lambda_{\rm m}$ and robust in \emph{sign} under every convention. Across the full skeleton the length per core approaches the classical value, but much of that skeleton is subcritical. Restricting the comparison to filament above the critical line mass moves the length per core downwards, though by less than the raw calculation suggests once the cores' own contribution to the selecting column density is removed: on a core-masked background the eligible statistic sits between the local and the global one, remaining sub-classical in the two clouds (Orion~B and Aquila) that retain usable core samples, while Perseus and Taurus become too sparse to decide (Section~\ref{sec:eligible}). The apparent local/global tension is therefore partly, but not wholly, a question of which filament was ever eligible to fragment; these eligible-length numbers are provisional pending region-specific kinematics.

\begin{itemize}
    \item \textbf{The deficit of the smallest separations is not shown to be physical.} The inferred exclusion scale, $0.011$--$0.044$\,pc, lies entirely within the range spanned by the \textit{Herschel} beam and the \textsc{getsources} deblending limit across our distances. It is therefore consistent with being wholly instrumental, we can place no lower limit on a physical contribution, and the gap-model fits should not be read as measuring a physical hard core. Settling this requires synthetic source pairs injected into the maps and re-extracted through the HGBS pipeline, which we have not done.

    \item \textbf{The cores are not a periodic comb, and occurrence is inhomogeneous.} No statistically significant classical-scale periodicity is detected (Fig.~\ref{fig:paircorr}) and both periodic and pure-Poisson gap models are rejected. Core occurrence is strongly inhomogeneous ($6$--$27$ per cent occupancy). Relative to a column-density-conditioned null, Taurus is the only one of the four clouds in which the cores are clumped beyond the structure of the filament they sit on. We deliberately do not phrase this as ``three of four clouds show no excess clumping'': with $N=4$, and with a null that is conservative by construction, the correct statement is that we detect excess clumping in one cloud and fail to detect it in the other three, which is not the same as establishing its absence there. The clouds differ from one another (Table~\ref{tab:percloud}). With four independent units our power to generalise is weak: a single discrepant cloud is $25$ per cent of the sample, and we could not have distinguished a genuine population split from chance unless roughly half the clouds behaved differently. ``Three of four'' should therefore be read as a description of this sample, not as a population fraction.
    \item \textbf{The local (conditional) spacing is sub-classical; the global line density is not.} Measured inside core-bearing stretches, the adjacent-core (ACS) spacing is sub-classical under every skeleton and width convention, a factor $\approx2.4$--$4.3$ across the spacing-estimator injection--recovery-validated DisPerSE conventions, more under FilFinder, the \emph{sign} robust, the \emph{magnitude} set by the skeleton algorithm (DisPerSE vs FilFinder, a factor of two), the dominant limit on any comparison with theory. The primary observed result is the convention envelope, $s_{\rm ACS,med}/W\simeq0.8$--$1.9$ for DisPerSE and $\simeq0.4$--$1.9$ including FilFinder.
    \item \textbf{The global length per core is near-classical only if subcritical filament is counted.} The core line density $N/L$, expressed through its reciprocal $L/N$, is only mildly sub-classical (classical in Perseus, near-classical in Taurus, sub-classical only in Orion~B and Aquila), so the local deficit reflects uneven occupancy (core-poor filament) plus the mean/median skew, not a globally elevated number of cores per unit length. The prestellar-only subset ($s_{\rm ACS,med}/W=2.10$ vs $2.08$ all-cores) confirms it is not protostellar contamination.
    \item \textbf{The all-pairs median is the wrong estimator.} It is a region-extent-biased estimator of local spacing, converging to $\sim L_{\rm fil}/3$; the ACS spacing is the estimator we adopt. The pipeline recovers injected spacings to $<0.5\%$ and, on a relaxed force-balanced equilibrium cylinder, returns the classical $22H\approx5$--$6\,W_{\rm FWHM}$, whereas for contracting filaments the same pipeline finds no resolved maximum down to the converged limit, implying $\lambda_{\rm turn}/W_{\rm FWHM}\lesssim1.4$ (equilibrium-recovery control; Fig.~\ref{fig:eqrecovery}).
    \item \textbf{The simulations are a counterexample, not an explanation.} They demonstrate that sub-equilibrium fragmentation scales are dynamically possible, but they do not reproduce the observed core point process and therefore do not identify the origin of the observed HGBS spacing distribution. A contracting filament is not constrained to fragment at $5$--$6\,W$ but reaches a resolution-stable upper bound $\lambda_{\rm turn}\lesssim1\,\lambda_{\rm J}$ with no resolved turnover; a deliberate consequence of the isothermal EOS, which lacks a small-scale cutoff and cannot predict a characteristic core mass or CMF peak. On the line-density comparison this upper bound \emph{over}-produces cores unless most of the filament is core-poor, consistent with the uneven-occupancy picture; the numerical bound $\lambda_{\rm turn}/W\lesssim1.0$--$1.3$ gives order-of-magnitude consistency only.
    \item \textbf{The physical cause is undetermined.} Single-filament dynamical fragmentation, hierarchical-fibre projection ($\lambda/W_{\rm fil}\approx(5$--$6)r/N$, complementary not exclusive) and helical magnetic tension all remain viable and current data cannot discriminate them. A resolved ideal perpendicular field drives monolithic radial collapse (the short beading a Truelove artefact) and a seeded helical field only pinches $W$, so the force-balanced Fiege--Pudritz helical/toroidal equilibrium is the important \emph{untested} candidate.
    \item \textbf{Decisive future tests.} Fibre-resolved ACS spacing on the full HGBS catalogues with high-resolution kinematics (ALMA/NOEMA) and polarimetry of filament interiors, a perpendicular-field resolution study, and MHD runs with a spatially varying temperature field (a single illustrative gradient run already shows the isothermal assumption can be locally non-conservative; Section~\ref{sec:limitations}).
    \end{itemize}

\smallskip\noindent With only four nearby, low-mass clouds, these results are a characterisation of the Gould Belt filament population, not a universal fragmentation law; extending the homogeneous ACS analysis to further HGBS regions and independent surveys is the natural next step.

\section*{ACKNOWLEDGMENTS}

We thank the Herschel Gould Belt Survey team for the datasets. The Athena++ simulations were run on commercially procured cloud infrastructure (Google Cloud E2), funded from the author's institutional research allocation; the full set of $2251$ runs (Appendix~\ref{sec:campaigns}) consumed of order $10^5$ CPU-hours.

\section*{DATA AVAILABILITY}

The supporting archive, (i) the full run manifest (all 2251 runs with their $(f,\beta,\mathcal{M},\theta,\eta_{\rm AD})$ parameters and outcomes); (ii) the Athena++ problem-generator source (\texttt{filament\_ambient.cpp}, including the divergence-free helical construction) and normalisation scripts; (iii) HDF5 snapshots of the growth-rate convergence subset ($128^3$--$512^3$) and the equilibrium-recovery control (remaining snapshots on request); and (iv) the analysis pipeline (DisPerSE v0.9.24 parameters, the bridging/deblending and adjacent-core-spacing scripts, and the derived $\lambda/W$ catalogues); is deposited in a \emph{public}, openly inspectable repository and available \emph{now} (\texttt{github.com/Tilanthi/ASTRA-dev}, under \texttt{simulations/}), with a citable Zenodo DOI (\texttt{10.5281/zenodo.14872301}) reserved for the accepted release and cited here, and versioned to the accepted release on publication; the repository is public and inspectable at submission rather than only on acceptance. The adjacent-core-spacing (ACS) estimator was re-implemented internally by the author (a second, from-scratch implementation of the estimator) and cross-checked against the deposited counts and medians (Table~\ref{tab:nnpipeline}); this is an internal consistency check, not third-party verification. We invite referees and readers to run the deposited pipeline; the released code and catalogues are the appropriate route to independent external verification, which has not yet been carried out. To be explicit about what has and has not been independently reproduced: the adjacent-core spacing, the pipeline counts of Table~\ref{tab:nnpipeline} and the region medians have been re-derived by that second implementation; the skeletonisation, the injection--recovery corrections, the MHD growth-rate extraction and the forward-modelling to observable units rest on a single pipeline and have not been re-run by anyone else. We would welcome that check. HGBS column-density and skeleton maps are available from the Herschel Science Archive.

\bibliographystyle{mnras}
\appendix

\section{The limited HGBS sample}
\label{sec:limited_sample}
The four \emph{limited} regions fail at least one of the primary-selection criteria (Section~\ref{sec:distance_uncertainties}) and are excluded from the ACS measurement; they are retained only for the legacy pairwise-median comparison (Table~\ref{tab:literature}). Table~\ref{tab:sample_limited} lists them together with the full 8-region core-count-weighted mean.

\begin{table}
\caption{The four \emph{limited} HGBS regions (excluded from the primary ACS measurement) and the full 8-region core-count-weighted mean. Columns as in Table~\ref{tab:sample}.}
\label{tab:sample_limited}
\footnotesize
\setlength{\tabcolsep}{3pt}
\begin{tabular}{lcccc}
\toprule
Region & $D$ & $N_{\rm core}$ & Pairwise med. & S.E. \\
       & (pc) &  & (pc) & (pc) \\
\midrule
Ophiuchus& 137 & 513   & 0.206 & 0.053 \\
Serpens  & 458 & 194   & 0.331 & 0.097 \\
TMC1     & 135 & 178   & 0.195 & 0.056 \\
CRA      & 150 & 239   & 0.248 & 0.072 \\
\midrule
\textbf{Full-sample w. Mean}$^\dagger$ &  & \textbf{5,069} & \textbf{0.279} & \textbf{0.009} \\
\bottomrule
\end{tabular}
\par\vspace{2pt}
\footnotesize $^\dagger$Core-count-weighted mean over all eight regions. The $0.009$\,pc figure is a formal core-count-weighted standard error (treating cores as independent); it is \emph{not} the relevant uncertainty, since the four primary clouds are the independent units (Section~\ref{sec:error_budget}), and plays no role in the comparison with theory.
\end{table}

\section{The $L/3$ bias of the pairwise-median spacing}
\label{sec:pwmedian}
For $N$ points uniformly distributed on a filament of length $L$, a single pairwise distance $D=|X_i-X_j|$ has the triangular PDF $f(d)=2(L-d)/L^2$, with mean $\langle D\rangle=L/3$ and median $d^\star=L(1-1/\sqrt{2})\approx0.293\,L\approx0.88\,(L/3)$. The median of all $N(N-1)/2$ pairwise distances converges to $d^\star$ as $N\to\infty$, for a random or a perfectly periodic distribution alike, so the pairwise median measures the region extent, not the fragmentation wavelength. A Monte-Carlo over $N=5$--$500$ gives median$/L\approx0.293\,(1+0.43/N)$, i.e.\ $\le2\%$ inflation for $N\gtrsim20$, so the convergence is reached at the core counts of interest. For branched skeletons the statistic measures $\approx0.3\times$ the component's arclength extent, and for a realistic filament length distribution it tracks the rms-length$/3$ (full tables in the Zenodo archive). Computing the pairwise median over \emph{all} cores in each primary region (without per-filament segmentation) returns $\sim L_{\rm region}/3$ ($12.5$, $7.5$, $8.0$, $2.5$\,pc for Orion~B, Aquila, Perseus, Taurus), the region extent, confirming directly that the statistic measures segmentation scale rather than the fragmentation wavelength.

\section{Validation of the measurement and the width normalisation}
\label{sec:val_appendix}

This appendix documents the checks summarised in Section~\ref{sec:validation_endtoend}. To verify the routines extracting $\lambda$ and $W$ carry no hidden normalisation bias, we passed synthetic filament cubes with \emph{known} transverse FWHM and axial bead spacing through the identical pipeline; the recovered $\lambda/W$ matches the input to better than $0.5\%$ over $\lambda/W_{\rm in}=1.4$--$5.7$ (Fig.~\ref{fig:pipeval}), confirming the pipeline recovers injected spacings with no measurable normalisation bias, and validating the width bookkeeping of Eq.~(\ref{eq:t1_apply}).

\begin{figure}
\includegraphics[width=0.46\textwidth]{figures/fig_pipeline_validation.pdf}
\caption{Closed-loop validation: the pipeline recovers the injected $\lambda/W$ to within $0.4\%$ over $\lambda/W=1.4$--$5.7$.}
\label{fig:pipeval}
\end{figure}

\emph{Growth-rate measurement and convergence.} Seeding eight ambient-confined filaments ($f=1.1,1.2,1.5$; $\beta=1$; two seeds) with a broadband perturbation ($\delta v/c_s=10^{-3}$, twelve axial modes) and fitting each mode's growth over $\delta_{\rm rms}=3\times$seed$\to0.3$ gives $\Gamma(\lambda)$ (Fig.~\ref{fig:growthcurves}). It is resolution-converged for $\lambda\gtrsim1\,\lambda_{\rm J}$ (a plateau, not a turnover), while modes at $\lambda\lesssim0.9\,\lambda_{\rm J}$ grow with transverse resolution and are numerical; the converged edge does not migrate, so $\lambda_{\rm turn}\lesssim1\,\lambda_{\rm J}$ is resolution-stable and an independent seed reproduces it to $3\%$. Across the eight runs $\lambda_{\rm turn}=0.9$--$1.1\,\lambda_{\rm J}$ ($\lambda_{\rm turn}/W_{\rm FWHM}=1.5$--$2.0$; Table~\ref{tab:wave}), so $\lambda_{\rm turn}/\lambda_{\rm IM92}=0.4$--$0.5$; extending to $f=1.8,2.0$ gives the same scale, covering the HGBS range. Periodic and reflecting boundaries are indistinguishable (KS $p=0.17$).

\emph{Comparison with linear theory, and the gravity solver.} Figure~\ref{fig:dispersion} overplots the from-scratch--solved linear dispersion relation of the isothermal equilibrium cylinder (Ostriker, $m=0$; reproducing \citealt{Nagasawa1987,Inutsuka1992,Fischera2012} to $<1\%$), which peaks at $\lambda_{\rm IM92}=22H=2.3\,\lambda_{\rm J}$ and cuts off at $\approx1.16\,\lambda_{\rm J}$, whereas the collapsing filament's spectrum rises monotonically through that cutoff without turning over. Self-gravity uses a periodic FFT with reflecting transverse hydrodynamic boundaries; a transverse domain-doubling and an isolated-boundary re-solution of Poisson's equation both leave $\lambda_{\rm turn}$ unchanged ($<1\%$), so the selected wavelength is unbiased by the gravity boundary treatment.

\begin{figure}
\includegraphics[width=0.48\textwidth]{figures/fig_dispersion.pdf}
\caption{The from-scratch--solved eigenvalue dispersion relation of the isothermal equilibrium cylinder (\citealt{Nagasawa1987,Inutsuka1992}), validated to $<1\%$ against \citet{Nagasawa1987} and \citet{Fischera2012} (not a fit): it peaks at $\lambda_{\rm IM92}\approx2.3\,\lambda_{\rm J}$ and vanishes at the cutoff $\lambda_{\rm cr}\approx1.2\,\lambda_{\rm J}$, overplotted on the measured growth-rate spectrum ($f=1.1$, three resolutions, normalised at $\lambda=1.14\,\lambda_{\rm J}$). The measured spectrum instead rises monotonically into the grid-scale region (shaded), confirming the small-scale rise is numerical, so the measured $\approx1.1\,\lambda_{\rm J}$ wavelength is an order-$\lambda_{\rm J}$, sub-classical value, not a precisely-located maximum.}
\label{fig:dispersion}
\end{figure}

\emph{Nonlinear bead spacing and timescale separation.} Evolved into the non-linear regime the near-critical runs develop discrete beads at $\lambda_{\rm bead}\approx1.1$--$1.6\,\lambda_{\rm J}$ (before merging reduces the count), overlapping the growth-rate range $\lambda_{\rm turn}=0.89$--$1.14\,\lambda_{\rm J}$ (Table~\ref{tab:wave}). Although the growth is fast ($\tau_{\rm growth}/\tau_{\rm bg}\approx0.3$--$0.8$), the spectrum is fit-window invariant: restricting to the early, near-static phase ($\delta_{\rm rms}<0.05$, $N_J\gtrsim5$) reproduces the full-window spectrum to $<1\%$, so the selected wavelength is set at the low-contrast start of the window.

\emph{Pipeline benchmark against the analytic Jeans rate.} A near-uniform self-gravitating box ($\rho\simeq\rho_0$) reproduces the analytic uniform-medium rate $\Gamma(k)=\sqrt{4\pi G\rho_0-c_s^2k^2}$ to $1$--$7\%$ for the resolved modes, confirming the machinery; the large filament rates ($\Gamma\approx18$--$26$) then simply reflect the \emph{local} free-fall rate at the $\sim8$--$17\times$ elevated central density ($\Gamma(\rho_0)\approx6.3$; consistent with the freeze-out $\rho_{\rm eff}$ of Section~\ref{sec:validation_endtoend}), not a super-free-fall anomaly.

\emph{Robustness to line-mass, initial profile and field strength.} The order-$\lambda_{\rm J}$ result is not specific to one configuration: across $\beta=0.3$--$1.0$ ($f=1.1$) $\lambda_{\rm turn}$ stays order $\lambda_{\rm J}$ (the $\beta=1.0$ ladder gives $1.14\,\lambda_{\rm J}$, Table~\ref{tab:wave}; run-to-run scatter $\approx0.1$--$0.3\,\lambda_{\rm J}$, no resolved $\beta$-trend), and an Ostriker ($p=4$) profile gives $\lambda_{\rm turn}=0.80\,\lambda_{\rm J}$ at $128^3$ and $256^3$, so the scale is not a profile-index artefact; though that run too fragments from a contracted state, confirming the sub-classical scale as a generic outcome of non-equilibrium fragmentation (the relaxed control below tests a true equilibrium).

\emph{Recovery of the classical limit for a genuine equilibrium; the primary control.} The critical test is whether the \emph{identical} pipeline recovers the classical $\lambda_{\rm m}=22H\approx5$--$6\,W_{\rm FWHM}$ for a force-balanced equilibrium cylinder. It does. Because the periodic FFT solver enforces the Jeans swindle at the box-mean density, the analytic Ostriker profile is not the solver's equilibrium (initialised directly it breathes by factors of two), so we relaxed a pressure-confined Ostriker profile to hydrostatic balance with an explicit velocity drag, verified by a no-perturbation control (central density drifts only $\sim3\%$, in the $k=0$ mode), and seeded it with the identical 12-mode perturbation ($L_x=16\,\lambda_{\rm J}$, two resolutions). Its $\Gamma(\lambda)$ shows the textbook peak-and-decline with a genuinely resolved maximum at $\lambda_{\rm peak}\approx22H$: $\lambda_{\rm peak}/W_{\rm FWHM}=5.0$--$5.5$, $\lambda_{\rm peak}/22H=1.07$--$1.10$ (converged; $22H/W_{\rm FWHM}=4.96$--$5.03$ matches the analytic Ostriker $5.07$ to $\approx1\%$); the $\approx7$--$10\%$ long peak is a conservative bias. The \emph{same} pipeline on the contracting filaments gives $\lambda_{\rm turn}\lesssim1\,\lambda_{\rm J}$, i.e.\ $\lambda_{\rm turn}/W_{\rm FWHM}\lesssim1.4$, ($\approx0.4\times22H$) with a monotonically rising $\Gamma$ (Fig.~\ref{fig:eqrecovery}). The two differ only in their \emph{prepared initial state}, not the drag-free measurement, so the pipeline is demonstrably unbiased and the difference is associated with the non-equilibrium contracting state in these experiments; the closed-loop synthetic test independently recovers injected $\lambda/W$ up to $5.7$ to $<0.5\%$. This is, however, a \emph{two-point} comparison (one relaxed equilibrium versus one contracting run): it establishes the \emph{sign} of the contraction effect but not its functional dependence on the degree of non-equilibrium, for which a relaxation ladder, a near-equilibrium filament grown by controlled accretion, is needed to establish causation. The control validates the \emph{sign} of the deficit; the contracting $\lambda_{\rm turn}$ itself remains an upper limit.

\begin{table}
\caption{Growth-rate measurements. $\lambda_{\rm turn}$ is the converged-band-edge \emph{upper bound} on where a turnover in $\Gamma(\lambda)$ could lie ($\Gamma(\lambda)$ shows no resolved turnover; the entry is not a measured peak), $\Gamma_{\rm peak}$ its growth rate (in units of $t_{\rm J}^{-1}$), $W$ the per-run Gaussian formation FWHM ($\approx0.55$--$0.67\,\lambda_{\rm J}$; all lengths in $\lambda_{\rm J}$). The $f=1.8,2.0$ rows extend the same method into the supercritical regime (single resolution) and give the same sub-classical scale as the bead-count ensemble ($\lambda/\lambda_{\rm J}\approx0.98$), so the two methods agree at the $\sim\lambda_{\rm J}$ level; their $\lambda_{\rm turn}\approx0.80\,\lambda_{\rm J}$ lies just below the converged-band edge ($\approx0.9$--$1.0\,\lambda_{\rm J}$), so it too is an upper bound and the apparent $\lambda_{\rm turn}(f)$ decline is an upper-bound trend, not a resolved dispersion relation. The discrete axial modes ($\lambda=L_x/n$, $\Delta\lambda\approx0.15\,\lambda_{\rm J}$ near $\lambda\sim1$) give $\lambda_{\rm turn}$ a quantisation uncertainty $\pm0.07\,\lambda_{\rm J}$; the final quoted digit is not significant, and the box cannot be lengthened to refine it because $L_x\geq24\,\lambda_{\rm J}$ introduces a seeding artefact.}
\label{tab:wave}
\footnotesize
\setlength{\tabcolsep}{4pt}
\begin{tabular}{cccccc}
\toprule
$f$ & resolution & seed & $\lambda_{\rm turn}$ (u.b.) & $\Gamma_{\rm peak}$ & $\lambda_{\rm turn}/W$ \\
\midrule
1.1 & $128^3$ & 42  & 1.14 & 17.8 & 1.95 \\
1.1 & $256^3$ & 42  & 1.14 & 18.5 & 1.99 \\
1.1 & $512^3$ & 42  & 1.14 & 18.7 & 2.00 \\
1.2 & $256^3$ & 42  & 1.00 & 19.1 & 1.81 \\
1.2 & $256^3$ & 137 & 0.89 & 20.1 & 1.60 \\
1.5 & $256^3$ & 42  & 0.89 & 21.2 & 1.47 \\
1.5 & $256^3$ & 137 & 0.89 & 22.1 & 1.47 \\
1.5 & $512^3$ & 42  & 0.89 & 21.5 & 1.47 \\
1.8 & $256^3$ & 42  & 0.80 & 24.6 & 1.19 \\
2.0 & $256^3$ & 42  & 0.80 & 26.3 & 1.20 \\
\bottomrule
\end{tabular}
\end{table}

A complementary end-to-end grid of 21 ambient-confined runs ($f=1.5$--$3.0$, $\beta=0.5$--$2.0$, two seeds; campaigns~11--13), forward-modelled through beam convolution and Plummer fitting, gives a consistent picture: $\lambda/\lambda_{\rm J}(\rho_0)\approx0.9$--$1.2$, with boundary condition, profile and seed each moving the value by $\lesssim40\%$ and a factor-$1.6$ contrast variation shifting it by only $\sim0.3$, so the wavelength is a property of the filament, not the confinement.

% ===== Reintegrated supporting appendices =====
\section{Simulation Validation Details}
\label{sec:validation}

To ensure robustness against numerical effects we performed three validation tests (83 simulations). For resolution, six parameter points run at both $128^3$ and $256^3$ with matched problem-generator settings, initial conditions, seeds and boundaries all fragment at both resolutions; the mean timescale ratio is
\begin{equation}
\frac{t_{\rm frag}(256^3)}{t_{\rm frag}(128^3)} = 0.915 \pm 0.012,
\end{equation}
an $8.5\%\pm1.2\%$ offset (from the ratio $0.915\pm0.012$), all points within a $\pm11\%$ band (Fig.~\ref{fig:resolution}), so the fragmentation classification is resolution-independent. Replacing the Gaussian-profile initial condition with uniform density yields identical classifications at all 10 tested points.

A fragmentation-wavelength measurement is trustworthy only if the local Jeans length stays resolved ($\gtrsim4$ cells; Truelove et al.\ 1997) up to the measurement epoch, since an under-resolved Jeans length spuriously enhances small-scale power; the grid-scale, resolution-divergent rise we identify as numerical (Fig.~\ref{fig:growthrate}a). Tracking the minimum $N_J$ along each collapse (Table~\ref{tab:truelove}): for the \emph{longitudinal} runs the inter-bead wavelength ($\lambda\approx\lambda_{\rm J}$, $\sim32$ cells) is set at $C\approx20$--$40$ where $N_J\approx5$--$7$, dropping below 4 only at $C>170$, after the spacing is set; for the perpendicular \emph{ambipolar} runs, beading is measured only at $C\approx250$--$560$ where $N_J\approx1$--$2$ (Truelove violated), so those short values are resolution-limited upper limits. We therefore base the quoted wavelength on the growth-rate spectrum ($N_J\gtrsim5$, converged), not late-time peak counts. Because the production grids use isotropic (cubic) cells, $N_J$ is identical in the axial and transverse directions; the load-bearing longitudinal runs have $N_J\approx5$--$7$ in \emph{both}, while the axial-only refinement run ($1024\times128^2$; Section~\ref{sec:optional_tests}) raises the \emph{axial} $N_J$ to $\gtrsim10$ with the transverse value held at $\approx5$--$7$ and leaves the physical-band plateau unmoved, confirming the plateau is fixed by the transverse-resolved physics rather than by axial mode quantisation. The explicit rule we adopt is that a perpendicular-field filament must keep $N_J>4$ cells per local Jeans length \emph{throughout} the advanced collapse to suppress grid-scale numerical fragmentation \citep{Truelove1997}; the short perpendicular beading appears only once this criterion is violated ($N_J\lesssim1$--$2$), which is exactly why we classify it as a numerical artefact (Section~\ref{sec:optional_tests}).

\begin{table}
\caption{Jeans (Truelove) resolution: minimum cells per local Jeans length $N_J$ at the epoch the spacing is measured, and the density contrast $C$ at that epoch.}
\label{tab:truelove}
\footnotesize
\setlength{\tabcolsep}{3pt}
\begin{tabular}{@{}p{3.0cm}p{1.9cm}p{2.5cm}@{}}
\toprule
Configuration & $C$ at meas. & min $N_J$ \\
\midrule
Near-critical longitudinal (calibration) & $20$--$40$ & $5$--$7$ (OK) \\
Supercritical longitudinal ensemble & early beading, moderate $C$ & $\gtrsim4$ (OK) \\
Ideal perp., converged ladder ($512^3$; \S\ref{sec:optional_tests}) & $\approx3$ / $\approx30$ & $170$ / $40$ (OK; monolithic) \\
Ideal perp., short-$\lambda$ artefact onset & $>10^2$ (grid-dep.) & $\lesssim1$ (violated) \\
Perp.\ ambipolar (tension side) & $250$--$560$ & $1$--$2$ (violated) \\
\bottomrule
\end{tabular}
\end{table}

\begin{figure}
\includegraphics[width=0.48\textwidth]{figures/figR1_resolution_scatter.pdf}
\caption{Resolution convergence with matched problem generator: $t_{\rm frag}(256^3)$ versus $t_{\rm frag}(128^3)$ for six parameter points, all within a $\pm11\%$ band (mean ratio $0.915\pm0.012$).}
\label{fig:resolution}
\end{figure}

\subsubsection{Equation of State Sensitivity}

Testing mildly sub-isothermal equations of state ($\gamma=0.9,0.8$) at 5 representative points, $\gamma<1$ systematically accelerates fragmentation (Fig.~\ref{fig:eos_sensitivity}): at $\gamma=0.8$ all 5 points fragment at $t_{\rm frag}\approx0.66$--$0.68\,t_{\rm J}$, roughly half the isothermal timescale, while $\gamma=0.9$ delays it (1 of 5 fragments by the integration limit). The point is that $\gamma<1$ shortens $t_{\rm frag}$ monotonically; we draw no stability inference from the finite integration window.

The opposite limit confirms the mechanism: an \emph{adiabatic} equation of state ($\gamma=5/3$; campaign~8, 30 runs) suppresses fragmentation entirely (none to $t>30$--$40\,t_{\rm J}$), because once the gas becomes adiabatic at a density contrast $C\gtrsim$ a few, compressive heating restores Jeans support before the axial mode can grow. At HGBS filament conditions ($n\approx10^4\,{\rm cm^{-3}}$, $T\approx10$\,K) efficient dust--gas thermal coupling and molecular-line cooling hold the gas close to isothermal, with an effective polytropic index $\gamma_{\rm eff}\lesssim1$ in this density--temperature regime \citep{Larson1985,Glover2012}; the real gas is therefore expected to have $\gamma\leq1$. Our tests demonstrate that $\gamma\leq1$ shortens the fragmentation \emph{time} $t_{\rm frag}$; they do \emph{not} demonstrate a change in the fragmentation \emph{wavelength} (the near-critical $\lambda_{\rm turn}$ is unchanged within the mode-quantisation step). We therefore do \emph{not} claim the isothermal suite is an upper bound on $\lambda/W$; we claim only that real filaments, being $\gamma_{\rm eff}\lesssim1$ in this density--temperature regime, are if anything more fragmentation-prone in \emph{time}. The $\gamma<1$ and $\gamma=5/3$ limits bracket the isothermal case and identify the equation of state as what would set the small-scale fragmentation cutoff relevant to the opacity limit (Section~\ref{sec:limitations}).

\begin{figure}
\includegraphics[width=0.48\textwidth]{figures/fig6_phase3_eos_sensitivity.pdf}
\caption{Equation of state sensitivity test. Left: fragmentation time or final simulation time as a function of $\gamma$ (all parameter points). Right: $t_{\rm frag}/t_{\rm final}$ versus $\gamma$ with lines connecting the same $(f, \beta, \mathcal{M})$ parameter point across $\gamma$ values. Sub-isothermal EOS ($\gamma < 1$) systematically reduces $t_{\rm frag}$.}
\label{fig:eos_sensitivity}
\end{figure}

\section{Specification of the inhomogeneous null and the freeze-out argument}
\label{sec:nullspec}

\emph{The column-density-conditioned null.} The null used in Sections~\ref{sec:intermittency} and \ref{sec:clustering} asks whether cores are clumped \emph{beyond} the clumping already present in the filament they lie on. It is an inhomogeneous Poisson process on the one-dimensional skeleton with intensity
\begin{equation}
\lambda(x)\;\propto\;\left[\max\!\left(N_{\rm H_2}(x)-N_0,\,0\right)\right]^{\alpha},
\label{eq:nullintensity}
\end{equation}
where $x$ is arclength along the spine, $N_{\rm H_2}(x)$ the column density of the skeleton pixel at $x$, $N_0=10^{21}\,{\rm cm^{-2}}$ a floor that prevents low-column pixels from dominating, and $\alpha$ a single exponent fitted per sample. The normalisation is fixed by requiring $\int\lambda\,{\rm d}x=N_{\rm obs}$, so each realisation contains the observed number of cores. The exponent is fitted by maximum likelihood on the observed core positions, and, to avoid the circularity of fitting and testing on the same data, we refit it under leave-one-cloud-out cross-validation: for each cloud, $\alpha$ is estimated from the other three and applied to the held-out cloud. The fitted values are $\alpha=1.6$--$2.4$ with a cross-validated mean of $2.0$, and the conclusions are unchanged over $\alpha=1$--$3$.

For each cloud we draw $10^4$ realisations, compute the same Fano factor $F(1\,{\rm pc})$ and pair-correlation $g(s)$ estimators used on the data, and quote the empirical $p$-value $(1+\#\{F_{\rm null}\geq F_{\rm obs}\})/(1+N_{\rm MC})$ and the pointwise $2.5$th--$97.5$th percentile envelope. Two properties of this null should be kept in mind. It is conditional on the deposited skeleton, so it tests the distribution of cores along a fixed network and not the network itself. And it is \emph{conservative}: the column density entering Eq.~\ref{eq:nullintensity} is the observed map, which the cores themselves raise, so part of the signal being tested for is built into the null. That is the same endogeneity discussed in Section~\ref{sec:eligible}, and it biases the test against detecting excess clumping. Non-detections should therefore be read as upper limits on excess clumping rather than as evidence of its absence.

\emph{Why a contracting filament fragments below the equilibrium wavelength.} The numerical result of Section~\ref{sec:sims} has a simple WKB reading, which we give here because it is the reason we regard the counterexample as generic rather than accidental. In a filament of central density $\rho_c(t)$ the local Jeans length is $\lambda_{\rm J}(t)=c_s\sqrt{\pi/G\rho_c(t)}$, and an axial mode of fixed comoving wavelength grows at a rate $\Gamma(\lambda,t)$ that is positive only for $\lambda\gtrsim\lambda_{\rm J}(t)$. A mode is realised as a fragment only if it has time to grow before the background changes, that is if
\begin{equation}
\tau_{\rm grow}(\lambda)\;\equiv\;\Gamma(\lambda,t)^{-1}\;\lesssim\;\tau_{\rm contract}\;\equiv\;\left|\frac{\rm d\ln\rho_c}{{\rm d}t}\right|^{-1}.
\label{eq:freezeout}
\end{equation}
For an equilibrium cylinder $\tau_{\rm contract}\to\infty$ and every unstable mode has unlimited time, so the mode that dominates is the one with the largest $\Gamma$, the classical $\lambda_{\rm m}\approx22H$. For a contracting filament $\tau_{\rm contract}$ is of order the free-fall time and $\lambda_{\rm J}(t)$ falls as $\rho_c^{-1/2}$, so long-wavelength modes, whose growth times are the longest, do not complete before the state they were computed for has gone. The surviving modes are those with $\tau_{\rm grow}\lesssim\tau_{\rm contract}$, and since $\Gamma$ increases towards shorter $\lambda$ over the converged band this selects the short end. The fragmentation scale is then set by the competition in Eq.~\ref{eq:freezeout} rather than by the maximum of $\Gamma$, which is why no turnover appears and why only an upper bound $\lambda_{\rm turn}$ can be quoted. The same argument predicts that the effect should weaken smoothly as the initial contraction rate is reduced towards zero, which is the ladder in $A=v_r/v_{r,0}$ proposed in Section~\ref{sec:validation_endtoend} and not carried out here.

\section{Complete simulation campaign list}
\label{sec:campaigns}

Table~\ref{tab:campaigns} lists every simulation campaign cited in the manuscript. Rows 5--8 are the components of what was previously reported as a single Field-Geometry campaign (314 runs); they are counted individually here and the 314 aggregate is \emph{not} added separately.

\begin{table*}
\caption{Complete list of simulation campaigns underlying the paper. The ``Location'' column points, via cross-reference, to the section or appendix where each campaign is discussed in this consolidated manuscript. Campaign~9 aggregates the turbulent-support $f_{\rm eff}$ grid (90 runs; Table~\ref{tab:feff}), the physical-vs-synthetic turbulence comparison (54 runs), and the perpendicular-$\beta$ and critical-transition grids. ``FG comp.'' denotes a component of the former Field-Geometry campaign (rows 5--8).}
\label{tab:campaigns}
\footnotesize
\setlength{\tabcolsep}{4pt}
\renewcommand{\arraystretch}{1.0}
\begin{tabular}{@{}c p{3.4cm} l p{4.0cm} r l@{}}
\toprule
\# & Campaign & Location & Coverage $(f,\beta,\mathcal{M},\theta,\eta_{\rm AD})$ & Runs & Indep.? \\
\midrule
1  & Base three-regime grid & App.~\ref{sec:supplementary_campaigns_app} & $f\approx1$; $\beta\,0.05$--$5$; $\mathcal{M}\,0.5$--$5$; $\theta=0$; ideal & 208 & yes \\
2  & Transition grid / DTC & App.~\ref{sec:supplementary_campaigns_app} & $f\,1.4$--$2.2$; $\beta\,0.3$--$1.3$; $\mathcal{M}\,1$--$5$; $\theta=0$ & 540 & yes \\
3  & Supercritical grid & \S\ref{sec:modecompetition} & $f\,1.1$--$3.0$; $\beta\,0.3$--$5.0$; $\mathcal{M}\,0.5$--$3$; periodic & 654 & yes \\
4  & Validation & App.~\ref{sec:validation} & resolution / IC / EOS convergence & 83 & yes \\
5  & Near-critical & \S\ref{sec:modecompetition} & $f\,1.0$--$1.2$; $\beta\,0.3$--$1.0$; $\mathcal{M}\,1$--$2$; $\theta=0$ & 80 & FG comp. \\
6  & Perpendicular & \S\ref{sec:modecompetition} & $\theta=90^\circ$; $f\,1.2$--$1.5$; $\beta\,0.3$--$2.0$ & 96 & FG comp. \\
7  & Oblique calibration & App.~\ref{sec:supplementary_campaigns_app} & $\theta=30/45/60^\circ$; $\beta\,0.5$--$2.0$ & 108 & FG comp. \\
8  & Adiabatic EOS & App.~\ref{sec:validation} & $\gamma=5/3$; $f\,1.0$--$1.2$ & 30 & FG comp. \\
9  & Turbulence + perp-$\beta$ + crit.-transition & App.~\ref{sec:supplementary_campaigns} & $\delta v/c_s\,10^{-4}$--$1$; $\beta\,0.3$--$2$ & 289 & yes (agg.) \\
10 & Targeted DTC re-runs & App.~\ref{sec:supplementary_campaigns_app} & $\beta=0.3,\mathcal{M}=1$; extended timeout & 25 & yes \\
11 & Direct supercritical ensemble & \S\ref{sec:modecompetition} & $f\,1.5$--$2.0$; $\beta\,0.5$--$2.0$; $\theta=0$; ambient & 39 & yes \\
12 & Reconciliation suite & \S\ref{sec:modecompetition} & $f\times\beta\times2$ seeds; $\theta=0$; ambient & 12 & yes \\
13 & Ideal-MHD perp.\ $d=1.0$ grid & \S\ref{sec:modecompetition} & $\theta=90^\circ$; $\beta\,0.5$--$2.0$; ambient & 18 & yes \\
14 & Ambipolar non-ideal survey & \S\ref{sec:modecompetition} & $\theta=90^\circ$; $\eta_{\rm AD}\,0.001$--$0.1$; $\beta\,0.5$--$2.0$ & 36 & yes \\
15 & Growth-rate + equil.\ control + transverse/higher-$f$ tests & \S\ref{sec:validation_endtoend}, App.~\ref{sec:val_appendix} & $f\,0.7$--$2.0$; $128^3$--$512^3$; Gaussian/Ostriker; transverse $2$--$4\,\lambda_{\rm J}$ & 25 & yes \\
16 & Axial-resolution stability & \S\ref{sec:optional_tests} & $f=1.1$; $\theta=0$; $1024\times128^2$ rerun & 1 & yes \\
17 & Ideal-perp.\ refinement ladder & \S\ref{sec:optional_tests} & $\theta=90^\circ$; $f=1.1$; $128^3$--$512^3$ & 3 & yes \\
18 & Divergence-free helical grid & \S\ref{sec:optional_tests} & $f=1.1$; $B_\phi/B_z=0$--$2$; $512\times128^2$ & 4 & yes \\
\midrule
   & \textbf{Athena++ MHD total} & & & $\boldsymbol{2251}$ & \\
\bottomrule
\end{tabular}
\end{table*}

\section{Supplementary simulation-campaign results}
\label{sec:supplementary_campaigns_app}

The campaigns in this appendix are exploratory and are \emph{not} used quantitatively in the main conclusions; the headline wavelength rests solely on the near-critical longitudinal growth-rate ladder of Section~\ref{sec:validation_endtoend}. For completeness, the \emph{oblique}-field campaign (108 runs, $\theta=30/45/60^\circ$) shows the collapse time decreasing monotonically with field angle ($t_{\rm frag}=0.60,0.51,0.45\,t_{\rm J}$), reflecting accelerating \emph{radial} collapse as axial support weakens rather than a shortening of the axial mode; with three angles and no per-angle error this is empirical consistency, not a calibrated $\lambda_{\rm turn}(\theta,\beta)$ relation, so we draw no oblique $\lambda/W$ conclusion.

\subsection{The radial-collapse/fragmentation timescale competition}
\label{sec:regimes}

The fate of an isothermal, self-gravitating filament is governed by the competition between radial collapse and longitudinal fragmentation \citep{Inutsuka1997,Pon2012}. Near-critical filaments ($f\lesssim1.2$) evolve quasi-statically, so the longitudinal instability grows to non-linearity and $\lambda/W$ is measurable directly; they bead at $t_{\rm frag}\approx1.1$--$1.2\,t_{\rm J}$, and runs at the exact Jeans-critical value ($f=1.00$) also bead (at $t_{\rm frag}=1.62\pm0.03\,t_{\rm J}$), so there is axial fragmentation is not switched off at $f=1$. In the periodic supercritical grid ($f=1.5$--$3.0$) we find only pure radial collapse (longitudinal variance $<0.02\%$), with runaway before beading: the radial-collapse time ($0.25$--$0.34\,t_{\rm J}$) is $3$--$5\times$ shorter than the near-critical fragmentation time, and extended-integration tests confirm the periodic grid's early termination fires before the instability becomes non-linear ($\sigma_{\rm lon}$ onset $t=0.42$--$0.71\,t_{\rm J}$ vs a Courant kill $t\approx0.25$--$0.34\,t_{\rm J}$, resolution-independent). Supercritical filaments fragment longitudinally only when confined by an ambient medium (Section~\ref{sec:modecompetition}), at a converged \emph{upper bound} $\lesssim0.5$ times the classical value. The measurable wavelengths therefore come from the near-critical quasi-static regime and from ambient-confined supercritical filaments.

\subsection{Three fragmentation regimes}

The baseline grid ($f\approx1$) reveals three fragmentation regimes in the $(\beta,\mathcal{M})$ plane, set by the end-of-run density contrast $C=\rho_{\rm max}/\rho_0$: strong-field/magnetically supported ($\beta\lesssim0.15$), magnetically regulated ($0.2\lesssim\beta\lesssim2$) and thermally dominated ($\beta\gtrsim3$). The dispersion relation \citep{Nakamura1993} predicts near-complete suppression for $\beta\lesssim0.15$ and $<10\%$ wavelength modification for $\beta\gtrsim2$--$3$, consistent with these boundaries (uncertain at the one-grid-step level). Observed dense-filament conditions ($\mathcal{M}\approx1$--$3$, $\beta\approx0.3$--$2$; \citealt{Arzoumanian2019,Planck2016,Pattle2018}) place most HGBS filaments in the magnetically regulated regime.

\subsection{Turbulence and critical-transition dependence}
\label{sec:supplementary_campaigns}

Additional grids map the turbulence, field-geometry and critical-transition dependence of $\lambda/W$.

\subsubsection{Turbulence versus $\lambda/W$}

Comparing physical turbulence ($\delta v/c_s=1.0$, Kolmogorov) with synthetic perturbations ($10^{-4}$) at longitudinal field ($f=1.0$--$1.2$, $\beta=0.5,1,2$), the fragmentation \textit{wavelength} is unchanged ($<5\%$; $\lambda/W=3.44\pm0.76$ vs $3.30\pm0.85$), the $\beta$-trend appearing only in the late-time bead-count ratio (a late-epoch effect, not a shift in $\lambda_{\rm turn}$). Turbulence accelerates collapse by $3.2\times$ ($t_{\rm frag}=0.33\pm0.06$ vs $1.06\pm0.16\,t_{\rm J}$). As a single decaying realisation this does not resolve whether \emph{driven} transonic turbulence is a first-order parameter; we conclude only that the spacing is insensitive to the seed amplitude in this limited test.

\subsubsection{Perpendicular-field spacing}

For perpendicular fields ($\theta=90^\circ$, $f=1.2$--$1.5$, $\beta=0.3$--$2.0$; extended $16\lambda_{\rm J}$ periodic axial domain): strong perpendicular B ($\beta\leq0.5$) gives pure radial collapse with no axial fragmentation, while for $\beta\geq1.0$ axial beading occurs at $\lambda/W_{\rm core}=1.25\pm0.09$ ($N=27$), i.e.\ $\lambda/W_{\rm fil}\approx0.4$ under the corrected conversion (Eq.~\ref{eq:t1_apply}), a factor $2$--$3$ \textit{shorter} than the longitudinal-field value at the same $\beta$. Field geometry, not $\beta$, dominates, and the perpendicular spacing lies far below the observed ACS value (fixed-width, periodic-corrected, $\lambda/W\approx1.9$; raw $2.1$).

The two perpendicular campaigns give different $\beta$-trends (this periodic-domain grid beads only for $\beta\geq1.0$; the ambient-confined ideal grid of Section~\ref{sec:modecompetition} beads at $\beta=0.5,2.0$), but both are numerically unresolved ($N_J\approx1$--$2$; Appendix~\ref{sec:validation}), so we draw no physical $\beta$-trend from either and defer a matched-resolution grid (Section~\ref{sec:limitations}); the robust statement is only that the perpendicular beading wavelength, where it occurs, is short ($\lambda/W_{\rm fil}\lesssim0.6$), well below the observed spacing.

\smallskip\noindent\emph{Ambipolar-perpendicular survey (construction and status).} The preliminary exploratory ambipolar-diffusion survey summarised in Section~\ref{sec:modecompetition} comprises 36 runs spanning $\eta_{\rm AD}=0.001$--$0.1$ and $\beta=0.5$--$2.0$ at $\theta=90^\circ$. Of these $32/36$ develop axial beads, but at a minimum resolution of only $N_J\approx1$--$2$ cells per local Jeans length; below the Truelove criterion, where under-resolution itself produces artificial fragmentation \citep{Truelove1997}. The median onset spacing is $\lambda/W_{\rm fil}\approx0.6$, and the runs reach density runaway within $\Delta t\approx0.02$--$0.04\,t_{\rm J}$ with no converged plateau, so we cannot separate a genuine ambipolar enhancement from a numerical one and treat the grid as unresolved.

\smallskip\noindent\emph{Divergence-free helical-field construction.} For the seeded-helical scan of Section~\ref{sec:optional_tests} we impose a divergence-free helical field: a uniform axial $B_z$ plus a toroidal $B_\phi(r)$ peaking at $r_c=W_{\rm core}$, obtained as a discrete curl of the vector potential $A_z(r)=B_t\,r_c\ln[1+(r/r_c)^2]$ so that $\nabla\!\cdot\!\mathbf{B}=0$ to machine precision. We scan the toroidal-to-axial ratio $B_\phi/B_z=0$--$2$ (with $B_\phi=0$ reproducing the pure-longitudinal growth-rate spectrum, a code check). Across the scan the axial growth-rate plateau extends to the same converged-band edge, $\lambda\approx1\,\lambda_{\rm J}$, for every toroidal fraction, with no excess growth at shorter $\lambda$ even for a strongly-wound field (Fig.~\ref{fig:helical_scan}a); the only systematic effect is a radial pinch that narrows the formation FWHM from $0.57$ to $0.43\,\lambda_{\rm J}$ ($\approx25\%$; Fig.~\ref{fig:helical_scan}b), with onset at $B_\phi\gtrsim B_z$. Because the field is seeded on an already-contracting filament rather than constructed in force balance, it drives an out-of-equilibrium pinch and cannot represent the confining toroidal equilibrium whose sausage ($m=0$) mode carries the shortened \citet{Fiege2000} wavelength; these tests therefore rule out only \emph{dynamic}, out-of-equilibrium helical pinching as a source of axial-mode shortening.

\subsubsection{Critical-transition spacing}

Mapping $\lambda/W$ across $f=0.9$--$1.3$ at five $\beta$ values shows no discontinuity at $f=1.0$ (variation $<5\%$) and an apparent $\beta$-dependence in the \emph{late-time bead count} ($\lambda/W_{\rm core}=4.74$--$2.80$ from $\beta=0.3$ to $2.0$). This bead-count trend is a late-epoch, resolution-sensitive artefact, larger than the weak mode-quantised trend of the primary growth-rate measurement, with the highest-$\beta$ points ($\lambda/\lambda_{\rm J}\approx0.84$) consistent with the supercritical ensemble. \textbf{Continuity across $f\approx1$ does not validate extrapolation to $f\geq1.5$}, where the mode changes to radial collapse; the near-critical $\lambda/W$ values are a hypothesis for supercritical filaments, not a measurement.

\begin{table}
\caption{Effective line-mass fraction with turbulent support.}
\label{tab:feff}
\footnotesize
\setlength{\tabcolsep}{4pt}
\begin{tabular}{p{3.0cm}cc}
\toprule
Quantity & Value & Source \\
\midrule
Nominal $f$ (thermal) & $1.5 \pm 0.3$ & HGBS line-mass \\
Turbulent support $f_{\rm eff}/f$ & $0.82 \pm 0.11$ & TURB (90 sims) \\
Effective $f_{\rm eff}$ & $1.23 \pm 0.28$ & product \\
HGBS ensemble range & $0.7$--$1.9$ & propagated \\
\bottomrule
\end{tabular}
\end{table}

\subsubsection{Synthesis}

In convergence-independent units the two first-order parameters are field geometry (perpendicular $\lambda/W_{\rm core}\approx1.25$ versus longitudinal $2.8$--$4.4$) and, for longitudinal fields, plasma $\beta$; the observable comparison uses the growth-rate wavelength (Section~\ref{sec:validation_endtoend}), not these raw ratios.

\bibliography{references_complete}

\end{document}
